% PALAEONTOLOGY — Geology Upper Sixth — Cours signé OnBuch+
% Official MINESEC syllabus. Compiler avec : tectonic GEO02-palaeontology.tex
\newcommand{\DOCMATIERE}{Geology}
\newcommand{\DOCNIVEAU}{Upper Sixth}
\newcommand{\DOCMODULE}{Upper Sixth — Geology}
\newcommand{\DOCLECON}{2}
\newcommand{\DOCTITRE}{PALAEONTOLOGY}
% OnBuch+ — préambule « cours signés » ENRICHI (compilé avec tectonic / XeLaTeX)
% Système visuel « Pop » d'OnBuch+ : fond crème (jamais blanc), bordures encre
% épaisses, ombres « dures » décalées, titres Archivo Black en capitales.
% Version MATHÉMATIQUES : graphes (pgfplots/tikz), boîtes pédagogiques
% (définition, propriété, méthode, attention, le savais-tu, exercice résolu,
% exercice, TP, bilan), page de garde et signature OnBuch+ ancrée en bas.
%
% Macros attendues dans main.tex AVANT \input{preamble} :
%   \DOCMATIERE  \DOCNIVEAU  \DOCMODULE  \DOCLECON (numéro)  \DOCTITRE
\documentclass[11pt]{article}

\usepackage{fontspec}
\usepackage[english]{babel}
\usepackage[a4paper,margin=2cm,top=3.4cm,bottom=2.8cm,headheight=1.4cm,footskip=1.3cm]{geometry}
\usepackage{amsmath,amssymb,amsfonts}
\usepackage{mathtools} % \xmapsto, \coloneqq... souvent utilisés par les modèles
\usepackage{cancel} % simplifications barrées (bilans de forces)
% \abs{z} (module, valeur absolue) et \norm{u} (norme), utilisés par les modèles
\providecommand{\abs}{}\let\abs\relax\DeclarePairedDelimiter{\abs}{\lvert}{\rvert}
\providecommand{\norm}{}\let\norm\relax\DeclarePairedDelimiter{\norm}{\lVert}{\rVert}
\usepackage[table]{xcolor} % [table] : fournit aussi \rowcolors (alternance de lignes)
\usepackage{pagecolor}
\usepackage{tikz}
\usetikzlibrary{arrows.meta,positioning,calc,shapes.geometric,shapes.misc,shapes.arrows,decorations.pathmorphing,decorations.pathreplacing,decorations.markings,patterns,shadows,mindmap,trees,fit,backgrounds,matrix,intersections,angles,quotes}
\usepackage{pgfplots}
\usepackage[version=4]{mhchem} % équations nucléaires \ce{^{235}_{92}U}
\usepackage[european,siunitx]{circuitikz} % schémas électriques
\pgfplotsset{compat=1.18}
\usepgfplotslibrary{fillbetween,groupplots} % aires entre courbes (calcul intégral)
\usepackage{enumitem}
\usepackage{fancyhdr}
% [most] charge toutes les bibliothèques tcolorbox (listings, minted, raster,
% documentation...) et gonfle inutilement la mémoire TeX sur un document long
% (~300 boîtes) : on ne charge que ce qui est réellement utilisé.
\usepackage[skins,breakable]{tcolorbox}
\usepackage{graphicx}
\usepackage{colortbl}
\usepackage{booktabs}
\usepackage{tabularx}
\usepackage{diagbox} % case d'en-tête diagonale des tableaux à double entrée
% Colonnes extensibles centrée / gauche / droite (C, L, R), souvent utilisées par les modèles
\newcolumntype{C}{>{\centering\arraybackslash}X}
\newcolumntype{L}{>{\raggedright\arraybackslash}X}
\newcolumntype{R}{>{\raggedleft\arraybackslash}X}
\usepackage{array}
\usepackage{multicol}
\usepackage{siunitx}
% parse-numbers=false : accepte aussi \SI{2,5 \times 10^{-3}}{...}
\sisetup{locale=UK,output-decimal-marker={.},group-minimum-digits=4,parse-numbers=false}
\DeclareSIUnit\ohm{\ensuremath{\symup{\Omega}}} % U+2126 absent de la police
\DeclareSIUnit\division{div} % réglages d'oscilloscope (ms/div, V/div)
\DeclareSIUnit\tr{tr} % tours (vitesse de rotation, ex. \SI{1500}{\tr\per\minute})
\usepackage{qrcode}
% \degree : commande fréquemment utilisée par les modèles pour un angle ou une
% température en dehors de \SI{}{} (non fournie par les paquets chargés).
\providecommand{\degree}{^{\circ}}
% \qed (fin de démonstration), utilisé par les modèles sans amsthm
\providecommand{\qed}{\hfill$\square$}
% \barre : double barre des valeurs interdites dans un tableau de variations (tkz-tab)
\providecommand{\barre}{\ensuremath{\|}}
% environnement proof (amsthm non chargé) utilisé par les modèles
\ifdefined\proof\else\newenvironment{proof}[1][Proof]{\par\noindent\textit{#1.}\ }{\qed\par}\fi
\usepackage{lastpage}
\usepackage{hyperref}
\hypersetup{hidelinks}

% ============================================================
% Palette « Pop » OnBuch+ — mêmes hex que la classe P (plus_ui.dart)
% ============================================================
\definecolor{poporange}{HTML}{F59321}
\definecolor{popdark}{HTML}{E07A0C}
\definecolor{popcream}{HTML}{FFF6E8}
\definecolor{popink}{HTML}{1C1714}
\definecolor{poppurple}{HTML}{7A5AE0}
\definecolor{popgreen}{HTML}{2E9E6A}
\definecolor{poppink}{HTML}{E0457B}
\definecolor{popblue}{HTML}{2F7DE1}
\definecolor{popgold}{HTML}{C98A12}
\definecolor{popmuted}{HTML}{8A7F76}
\definecolor{popline}{HTML}{EADFD2}
\definecolor{popgray}{HTML}{8A7F76}
\colorlet{popgrey}{popgray}
% Alias tolérants (le modèle nomme parfois les couleurs différemment)
\colorlet{popwhite}{white}
\colorlet{popblack}{popink}
\colorlet{popred}{poppink}
\colorlet{popyellow}{popgold}
% Teintes claires pour fonds de tableaux / zones
\colorlet{popblueL}{popblue!10!white}
\colorlet{poporangeL}{poporange!14!white}
\colorlet{popgreenL}{popgreen!12!white}
\colorlet{poppurpleL}{poppurple!12!white}
\colorlet{poppinkL}{poppink!12!white}
\colorlet{popgoldL}{popgold!14!white}

\pagecolor{popcream}

% ============================================================
% Typographie
% ============================================================
\defaultfontfeatures{Ligatures=TeX}
\setmainfont{PlusJakartaSans-Medium.ttf}[Path=../../../fonts/,BoldFont=PlusJakartaSans-Bold.ttf]
% Maths en sans-serif (Fira Math) avec chiffres et lettres droites de Plus
% Jakarta, pour que les formules se fondent dans le texte.
\usepackage[math-style=ISO,bold-style=ISO]{unicode-math}
\setmathfont{FiraMath-Regular.otf}[Path=../../../fonts/]
\setmathfont{PlusJakartaSans-Medium.ttf}[Path=../../../fonts/,range={up/num,up/{Latin,latin},\mathup}]
% (icomma retiré : décimales avec point en anglais)
% Caractères mathématiques Unicode absents des polices de texte (Plus Jakarta,
% Archivo Black) : rendus par leur équivalent mathématique, en texte comme en
% formule — sinon ils s'affichent en carré vide (ex. « Arithmétique dans ℤ »).
\usepackage{newunicodechar}
\newunicodechar{ℝ}{\ensuremath{\mathbb{R}}}
\newunicodechar{ℤ}{\ensuremath{\mathbb{Z}}}
\newunicodechar{ℂ}{\ensuremath{\mathbb{C}}}
\newunicodechar{ℕ}{\ensuremath{\mathbb{N}}}
\newunicodechar{ℚ}{\ensuremath{\mathbb{Q}}}
\newunicodechar{𝔻}{\ensuremath{\mathbb{D}}}
\newunicodechar{α}{\ensuremath{\alpha}}
\newunicodechar{β}{\ensuremath{\beta}}
\newunicodechar{γ}{\ensuremath{\gamma}}
\newunicodechar{δ}{\ensuremath{\delta}}
\newunicodechar{ε}{\ensuremath{\varepsilon}}
\newunicodechar{θ}{\ensuremath{\theta}}
\newunicodechar{λ}{\ensuremath{\lambda}}
\newunicodechar{μ}{\ensuremath{\mu}}
\newunicodechar{σ}{\ensuremath{\sigma}}
\newunicodechar{τ}{\ensuremath{\tau}}
\newunicodechar{φ}{\ensuremath{\varphi}}
\newunicodechar{ψ}{\ensuremath{\psi}}
\newunicodechar{ω}{\ensuremath{\omega}}
\newunicodechar{Σ}{\ensuremath{\Sigma}}
% majuscules produites par \MakeUppercase dans les titres (α-aminés -> Α)
\newunicodechar{Α}{\ensuremath{\alpha}}
\newunicodechar{Β}{\ensuremath{\beta}}
\newunicodechar{Γ}{\ensuremath{\Gamma}}
\newunicodechar{Δ}{\ensuremath{\Delta}}
\newunicodechar{Ε}{\ensuremath{\varepsilon}}
\newunicodechar{Θ}{\ensuremath{\Theta}}
\newunicodechar{Λ}{\ensuremath{\Lambda}}
\newunicodechar{Μ}{\ensuremath{\mu}}
\newunicodechar{Τ}{\ensuremath{\tau}}
\newunicodechar{Φ}{\ensuremath{\Phi}}
\newunicodechar{Ψ}{\ensuremath{\Psi}}
\newunicodechar{Ω}{\ensuremath{\Omega}}
\newunicodechar{↦}{\ensuremath{\mapsto}}
\newunicodechar{∈}{\ensuremath{\in}}
\newunicodechar{∉}{\ensuremath{\notin}}
\newunicodechar{∧}{\ensuremath{\wedge}}
\newunicodechar{⇔}{\ensuremath{\Leftrightarrow}}
\newunicodechar{⟺}{\ensuremath{\Longleftrightarrow}}
\newunicodechar{⇒}{\ensuremath{\Rightarrow}}
\newunicodechar{⟹}{\ensuremath{\Longrightarrow}}
\newunicodechar{∘}{\ensuremath{\circ}}
\newunicodechar{∪}{\ensuremath{\cup}}
\newunicodechar{∩}{\ensuremath{\cap}}
\newunicodechar{≡}{\ensuremath{\equiv}}
\newunicodechar{⊂}{\ensuremath{\subset}}
\newunicodechar{⊥}{\ensuremath{\perp}}
\newunicodechar{∃}{\ensuremath{\exists}}
\newunicodechar{∀}{\ensuremath{\forall}}
\newunicodechar{∛}{\ensuremath{\sqrt[3]{\,}}}
\newunicodechar{∜}{\ensuremath{\sqrt[4]{\,}}}
\newunicodechar{⃗}{\ensuremath{{}^{\to}}}
\newunicodechar{ⁿ}{\ifmmode^{n}\else\textsuperscript{\ensuremath{n}}\fi}
\newunicodechar{ˣ}{\ifmmode^{x}\else\textsuperscript{\ensuremath{x}}\fi}
\newunicodechar{ᵇ}{\ifmmode^{b}\else\textsuperscript{\ensuremath{b}}\fi}
\newunicodechar{ʳ}{\ifmmode^{r}\else\textsuperscript{\ensuremath{r}}\fi}
\newunicodechar{ᵘ}{\ifmmode^{u}\else\textsuperscript{\ensuremath{u}}\fi}
\newunicodechar{ʸ}{\ifmmode^{y}\else\textsuperscript{\ensuremath{y}}\fi}
\newunicodechar{ᵃ}{\ifmmode^{a}\else\textsuperscript{\ensuremath{a}}\fi}
\newunicodechar{ᵉ}{\ifmmode^{e}\else\textsuperscript{\ensuremath{e}}\fi}
\newunicodechar{ᵏ}{\ifmmode^{k}\else\textsuperscript{\ensuremath{k}}\fi}
\newunicodechar{ᵖ}{\ifmmode^{p}\else\textsuperscript{\ensuremath{p}}\fi}
\newunicodechar{ᴺ}{\ifmmode^{N}\else\textsuperscript{\ensuremath{N}}\fi}
\newunicodechar{ᶜ}{\ifmmode^{c}\else\textsuperscript{\ensuremath{c}}\fi}
\newunicodechar{ʰ}{\ifmmode^{h}\else\textsuperscript{\ensuremath{h}}\fi}
\newunicodechar{ᵠ}{\ifmmode^{q}\else\textsuperscript{\ensuremath{q}}\fi}
\newunicodechar{ₙ}{\ifmmode_{n}\else\textsubscript{\ensuremath{n}}\fi}
\newunicodechar{ₐ}{\ifmmode_{a}\else\textsubscript{\ensuremath{a}}\fi}
\newunicodechar{ᵢ}{\ifmmode_{i}\else\textsubscript{\ensuremath{i}}\fi}
\newunicodechar{ᵣ}{\ifmmode_{r}\else\textsubscript{\ensuremath{r}}\fi}
\newunicodechar{ₖ}{\ifmmode_{k}\else\textsubscript{\ensuremath{k}}\fi}
\newunicodechar{ᵦ}{\ifmmode_{\beta}\else\textsubscript{\ensuremath{\beta}}\fi}
\newunicodechar{ₓ}{\ifmmode_{x}\else\textsubscript{\ensuremath{x}}\fi}
\newfontfamily\popdisplay{ArchivoBlack-Regular.ttf}[Path=../../../fonts/]
\newfontfamily\poplogo{SpaceGrotesk-Bold.ttf}[Path=../../../fonts/]
\newfontfamily\popsemi{PlusJakartaSans-SemiBold.ttf}[Path=../../../fonts/]
\newfontfamily\popxbold{PlusJakartaSans-ExtraBold.ttf}[Path=../../../fonts/]
\newfontfamily\popxboldls{PlusJakartaSans-ExtraBold.ttf}[Path=../../../fonts/,LetterSpace=3.0]

\setlist[enumerate]{leftmargin=1.7em,itemsep=5pt,topsep=4pt}
\setlist[itemize]{leftmargin=1.7em,itemsep=4pt,topsep=4pt,label={\color{poporange}\textbullet}}
\setlength{\parindent}{0pt}
\setlength{\parskip}{5pt}
\renewcommand{\arraystretch}{1.25}

% pgfplots : style Pop par défaut
\pgfplotsset{
  every axis/.append style={axis line style={popink,line width=1pt},tick style={popink},
    grid style={popline,line width=0.5pt},label style={font=\small},tick label style={font=\footnotesize}},
  popcurve/.style={poporange,line width=1.8pt,no markers,smooth},
}
% Même style disponible dans un \draw TikZ simple (hors pgfplots)
\tikzset{popcurve/.style={poporange,line width=1.8pt}}
\tikzset{popgrid/.style={popline,very thin}}
\colorlet{popcurve}{poporange} % le nom du style est parfois utilisé comme couleur

% ============================================================
% Pill / squiggle
% ============================================================
\newcommand{\poppill}[3]{%
  \tikz[baseline=-0.55ex]\node[draw=popink,line width=1.2pt,rounded corners=2.6pt,%
    fill=#2,inner xsep=6.5pt,inner ysep=3pt]{%
    \popxboldls\fontsize{7}{7}\selectfont\color{#3}\MakeUppercase{#1}};%
}
\newcommand{\popsquiggle}[1]{%
  \tikz[baseline=-0.4ex]{
    \draw[#1,line width=2.6pt,line cap=round]
      plot [smooth] coordinates {(0,0.09) (0.34,0.01) (0.68,0.21) (1.02,0.01) (1.36,0.10)};
  }%
}
% Mot-clé surligné (marqueur orange sous le texte)
\newcommand{\cle}[1]{{\popxbold\color{popdark}#1}}

% ============================================================
% En-tête / pied
% ============================================================
\pagestyle{fancy}
\fancyhf{}
\renewcommand{\headrulewidth}{0pt}
\renewcommand{\footrulewidth}{0pt}
\lhead{\raisebox{-0.15ex}{\poplogo\fontsize{13}{13}\selectfont\color{popink}OnBuch\color{poporange}+}}
\rhead{\poppill{\DOCMATIERE\ \textbullet\ \DOCNIVEAU\ \textbullet\ Chapter \DOCLECON}{white}{popink}}
\lfoot{\popxbold\fontsize{7.6}{7.6}\selectfont\color{popmuted}\MakeUppercase{Signed course OnBuch+}\ \textbullet\ {\popsemi\DOCTITRE}}
\rfoot{\tikz[baseline=-0.6ex]\node[rounded corners=8.5pt,fill=popink,minimum height=17pt,minimum width=17pt,inner xsep=5pt,inner sep=0]{\popxbold\fontsize{8}{8}\selectfont\color{popcream}\thepage\,/\,\pageref*{LastPage}};}

% ============================================================
% Titres
% ============================================================
\newcommand{\chaptitre}[2]{%
  \begin{tcolorbox}[enhanced,colback=poporange,colframe=popink,boxrule=2.6pt,arc=11pt,
      drop shadow=popink,left=15pt,right=15pt,top=12pt,bottom=11pt,before skip=3pt,after skip=18pt]
    \poppill{#2}{popink}{popcream}\par\vspace{9pt}
    {\raggedright\hyphenpenalty=10000\exhyphenpenalty=10000\popdisplay\fontsize{22}{24}\selectfont\color{popcream}\MakeUppercase{#1}\par}
  \end{tcolorbox}
}
% Section numérotée : pastille orange avec le numéro + titre Archivo
\newcounter{popsec}
\newcommand{\coursec}[1]{%
  \stepcounter{popsec}\setcounter{popsub}{0}%
  \par\vspace{16pt}\noindent
  \begin{minipage}{\linewidth}
  \tikz[baseline=-0.7ex]\node[draw=popink,line width=1.6pt,fill=poporange,rounded corners=4pt,minimum size=22pt,inner sep=0]{\popdisplay\fontsize{12}{12}\selectfont\color{popcream}\thepopsec};%
  \hspace{8pt}{\popdisplay\fontsize{15}{17}\selectfont\color{popink}\MakeUppercase{#1}}\par
  \vspace{4pt}\noindent{\color{poporange}\rule{36pt}{3.6pt}}
  \end{minipage}\par\nopagebreak\vspace{8pt}
}
\newcounter{popsub}
\newcommand{\courssub}[1]{%
  \stepcounter{popsub}%
  \par\vspace{9pt}\noindent{\popxbold\fontsize{11.5}{13}\selectfont\color{popink}{\color{poporange}\thepopsec.\thepopsub}\hspace{6pt}#1}\par\nopagebreak\vspace{4pt}
}

% ============================================================
% Boîtes pédagogiques — même recette : fond blanc, bordure épaisse colorée,
% ombre dure encre, badge plein. Titre optionnel [#1].
% ============================================================
\tcbset{popbase/.style={breakable,enhanced,colback=white,boxrule=2.3pt,arc=9pt,
  shadow={3pt}{-3pt}{0pt}{popink},left=14pt,right=14pt,top=11pt,bottom=11pt,before skip=11pt,after skip=15pt,
  fontupper=\popsemi\fontsize{10.6}{14.5}\selectfont\color{popink},
  fontlower=\popsemi\fontsize{10.6}{14.5}\selectfont\color{popink}}}
% Coche dessinée (le glyphe ✓ n'existe pas dans les polices du document)
\newcommand{\popcheck}[1]{\tikz[baseline=-0.5ex]\draw[#1,line width=1.6pt,line cap=round,line join=round](-0.13,0.0)--(-0.04,-0.09)--(0.13,0.1);}
\providecommand{\coche}{\popcheck{popgreen}}
\providecommand{\resultat}[1]{\par\smallskip\noindent\fcolorbox{popgreen}{popgreenL}{\parbox{\dimexpr\linewidth-2\fboxsep-2\fboxrule\relax}{\textbf{Result.} #1}}\par\smallskip}
\newcommand{\popboxhead}[3]{\poppill{#1}{#2}{white}\ifx\relax#3\relax\else\hspace{6pt}{\popxbold\fontsize{10.6}{12}\selectfont\color{popink}#3}\fi\par\vspace{7pt}}

\newtcolorbox{aretenir}[1][]{popbase,colframe=popgreen,before upper={\popboxhead{Key points}{popgreen}{#1}}}
\newtcolorbox{exemplebox}[1][]{popbase,colframe=poporange,before upper={\popboxhead{Exemple}{poporange}{#1}}}
\newtcolorbox{definition}[1][]{popbase,colframe=popblue,before upper={\popboxhead{Definition}{popblue}{#1}}}
\newtcolorbox{propriete}[1][]{popbase,colframe=poppurple,before upper={\popboxhead{Property}{poppurple}{#1}}}
\newtcolorbox{methode}[1][]{popbase,colframe=popgold,before upper={\popboxhead{Method}{popgold}{#1}}}
\newtcolorbox{attention}[1][]{popbase,colframe=poppink,before upper={\popboxhead{Warning}{poppink}{#1}}}
\newtcolorbox{experience}[1][]{popbase,colframe=popblue,colback=popblueL,before upper={\popboxhead{Experiment}{popblue}{#1}}}
% « Le savais-tu ? » : carte violette pleine (contexte, culture, Cameroun)
\newtcolorbox{savaistu}[1][]{popbase,colback=poppurple,colframe=popink,
  fontupper=\popsemi\fontsize{10.4}{14.2}\selectfont\color{white},
  before upper={\poppill{Did you know?}{white}{poppurple}\ifx\relax#1\relax\else\hspace{6pt}{\popxbold\color{white}#1}\fi\par\vspace{7pt}}}
% Exercice résolu : énoncé en haut, \tcblower puis la solution sur fond crème
\newcounter{popexo}\newcounter{popexr}
\newtcolorbox{exoresolu}[1][]{popbase,colframe=poporange,colbacklower=popcream,
  segmentation style={popink,line width=1.2pt,dashed},
  before upper={\stepcounter{popexr}\popboxhead{Worked example \thepopexr}{poporange}{#1}},
  before lower={\poppill{Solution}{popink}{popcream}\par\vspace{6pt}}}
% Exercice (non résolu en place) — la correction est dans la partie Corrigés
\newtcolorbox{exercice}[1][]{popbase,colframe=popink,
  before upper={\stepcounter{popexo}\popboxhead{Exercise \thepopexo}{popink}{#1}}}
% Corrigé
\newtcolorbox{corrige}[1][]{popbase,colframe=popgreen,colback=popgreenL,
  before upper={\popboxhead{Solution}{popgreen}{#1}}}
% Objectifs / prérequis / bilan
\newtcolorbox{objectifs}{popbase,colframe=popink,colback=poporangeL,
  before upper={\popboxhead{Chapter objectives}{popink}{}}}
\newtcolorbox{prerequis}{popbase,colframe=popink,colback=popblueL,
  before upper={\popboxhead{Prerequisites}{popblue}{}}}
\newtcolorbox{bilan}{popbase,colframe=popink,colback=white,boxrule=2.8pt,
  before upper={\popboxhead{Fiche bilan}{poporange}{L'essentiel à connaître pour l'examen}}}
% Figure encadrée avec légende
\newtcolorbox{popfigure}[1][]{enhanced,breakable=false,colback=white,colframe=popline,boxrule=1.4pt,arc=8pt,
  left=10pt,right=10pt,top=10pt,bottom=8pt,before skip=10pt,after skip=12pt,halign=center,
  lower separated=false,halign lower=center,
  fontlower=\popsemi\fontsize{9}{11.5}\selectfont\color{popmuted},
  #1}
\newcommand{\legende}[1]{\par\vspace{6pt}{\popsemi\fontsize{9}{11.5}\selectfont\color{popmuted}#1}}

% ============================================================
% Signature OnBuch+
% ============================================================
\newcommand{\poplogoinline}[1]{{\poplogo\fontsize{#1}{#1}\selectfont\color{popink}OnBuch\color{poporange}+}}
\newcommand{\signatureonbuch}{%
  \begin{tcolorbox}[enhanced,colback=popink,colframe=popink,boxrule=2.2pt,arc=10pt,
      drop shadow=poporange,left=14pt,right=14pt,top=10pt,bottom=10pt,before skip=0pt,after skip=0pt]
    \tikz[baseline=-0.7ex]\node[circle,fill=popgreen,draw=popcream,line width=1.3pt,minimum size=22pt,inner sep=0]{\popcheck{white}};%
    \hspace{9pt}\begin{minipage}[c]{0.62\linewidth}
      {\popxbold\fontsize{12}{13}\selectfont\color{popcream}Signed\ \textbullet\ {\poplogo OnBuch\color{poporange}+}}\par\vspace{2pt}
      {\raggedright\popsemi\fontsize{8}{10.5}\selectfont\color{popline}\MakeUppercase{OnBuch+ teaching team}\\\MakeUppercase{Follows the official MINESEC syllabus}\par}
    \end{minipage}\hfill
    {\poplogo\fontsize{22}{22}\selectfont\color{popcream}OnBuch\color{poporange}+}
  \end{tcolorbox}%
}

% ============================================================
% Page de garde
% #1 titre · #2 matière · #3 niveau · #4 module · #5 n° leçon · #6 durée ·
% #7 description / objectifs (texte ou itemize)
% ============================================================
\newcommand{\pagegarde}[7]{%
  \thispagestyle{empty}
  \newgeometry{a4paper,margin=1.7cm,top=1.6cm,bottom=1.4cm}%
  \begin{tikzpicture}[remember picture,overlay]
    \fill[poporange!18] ([xshift=-3.2cm,yshift=-3.6cm]current page.north east) circle (4.6cm);
    \fill[poppurple!14] ([xshift=2.4cm,yshift=7.6cm]current page.south west) circle (3.8cm);
  \end{tikzpicture}%
  {\poplogo\fontsize{30}{30}\selectfont\color{popink}OnBuch\color{poporange}+}\hfill
  \raisebox{0.6ex}{\poppill{Signed course \textbullet\ Premium}{popink}{popcream}}\par
  \vspace{1pt}\popsquiggle{poppurple}\par
  \vspace{8pt}
  \begin{tcolorbox}[enhanced,colback=poporange,colframe=popink,boxrule=2.8pt,arc=13pt,
      drop shadow=popink,left=18pt,right=18pt,top=12pt,bottom=13pt,after skip=10pt]
    \poppill{#2 \textbullet\ #3}{popink}{popcream}\hspace{5pt}\poppill{#4}{white}{popink}\par\vspace{8pt}
    {\popdisplay\fontsize{14}{14}\selectfont\color{popink}CHAPTER #5}\par\vspace{4pt}
    {\raggedright\hyphenpenalty=10000\exhyphenpenalty=10000\popdisplay\fontsize{25}{27}\selectfont\color{popcream}\MakeUppercase{#1}\par}
    \vspace{8pt}
    {\popxbold\fontsize{9.5}{9.5}\selectfont\color{popink}\MakeUppercase{Suggested time: #6}}
  \end{tcolorbox}
  \begin{tcolorbox}[enhanced,colback=white,colframe=popink,boxrule=2.2pt,arc=10pt,
      drop shadow=popink,left=16pt,right=16pt,top=10pt,bottom=10pt,after skip=8pt]
    \poppill{In this chapter}{poppurple}{white}\par\vspace{5pt}
    \popsemi\fontsize{9.3}{12.6}\selectfont\color{popink}#7
  \end{tcolorbox}
  \noindent\begin{minipage}[t]{0.487\linewidth}
    \begin{tcolorbox}[enhanced,colback=popgreen,colframe=popink,boxrule=2.2pt,arc=10pt,
        drop shadow=popink,left=11pt,right=11pt,top=10pt,bottom=10pt,height=3.1cm,valign=center]
      \tcbox[colback=white,colframe=popink,boxrule=1.3pt,arc=5pt,boxsep=0pt,nobeforeafter,
        left=3pt,right=3pt,top=3pt,bottom=3pt,valign=center]{\qrcode[height=1.9cm]{https://play.google.com/store/apps/details?id=com.luvvix.onbuch&hl=en}}%
      \hspace{8pt}\begin{minipage}[c]{0.5\linewidth}
        {\popdisplay\fontsize{11}{12.5}\selectfont\color{white}\MakeUppercase{Download OnBuch}}\par\vspace{4pt}
        {\raggedright\popsemi\fontsize{8.4}{11}\selectfont\color{white}Free \textbullet\ Google Play\\AI tutor, past papers, results\par}
      \end{minipage}
    \end{tcolorbox}
  \end{minipage}\hfill
  \begin{minipage}[t]{0.487\linewidth}
    \begin{tcolorbox}[enhanced,colback=poppurple,colframe=popink,boxrule=2.2pt,arc=10pt,
        drop shadow=popink,left=11pt,right=11pt,top=10pt,bottom=10pt,height=3.1cm,valign=center]
      \tikz[baseline=-0.7ex]\node[circle,fill=white,draw=popink,line width=1.3pt,minimum size=1.6cm,inner sep=0]{\popdisplay\fontsize{24}{24}\selectfont\color{poppurple}+};%
      \hspace{8pt}\begin{minipage}[c]{0.6\linewidth}
        {\popdisplay\fontsize{11}{12.5}\selectfont\color{white}\MakeUppercase{Go OnBuch+}}\par\vspace{4pt}
        {\raggedright\popsemi\fontsize{8.4}{11}\selectfont\color{white}All signed courses, unlimited AI tutor, PDF sheets — OnBuch+ tab in the app\par}
      \end{minipage}
    \end{tcolorbox}
  \end{minipage}
  \vspace{8pt}
  \popcontenu
  \vfill
  \signatureonbuch
  \restoregeometry
  \newpage
}

% Rangée « Ce cours contient » (page de garde)
\newcommand{\poptile}[3]{% #1 couleur #2 gros texte #3 légende
  \begin{tcolorbox}[enhanced,colback=#1,colframe=popink,boxrule=2pt,arc=9pt,shadow={2.5pt}{-2.5pt}{0pt}{popink},
      left=6pt,right=6pt,top=9pt,bottom=9pt,halign=center,valign=center,height=1.75cm,width=\linewidth,nobeforeafter]
    {\popdisplay\fontsize{12.5}{13}\selectfont\color{white}\MakeUppercase{#2}}\par\vspace{3pt}
    {\centering\popsemi\fontsize{7.8}{9.5}\selectfont\color{white}#3\par}
  \end{tcolorbox}}
\newcommand{\popcontenu}{%
  \par\noindent{\popxboldls\fontsize{7.5}{7.5}\selectfont\color{popmuted}THIS SIGNED COURSE CONTAINS}\par\vspace{7pt}
  \noindent\begin{minipage}[t]{0.235\linewidth}\poptile{popblue}{Course}{complete, illustrated, step by step}\end{minipage}\hfill
  \begin{minipage}[t]{0.235\linewidth}\poptile{poporange}{Methods}{and worked examples}\end{minipage}\hfill
  \begin{minipage}[t]{0.235\linewidth}\poptile{poppink}{Exercises}{exam style + solutions}\end{minipage}\hfill
  \begin{minipage}[t]{0.235\linewidth}\poptile{popgreen}{Summary sheet}{the essentials to remember}\end{minipage}\par}

% Sommaire
\newtcolorbox{sommaire}{enhanced,colback=white,colframe=popink,boxrule=2.2pt,arc=10pt,
  drop shadow=popink,left=18pt,right=18pt,top=13pt,bottom=13pt,before skip=6pt,after skip=20pt,
  before upper={\poppill{Contents}{poppurple}{white}\par\vspace{9pt}\popsemi\fontsize{10.6}{14.5}\selectfont\color{popink}}}

% Signature de fin de cours — poussée tout en bas de la dernière page
\newcommand{\findecours}{%
  \par\vfill\nopagebreak
  \begin{center}\popsquiggle{poporange}\end{center}\vspace{4pt}
  \signatureonbuch
}
\AtBeginDocument{\def\llbracket{\mathopen{[\mkern-3.5mu[}}\def\rrbracket{\mathclose{]\mkern-3.5mu]}}} % intervalles d'entiers (glyphe ⟦ absent de la police)
% Glyphes absents de Fira Math : redéfinis à partir de glyphes disponibles
\AtBeginDocument{%
  \def\setminus{\mathbin{\backslash}}\let\smallsetminus\setminus
  \def\vdots{\mathord{\vbox{\baselineskip3.2pt\lineskiplimit0pt\kern4pt\hbox{.}\hbox{.}\hbox{.}}}}%
  \def\ddots{\mathinner{\mkern1mu\raise6.5pt\hbox{.}\mkern2mu\raise3.6pt\hbox{.}\mkern2mu\raise0.7pt\hbox{.}\mkern1mu}}%
  \def\triangle{\mathord{\scalebox{0.9}{$\symup{\Delta}$}}}%
  \def\nabla{\mathord{\raisebox{\depth}{\scalebox{1}[-1]{$\symup{\Delta}$}}}}%
  \def\checkmark{\popcheck{popdark}}%
  \def\star{\mathbin{\ast}}\def\bigstar{\mathord{\ast}}%
  \def\diamond{\mathbin{\scalebox{0.6}{$\Diamond$}}}%
  \def\bigcup{\mathop{\mathchoice{\raisebox{-0.3em}{\scalebox{1.7}{$\cup$}}}{\scalebox{1.25}{$\cup$}}{\cup}{\cup}}\displaylimits}%
  \def\bigcap{\mathop{\mathchoice{\raisebox{-0.3em}{\scalebox{1.7}{$\cap$}}}{\scalebox{1.25}{$\cap$}}{\cap}{\cap}}\displaylimits}%
  \def\lesseqqgtr{\mathrel{\lessgtr}}\def\bigoplus{\mathop{\oplus}\displaylimits}%
}
\providecommand{\ssi}{if and only if} % abréviation fréquente
\AtBeginDocument{\providecommand{\wideparen}{\overparen}} % arc de cercle

% Fonctions usuelles que les modèles utilisent sans que amsmath les fournisse
\providecommand{\arsinh}{\operatorname{arsinh}}\providecommand{\arcosh}{\operatorname{arcosh}}
\providecommand{\artanh}{\operatorname{artanh}}\providecommand{\arsech}{\operatorname{arsech}}
\providecommand{\arcsch}{\operatorname{arcsch}}\providecommand{\arcoth}{\operatorname{arcoth}}
\providecommand{\arcsinh}{\operatorname{arsinh}}\providecommand{\arccosh}{\operatorname{arcosh}}
\providecommand{\arctanh}{\operatorname{artanh}}\providecommand{\sech}{\operatorname{sech}}
\providecommand{\csch}{\operatorname{csch}}\providecommand{\arcsec}{\operatorname{arcsec}}
\providecommand{\arccsc}{\operatorname{arccsc}}\providecommand{\arccot}{\operatorname{arccot}}
\providecommand{\sgn}{\operatorname{sgn}}\providecommand{\lcm}{\operatorname{lcm}}
\providecommand{\Var}{\operatorname{Var}}\providecommand{\Cov}{\operatorname{Cov}}

%% FIGFIX-BEGIN v5 — correctifs globaux des figures (docs/onbuch-generation/tools/figfix)
\usepackage{adjustbox}\usepackage{etoolbox}\tcbuselibrary{raster}
% toute figure TikZ/pgfplots plus large que la ligne est réduite à la largeur disponible
\BeforeBeginEnvironment{tikzpicture}{\begin{adjustbox}{max width=\linewidth}}
\AfterEndEnvironment{tikzpicture}{\end{adjustbox}}
% légendes pgfplots lisibles par-dessus les courbes
\pgfplotsset{every axis/.append style={legend style={fill=white,fill opacity=0.92,text opacity=1,draw=black!15,inner sep=3pt}}}
% étiquettes d'arêtes (syntaxe edge["..."]) avec fond blanc
\tikzset{every edge quotes/.append style={fill=white,inner sep=1.2pt}}
%% FIGFIX-END
\begin{document}

\pagegarde{PALAEONTOLOGY}{Geology}{Upper Sixth}{Upper Sixth — Geology}{2}{27 periods}{This chapter introduces palaeontology, the study of ancient life through fossils. Students explore the types of fossils, the conditions and modes of fossilization, the uses and limits of the fossil record, and learn to describe and classify the fossil groups most commonly preserved in sedimentary rocks.
\vspace{4pt}
\begin{multicols}{2}\small
\begin{itemize}
\item By the end of this chapter I can define palaeontology and state its scope and branches
\item By the end of this chapter I can distinguish body fossils, trace fossils and chemical fossils with examples
\item By the end of this chapter I can explain the conditions necessary for fossilization
\item By the end of this chapter I can describe the main modes of fossilization (permineralization, moulds and casts, carbonization, etc.)
\item By the end of this chapter I can assess the preservation potential of different organisms and environments
\item By the end of this chapter I can explain how fossils occur in rocks and how they are collected
\end{itemize}\end{multicols}}

\begin{sommaire}
\begin{multicols}{2}
\begin{enumerate}
\item Scope and Branches of Palaeontology
\item Types of Fossils I: Body Fossils
\item Types of Fossils II: Trace Fossils and Chemical Fossils
\item Conditions for Fossilization
\item Modes of Fossilization I: Unaltered Remains and Replacement Processes
\item Modes of Fossilization II: Moulds, Casts, Carbonization and Imprints
\item Preservation Potential and Occurrence of Fossils
\item Uses of Fossils, Gaps in the Fossil Record and Extinction
\item Classification of Fossils
\item Description of Commonly Preserved Fossils I: Microfossils, Sponges, Corals and Graptolites
\item Description of Commonly Preserved Fossils II: Molluscs, Brachiopods, Arthropods, Echinoderms, Vertebrates and Plants (Integration)
\item Integration activity
\item Methods and worked examples
\item Exercises
\item Solutions to the exercises
\item Summary sheet
\end{enumerate}
\end{multicols}
\end{sommaire}

\begin{objectifs}
\begin{itemize}
\item By the end of this chapter I can define palaeontology and state its scope and branches
\item By the end of this chapter I can distinguish body fossils, trace fossils and chemical fossils with examples
\item By the end of this chapter I can explain the conditions necessary for fossilization
\item By the end of this chapter I can describe the main modes of fossilization (permineralization, moulds and casts, carbonization, etc.)
\item By the end of this chapter I can assess the preservation potential of different organisms and environments
\item By the end of this chapter I can explain how fossils occur in rocks and how they are collected
\item By the end of this chapter I can state the geological, biological and economic uses of fossils
\item By the end of this chapter I can account for gaps in the fossil record and causes of extinction
\item By the end of this chapter I can classify fossils and describe at least fourteen commonly preserved fossil groups
\end{itemize}
\end{objectifs}

\begin{prerequis}
\begin{itemize}
\item Sedimentary rocks and their formation (Lower Sixth: sedimentology)
\item Stratigraphy: principles of superposition and relative dating (Lower Sixth)
\item The geological time scale and major eras (Lower Sixth)
\item Basic biology: cell structure, hard and soft body parts of organisms
\item Rock cycle and diagenesis (Lower Sixth)
\end{itemize}
\end{prerequis}

\newpage

\coursec{Scope and Branches of Palaeontology}

In 2019, workers quarrying limestone near Figuil in the North Region of Cameroon uncovered strange coiled shells embedded in the rock, hundreds of kilometres from the sea. How did marine creatures end up in a dry savanna landscape? Who were these organisms, how did their remains survive for millions of years, and what can they tell us about the ancient environments of Cameroon? Palaeontology, the science of fossils, holds the answers --- and it is the key geologists use to date rocks, reconstruct past climates and locate mineral resources.

This opening section defines palaeontology, fixes its vocabulary, explores its scope and limits, presents its main branches, and situates it within the Earth sciences. Everything else in this chapter --- types of fossils, fossilization, classification, description of fossil groups --- rests on these foundations.

\courssub{What is Palaeontology?}

\textbf{Intuition first.} Pick up a limestone pebble on any Cameroonian hill and split it: you may find the delicate imprint of a shell that no living animal has carried for a hundred million years. That imprint is a message from the deep past. Reading such messages --- identifying the organism, reconstructing how it lived, and deducing the environment in which it died --- is the daily work of the palaeontologist.

\begin{definition}[Palaeontology]
\cle{Palaeontology} is the scientific study of ancient life on Earth through the examination of \cle{fossils} preserved in sedimentary rocks (and exceptionally in other materials such as amber or ice). It seeks to identify past organisms, reconstruct their anatomy, mode of life and evolution, and interpret the environments and climates in which they lived.
\end{definition}

The word itself tells the story. It is built from three Greek roots:

\begin{center}
\begin{tabularx}{\linewidth}{|l|l|X|}
\hline
\rowcolor{popblueL} \textbf{Greek root} & \textbf{Meaning} & \textbf{Contribution to the word} \\
\hline
\textit{palaios} & ancient, old & the time dimension: life of the geological past \\
\hline
\textit{onta} (from \textit{on, ontos}) & beings, things that exist & the object: living organisms \\
\hline
\textit{logos} & study, discourse, science & the method: rational, scientific investigation \\
\hline
\end{tabularx}
\end{center}

Palaeontology is therefore literally ``the study of ancient beings''. The term was popularised in the early nineteenth century, when the growing collections of strange petrified shells and bones forced naturalists to recognise that the Earth had hosted a succession of living worlds very different from the present one.

\begin{definition}[Fossil]
A \cle{fossil} is any preserved evidence of past life --- remains of an organism (shell, bone, leaf), traces of its activity (footprints, burrows, coprolites) or chemical signatures of its presence --- found in rocks or sediments. By convention, remains are generally considered fossils when they are older than about $10\,000\ \mathrm{years}$, i.e.\ older than the end of the last glacial period; more recent remains are called \emph{subfossils}.
\end{definition}

\begin{attention}[Palaeontology is not archaeology]
A very common mistake at the GCE is to confuse \cle{palaeontology} with \cle{archaeology}. \textbf{Archaeology} studies \emph{human} artefacts and structures (pottery, tools, buildings, rock engravings) to reconstruct past human societies. \textbf{Palaeontology} studies the remains of \emph{ancient organisms} --- animals, plants, micro-organisms --- preserved in rocks. A fossilised dinosaur bone is palaeontology; a carved stone tool from an ancient village is archaeology. The two sciences meet only in the study of early humans and their fossil relatives (palaeoanthropology).
\end{attention}

\courssub{The Scope of Palaeontology: What Fossils Can and Cannot Tell Us}

\textbf{What palaeontology can tell us.} From fossils, palaeontologists can:

\begin{itemize}
\item \textbf{Identify extinct organisms} and describe their anatomy, size and growth;
\item \textbf{Reconstruct evolution}, showing how groups appeared, diversified and disappeared through geological time;
\item \textbf{Date rocks} (biostratigraphy): characteristic fossils allow sedimentary layers to be correlated from one region to another;
\item \textbf{Reconstruct past environments and climates} (palaeoecology, palaeoclimatology): marine shells at Figuil prove that the sea once covered the North Region of Cameroon;
\item \textbf{Reconstruct past geography}: identical fossils on both sides of the Atlantic helped to demonstrate that Africa and South America were once joined;
\item \textbf{Guide the search for resources}: microfossils date and correlate the sedimentary formations that host petroleum in the Douala and Rio del Rey basins.
\end{itemize}

\textbf{What palaeontology cannot tell us.} The fossil record is incomplete and biased, and honest scientists must state its limits:

\begin{itemize}
\item Soft-bodied organisms (jellyfish, worms, most internal organs) are almost never preserved, so whole branches of ancient life are invisible;
\item Colour, sounds, behaviour and most physiological details are lost, except where trace fossils or exceptional preservation give indirect clues;
\item The record is strongly biased towards marine organisms with hard parts (shells, skeletons) living in low-energy environments where burial is rapid;
\item Erosion, metamorphism and tectonic deformation have destroyed countless fossils, so the record is a sampling, not a complete archive.
\end{itemize}

\begin{aretenir}[Scope of palaeontology]
Palaeontology reveals the \emph{identity}, \emph{anatomy}, \emph{evolution}, \emph{age} and \emph{environment} of past life, but only for the small fraction of organisms that were fossilized and later discovered. Absence of a fossil is never proof of the absence of the organism.
\end{aretenir}

\courssub{The Branches of Palaeontology}

Because ancient life was as diverse as modern life, palaeontology is divided into specialised branches, each defined by the kind of fossil it studies.

\begin{definition}[Main branches of palaeontology]
\begin{itemize}
\item \cle{Palaeozoology}: the study of fossil \textbf{animals} (invertebrates such as ammonites and trilobites, and vertebrates such as fishes, dinosaurs and mammals);
\item \cle{Palaeobotany}: the study of fossil \textbf{plants} (leaves, wood, seeds, flowers);
\item \cle{Micropalaeontology}: the study of \textbf{microfossils} --- fossils so small that they must be studied under the microscope (foraminifera, ostracods, radiolarians, conodonts);
\item \cle{Palynology}: the study of fossil \textbf{pollen grains and spores} (and other resistant organic micro-remains such as dinoflagellate cysts);
\item \cle{Ichnology}: the study of \textbf{trace fossils} --- footprints, trackways, burrows, borings and coprolites --- which record the \emph{behaviour} of organisms rather than their bodies.
\end{itemize}
\end{definition}

\begin{popfigure}
\begin{tikzpicture}[
  grow=right,
  level 1/.style={sibling distance=2.6cm, level distance=3.4cm},
  level 2/.style={sibling distance=1.15cm, level distance=3.6cm},
  every node/.style={font=\small},
  edge from parent/.style={draw=popmuted, thick}
]
\node[draw=popdark, fill=popblue, text=white, rounded corners, align=center, font=\bfseries, text width=2.2cm] {Palaeontology}
  child { node[draw=popink, fill=popblueL, rounded corners, align=center, text width=2.2cm] {\textbf{Ichnology}\\ trace fossils}
    child { node[draw=popmuted, rounded corners, align=center, fill=white, text width=2cm] {footprints,\\ burrows, coprolites} } }
  child { node[draw=popink, fill=popblueL, rounded corners, align=center, text width=2.2cm] {\textbf{Palynology}\\ pollen \& spores}
    child { node[draw=popmuted, rounded corners, align=center, fill=white, text width=2cm] {pollen, spores,\\ cysts} } }
  child { node[draw=popink, fill=popblueL, rounded corners, align=center, text width=2.2cm] {\textbf{Micropalaeontology}\\ microfossils}
    child { node[draw=popmuted, rounded corners, align=center, fill=white, text width=2cm] {foraminifera,\\ ostracods} } }
  child { node[draw=popink, fill=popblueL, rounded corners, align=center, text width=2.2cm] {\textbf{Palaeobotany}\\ fossil plants}
    child { node[draw=popmuted, rounded corners, align=center, fill=white, text width=2cm] {leaves, wood,\\ seeds} } }
  child { node[draw=popink, fill=popblueL, rounded corners, align=center, text width=2.2cm] {\textbf{Palaeozoology}\\ fossil animals}
    child { node[draw=popmuted, rounded corners, align=center, fill=white, text width=2cm] {ammonites,\\ trilobites, bones} } };
\end{tikzpicture}
\legende{Concept map of the main branches of palaeontology and the objects each branch studies.}
\end{popfigure}

\courssub{Palaeontology and the Other Sciences}

Palaeontology sits at the crossroads of geology and biology, and it is the backbone of stratigraphy.

\begin{itemize}
\item \textbf{With biology:} fossils are studied with the tools of biology --- anatomy, taxonomy, ecology, genetics of living relatives. Palaeontology provides biology with the dimension of \emph{deep time}: it is the only direct evidence of evolution and extinction.
\item \textbf{With geology:} fossils occur almost exclusively in sedimentary rocks, so the palaeontologist must understand sedimentation, diagenesis and tectonics. Conversely, fossils give geology its clock.
\item \textbf{With stratigraphy:} the subdivision of stratigraphy called \cle{biostratigraphy} uses \emph{index fossils} --- species that are widespread, abundant, easily recognised and short-lived --- to define biozones and to correlate rock layers between distant outcrops. A foraminiferan assemblage in a borehole near Douala can be matched with the same assemblage in Nigeria or Gabon, proving the layers are of the same age.
\end{itemize}

\courssub{A Brief History of Palaeontology}

\textbf{From curiosity to science.} For centuries, fossils were collected as curiosities. Greek and Roman writers already noticed seashells on mountains, and some, like Xenophanes, correctly suggested that the land had once been covered by the sea. During the Renaissance, Leonardo da Vinci argued that fossil shells were the remains of real marine organisms, not ``sports of nature''. In the seventeenth century, Nicolas Steno showed that ``tongue stones'' were fossil shark teeth and formulated the principles of stratigraphy that allow fossils to be placed in a time sequence.

\textbf{The founding breakthrough: extinction.} Around 1800, the French anatomist \cle{Georges Cuvier} compared fossil elephant-like bones (mammoths, mastodons) with living elephants and proved they belonged to species that \emph{no longer exist}. This established the reality of \cle{extinction} --- a revolutionary idea --- and founded comparative anatomy and vertebrate palaeontology. Cuvier also showed that successive rock layers contain different fossil faunas, separated by sudden disappearances which he attributed to catastrophes.

\textbf{The nineteenth-century synthesis.} William Smith in England used fossils to correlate strata and drew the first geological map; Mary Anning's discoveries of marine reptiles on the Dorset coast transformed knowledge of Mesozoic seas; Charles Darwin, in 1859, placed fossils at the heart of the theory of evolution by natural selection, interpreting the fossil record as a history of descent with modification. In the twentieth century, micropalaeontology and palynology became industrial sciences through petroleum exploration, and radiometric dating gave the fossil record an absolute time scale.

\begin{popfigure}
\begin{tikzpicture}[scale=1, transform shape]
\draw[very thick, popdark, ->] (0,0) -- (13.2,0) node[right, font=\small\bfseries] {time};
\foreach \x/\yr in {0.4/{c. 500 BC}, 2.6/{1500}, 4.8/{1669}, 7.0/{1812}, 9.2/{1859}, 11.4/{1900s}}{
  \draw[popdark] (\x,0.12) -- (\x,-0.12);
  \node[below, font=\scriptsize, align=center] at (\x,-0.15) {\yr};
}
\node[above, font=\scriptsize, align=center, text width=2.0cm] at (0.4,0.2) {Xenophanes: shells on mountains};
\node[above, font=\scriptsize, align=center, text width=2.0cm] at (2.6,0.2) {Leonardo: fossils are real organisms};
\node[above, font=\scriptsize, align=center, text width=2.0cm] at (4.8,0.2) {Steno: shark teeth \& strata};
\node[above, font=\scriptsize, align=center, text width=2.2cm] at (7.0,0.2) {Cuvier: extinction proven};
\node[above, font=\scriptsize, align=center, text width=2.0cm] at (9.2,0.2) {Darwin: evolution by selection};
\node[above, font=\scriptsize, align=center, text width=2.2cm] at (11.4,0.2) {Micropalaeontology \& oil industry};
\end{tikzpicture}
\legende{Timeline of some key steps in the history of palaeontology.}
\end{popfigure}

\courssub{How a Palaeontologist Works}

\begin{methode}[From field collection to laboratory identification]
\begin{enumerate}
\item \textbf{Prospecting:} locate fossiliferous outcrops using geological maps; sedimentary rocks (limestone, shale, sandstone) are the targets.
\item \textbf{Collection:} extract specimens carefully with hammer, chisel and brush; wrap fragile fossils in paper or plaster jackets.
\item \textbf{Recording:} note the exact locality, the stratigraphic layer, the orientation of the fossil and associated sediments --- a fossil without its context loses most of its scientific value.
\item \textbf{Preparation:} in the laboratory, clean the fossil mechanically (needles, vibro-tools) or chemically (dilute acid for limestones); sieve and wash sediments to recover microfossils.
\item \textbf{Identification:} compare the specimen with reference collections and published descriptions; measure diagnostic characters; assign it to a taxonomic group (genus, species).
\item \textbf{Interpretation:} deduce the age of the rock, the palaeoenvironment (marine, lagoonal, continental), and the palaeoclimate; publish and store the specimen in a curated collection.
\end{enumerate}
\end{methode}

\courssub{Palaeontology in Cameroon}

\begin{savaistu}[Fossils at the service of Cameroon]
Cameroon possesses a rich and varied fossil heritage. The Cretaceous limestones of \cle{Figuil} (North Region), exploited by the cement industry, contain marine fossils --- ammonites, bivalves, gastropods --- proving that a shallow sea extended far into the continent at that time. In the \cle{Douala basin}, micropalaeontologists study foraminifera and ostracods from borehole samples to date and correlate sedimentary formations: this biostratigraphic work is an essential tool of \textbf{petroleum exploration} in the Douala and Rio del Rey basins. Palynology (fossil pollen and spores) is likewise used to characterise source rocks and reconstruct the ancient deltas that fed the basins. Palaeontology is thus not an abstract science in Cameroon: it contributes directly to the cement industry, to oil exploration and to the understanding of the country's geological history.
\end{savaistu}

\begin{exoresolu}[Applying the definition: fossil or not?]
\textbf{Statement.} A student finds the following objects: (a) a coiled ammonite shell in Cretaceous limestone at Figuil; (b) a broken clay pot in an abandoned village site; (c) a bird that died last year, dried in the sun on a rock; (d) a network of fossil burrows in a sandstone bed. For each object, state whether it is a fossil and which science studies it.
\tcblower
\textbf{Solution.}
\begin{itemize}
\item (a) The ammonite is the preserved remain of an ancient marine animal, millions of years old, embedded in rock: it is a \textbf{body fossil}, studied by \textbf{palaeozoology}.
\item (b) The clay pot is a human artefact: it is \textbf{not a fossil}; it belongs to \textbf{archaeology}.
\item (c) The dried bird is far younger than the conventional limit of about $10\,000\ \mathrm{years}$ and is not incorporated in sediment: it is a \textbf{recent remain (subfossil at most)}, not a true fossil.
\item (d) The burrows record the \emph{activity} of an ancient organism, not its body: they are \textbf{trace fossils}, studied by \textbf{ichnology}.
\end{itemize}
\end{exoresolu}

\begin{exercice}[Checking your understanding]
\begin{enumerate}
\item Give the etymology of the word ``palaeontology'' and explain how each Greek root contributes to its meaning.
\item State two things palaeontology can reveal about the past and two things it cannot, justifying each answer.
\item A borehole near Douala yields abundant microscopic shells of foraminifera. Which branch of palaeontology studies them, and why is this study economically important for Cameroon?
\item Explain in three sentences why Cuvier's work is considered a turning point in the history of palaeontology.
\end{enumerate}
\end{exercice}

\begin{corrige}[Exercise 1]
\begin{enumerate}
\item From \textit{palaios} (ancient), \textit{onta} (beings) and \textit{logos} (study): the scientific study of ancient beings, i.e.\ of past life through fossils.
\item \emph{Can:} identify extinct organisms and reconstruct evolution; date and correlate rocks and reconstruct palaeoenvironments (e.g.\ marine fossils at Figuil prove an ancient sea). \emph{Cannot:} preserve soft-bodied organisms and details such as colour or behaviour; give a complete record, because fossilization is rare and biased towards hard parts and rapid burial.
\item Micropalaeontology. Foraminifera are excellent index fossils: they allow the sedimentary formations of the Douala basin to be dated and correlated, which guides petroleum exploration.
\item Cuvier compared fossil bones with the anatomy of living animals and demonstrated that some fossil species no longer exist, establishing the reality of extinction. He thereby founded vertebrate palaeontology and comparative anatomy, and showed that successive strata contain different, successive faunas.
\end{enumerate}
\end{corrige}

\coursec{Types of Fossils I: Body Fossils}

In the previous section we defined palaeontology as the study of ancient life through its preserved remains, and we saw that the discipline branches into palaeozoology, palaeobotany and micropalaeontology. But what exactly does a palaeontologist study in the field? The answer is: \cle{fossils}. Yet not all fossils are of the same kind. A shell embedded in limestone, a footprint pressed into an ancient mudflat, and a thin film of organic molecules in shale are all fossils — but they record life in very different ways. In this section and the next, we establish the general classification of fossils, and then examine in depth the first and most familiar category: the \cle{body fossils}.

\courssub{General Classification of Fossils}

\begin{definition}[Fossil]
A \cle{fossil} is any preserved remain, impression, or trace of an organism that lived in the geological past, typically older than about 10\,000 years (the conventional Pleistocene–Holocene boundary), found in sedimentary rocks or other natural deposits (ice, amber, tar, peat).
\end{definition}

Fossils are classified into three great categories according to \emph{what} is preserved:

\begin{enumerate}
\item \cle{Body fossils}: the actual remains of the organism itself, or parts of it (a shell, a bone, a leaf).
\item \cle{Trace fossils} (ichnofossils): evidence of the \emph{activity} of an organism — footprints, burrows, borings, coprolites (fossil faeces) — without the body itself.
\item \cle{Chemical fossils} (biomarkers): organic molecules or distinctive chemical signatures left in rocks by living organisms, even when no visible structure survives.
\end{enumerate}

\begin{popfigure}
\begin{center}
\begin{tikzpicture}[
grow=down, level distance=1.6cm,
level 1/.style={sibling distance=4.6cm},
level 2/.style={sibling distance=2.2cm},
every node/.style={draw=popline, rounded corners=2pt, fill=popblueL, text=popink, font=\small, align=center},
edge from parent/.style={draw=popmuted, thick}]
\node[fill=popblue, text=white, font=\small\bfseries] {FOSSILS}
child { node[align=center]{Body fossils\\ \emph{remains of the body}}
  child { node[fill=white, align=center]{Hard parts\\ shells, bones, teeth} }
  child { node[fill=white, align=center]{Soft parts\\ (rare: amber, ice)} } }
child { node[align=center]{Trace fossils\\ \emph{activity}}
  child { node[fill=white] {Tracks, burrows} }
  child { node[fill=white] {Coprolites, borings} } }
child { node[align=center]{Chemical fossils\\ \emph{molecules}}
  child { node[fill=white] {Biomarkers} }
  child { node[fill=white] {Kerogen, isotopes} } };
\end{tikzpicture}
\end{center}
\legende{General classification of fossils into body fossils, trace fossils and chemical fossils.}
\end{popfigure}

This section is devoted to body fossils; trace and chemical fossils are treated in the next section.

\courssub{Body Fossils: Definition and Nature}

\begin{definition}[Body fossil]
A \cle{body fossil} is a fossil consisting of the actual remains of an organism's body, or of a part of it — such as a shell, bone, tooth, exoskeleton, piece of wood, spore or leaf — preserved in the rock record.
\end{definition}

The intuitive idea is simple: when an organism dies, its body may escape total destruction and become incorporated into sediment. What survives depends overwhelmingly on the \emph{composition} of the body part. Organisms are made of two broad categories of tissues:

\begin{itemize}
\item \cle{Hard parts}: mineralised or resistant structures — shells, bones, teeth, exoskeletons, wood, spores and pollen walls.
\item \cle{Soft parts}: muscles, skin, internal organs, blood, soft leaves and flowers.
\end{itemize}

\courssub{Hard Parts Commonly Preserved}

\begin{propriete}[Resistance of hard parts to decay]
Hard parts resist decay and mechanical destruction far better than soft tissues because of their composition: \cle{calcite} and \cle{aragonite} (\ce{CaCO3}) in mollusc and brachiopod shells and corals; \cle{silica} (\ce{SiO2}) in diatom and radiolarian skeletons and some sponges; \cle{calcium phosphate} (apatite, \ce{Ca5(PO4)3(OH)}) in vertebrate bones and teeth; and \cle{chitin} (a tough organic polymer) in arthropod exoskeletons. Soft tissues, made mainly of proteins, fats and water, are rapidly destroyed by bacteria, scavengers and oxidation, and are preserved only under exceptional conditions.
\end{propriete}

The principal hard parts encountered as body fossils are:

\begin{center}
\begin{tabularx}{\linewidth}{|l|X|l|}
\hline
\rowcolor{popblueL} \textbf{Hard part} & \textbf{Composition} & \textbf{Typical fossil examples} \\
\hline
Shells & Calcite or aragonite (\ce{CaCO3}) & Ammonites, bivalves, gastropods \\
\hline
Bones & Calcium phosphate (apatite) & Vertebrate skeletons \\
\hline
Teeth & Enamel (highly mineralised apatite) & Shark teeth, mammal molars \\
\hline
Exoskeletons & Chitin, often calcified & Trilobites, crabs, insects \\
\hline
Wood & Cellulose and lignin & Petrified tree trunks \\
\hline
Spores and pollen & Sporopollenin (extremely resistant) & Fossil pollen in shales \\
\hline
\end{tabularx}
\end{center}

Among these, \cle{teeth} deserve special mention: tooth enamel is the most highly mineralised substance produced by any organism (up to 96\% inorganic apatite), which is why teeth are often the \emph{only} remains of vertebrates in many deposits. Likewise, \cle{sporopollenin}, the substance of pollen and spore walls, is one of the most decay-resistant organic materials known, which explains the abundance of fossil pollen in fine-grained sediments.

\courssub{Soft Parts: Rare but Spectacular Preservation}

Soft parts normally disappear within days or weeks of death. Their preservation requires that decay be stopped almost immediately. Three exceptional situations allow this:

\begin{enumerate}
\item \cle{Mummification}: in extremely dry climates, rapid desiccation removes the water that bacteria need, so skin and dried muscles may survive (e.g. mummified mammals in desert caves).
\item \cle{Freezing}: in permanently frozen ground (permafrost), decay is halted by cold; entire woolly mammoths with skin, flesh and even stomach contents have been recovered from Siberian ice.
\item \cle{Amber}: insects or small organisms trapped in sticky tree resin are sealed away from air, water and bacteria; the resin hardens into amber, preserving delicate bodies in three dimensions, sometimes with fine hairs and wing veins intact.
\end{enumerate}

\begin{savaistu}[Amber: a window on ancient life]
Amber is fossilised tree resin. Because it entombs organisms whole and excludes oxygen and microbes, amber has preserved insects, spiders, lizards and even feathers in exquisite detail for tens of millions of years. Such fossils are body fossils of soft or delicate organisms that would otherwise leave no record at all.
\end{savaistu}

\courssub{Unaltered Original Material}

In some body fossils, the \emph{original material} survives essentially unchanged — no mineral replacement has occurred. This happens when the material is naturally resistant and burial conditions are gentle:

\begin{itemize}
\item \cle{Tooth enamel} and dense bone often retain their original calcium phosphate.
\item \cle{Shells in young (recent) deposits} may still consist of original aragonite or calcite, sometimes retaining traces of original colour patterns.
\item \cle{Insects in amber} retain their original chitinous exoskeletons and even internal tissues.
\item \cle{Pollen and spores} retain their original sporopollenin walls.
\end{itemize}

Unaltered remains are precious to scientists because the original chemistry — including isotopic signatures of diet and climate — is still readable.

\courssub{Typical Examples of Body Fossils}

\begin{exemplebox}[Classic body fossils]
\begin{itemize}
\item \cle{Ammonite shells}: coiled, chambered aragonitic shells of extinct cephalopods, extremely common in Mesozoic marine limestones.
\item \cle{Trilobite exoskeletons}: segmented, calcified chitinous shields of extinct marine arthropods, characteristic of Palaeozoic rocks.
\item \cle{Vertebrate bones}: limb bones, vertebrae, skulls and teeth of fishes, reptiles and mammals.
\item \cle{Plant leaves}: fronds and leaves preserved as compressions or impressions in fine shales.
\end{itemize}
\end{exemplebox}

\begin{popfigure}
\begin{center}
\begin{tikzpicture}[scale=1.0]
% ---- Ammonite (left) ----
\begin{scope}[shift={(-3.4,0)}]
\draw[thick, popink] (0,0) circle (1.7);
\draw[thick, popink] (0,0) circle (0.35);
\foreach \a in {0,45,...,315} {
  \draw[popmuted] (\a:0.35) -- (\a:1.7);
}
\draw[popblue, thick] (0:0.35) arc (0:270:0.35 and 0.35);
\draw[popblue, thick, domain=0:300, samples=60] plot ({ (0.35+0.0045*\x)*cos(-\x) }, { (0.35+0.0045*\x)*sin(-\x) });
\node[font=\small\bfseries, popink] at (0,-2.2) {Ammonite shell};
\draw[->, popmuted] (1.9,1.2) -- (1.15,1.15) node[pos=0, right, font=\scriptsize, align=left, text=popink] {ribbed shell wall};
\draw[->, popmuted] (-2.6,0.6) -- (-1.2,0.5) node[pos=0, left, font=\scriptsize, align=right, text=popink] {septa between\\ chambers};
\draw[->, popmuted] (-1.6,-1.5) -- (-0.25,-0.25) node[pos=0, left, font=\scriptsize, align=right, text=popink] {coiled whorls};
\end{scope}
% ---- Trilobite (right) ----
\begin{scope}[shift={(3.6,0)}]
\draw[thick, popink] (0,2.0) arc (90:270:0.9 and 0.55);
\draw[thick, popink] (0,2.0) arc (90:-90:0.9 and 0.55);
\draw[thick, popink] (0,1.45) -- (0,-1.9);
\draw[thick, popink] (-0.9,1.45) .. controls (-1.05,0.2) and (-1.05,-1.2) .. (-0.55,-1.9);
\draw[thick, popink] (0.9,1.45) .. controls (1.05,0.2) and (1.05,-1.2) .. (0.55,-1.9);
\draw[thick, popink] (-0.55,-1.9) -- (0.55,-1.9);
\foreach \y in {1.1,0.75,0.4,0.05,-0.3,-0.65,-1.0,-1.35,-1.65} {
  \draw[popmuted] (-0.95,\y) .. controls (0,\y-0.12) .. (0.95,\y);
}
\draw[popblue, thick] (0,1.45) -- (0,-1.9);
\node[font=\small\bfseries, popink] at (0,-2.5) {Trilobite exoskeleton};
\draw[->, popmuted] (2.3,1.7) -- (0.6,1.7) node[pos=0, right, font=\scriptsize, align=left, text=popink] {cephalon (head)};
\draw[->, popmuted] (2.3,0.2) -- (0.9,0.2) node[pos=0, right, font=\scriptsize, align=left, text=popink] {segmented thorax};
\draw[->, popmuted] (2.3,-1.9) -- (0.5,-1.85) node[pos=0, right, font=\scriptsize, align=left, text=popink] {pygidium (tail)};
\draw[->, popmuted] (-2.3,0.6) -- (-0.1,0.6) node[pos=0, left, font=\scriptsize, align=right, text=popink] {axial lobe};
\end{scope}
\end{tikzpicture}
\end{center}
\legende{Two classic body fossils: the coiled, chambered shell of an ammonite (left) and the three-lobed, segmented exoskeleton of a trilobite (right).}
\end{popfigure}

\courssub{Megafossils and Microfossils}

Body fossils are also classified by \emph{size}:

\begin{itemize}
\item \cle{Megafossils} (macrofossils): fossils visible to the naked eye — shells, bones, leaves, trunks. Most museum fossils are megafossils.
\item \cle{Microfossils}: fossils so small that they must be studied with a microscope, typically less than about $1\ \mathrm{mm}$.
\end{itemize}

\begin{definition}[Microfossil]
A \cle{microfossil} is a fossil, or fossil fragment, too small to be studied properly with the naked eye, requiring a microscope. Important groups include \cle{foraminifera} (single-celled organisms with chambered calcareous shells), \cle{ostracods} (tiny bivalved crustaceans), \cle{calcareous nannofossils} (coccoliths), \cle{conodonts} (phosphatic tooth-like elements), and \cle{pollen and spores}.
\end{definition}

Microfossils are of enormous economic importance. Because foraminifera, ostracods and nannofossils evolved rapidly and are abundant in marine sediments, they allow geologists to \emph{date} and \emph{correlate} the rock layers penetrated by boreholes. In \cle{oil exploration}, micropalaeontologists examine drill cuttings to identify the age and depositional environment of subsurface strata — a routine procedure in the petroleum industry, including in the sedimentary basins of the \cle{Rio del Rey} and \cle{Douala} basins of Cameroon, where planktonic foraminiferal zones and calcareous nannofossil bioevents help locate oil-bearing formations.

\courssub{Method: Classifying a Fossil Specimen}

\begin{methode}[How to classify a fossil into body, trace or chemical]
\begin{enumerate}
\item \textbf{Observe the specimen:} Is there a recognisable structure (shell, bone, leaf)? If yes $\rightarrow$ likely a \cle{body fossil}.
\item \textbf{If no body structure:} Does it show evidence of behaviour (track, burrow, boring, coprolite)? If yes $\rightarrow$ \cle{trace fossil}.
\item \textbf{If no visible structure at all:} Does laboratory analysis reveal organic molecules characteristic of life (hopanes, steranes, specific isotope ratios)? If yes $\rightarrow$ \cle{chemical fossil}.
\item \textbf{Check size:} If body fossil, is it visible to naked eye? $\rightarrow$ megafossil. Requires microscope? $\rightarrow$ microfossil.
\end{enumerate}
\end{methode}

\courssub{Worked Example and Common Mistakes}

\begin{exoresolu}[Identifying fossil types]
A student collects three specimens from a shale outcrop: (a) a complete clam shell; (b) a sinuous trail winding across a bedding plane; (c) a dark shale that yields organic molecules characteristic of ancient algae under laboratory analysis. Classify each specimen and justify your answer.
\tcblower
\textbf{Solution.}
\begin{enumerate}
\item The clam shell is a \cle{body fossil}: it is the actual mineralised hard part (calcareous shell) of the organism itself.
\item The sinuous trail is a \cle{trace fossil}: it records the \emph{activity} (locomotion) of an organism, not its body.
\item The algal molecules are \cle{chemical fossils} (biomarkers): no structure is preserved, only the chemical signature of life.
\end{enumerate}
The key test is always: \emph{what is preserved — the body, the behaviour, or the chemistry?}
\end{exoresolu}

\begin{attention}[Common mistake: "all fossils are bones"]
Many students imagine that fossils are mostly dinosaur bones. In reality, \textbf{bones are a small minority} of the fossil record. The overwhelming majority of fossils are \cle{shells} of marine invertebrates (molluscs, brachiopods) and \cle{plant remains}, because these hard parts are abundant, resistant, and produced in environments where burial is likely. Vertebrate skeletons are comparatively rare because vertebrates are far less numerous and often die in upland environments where erosion, rather than burial, dominates.
\end{attention}

\begin{attention}[Common mistake: confusing moulds/casts with body fossils]
A mould or cast is \emph{not} a body fossil — it is a \cle{mode of preservation} of a body fossil. The original body has dissolved away, leaving a cavity (mould) that may later be filled (cast). In the GCE exam, you must distinguish the \emph{type of fossil} (body, trace, chemical) from the \emph{mode of fossilization} (unaltered, replacement, mould/cast, carbonization, etc.).
\end{attention}

\begin{exercice}[Checking your understanding]
\begin{enumerate}
\item Define a body fossil and give three examples of hard parts commonly preserved.
\item Explain, in terms of chemical composition, why teeth and pollen are among the most frequently preserved body fossils.
\item Describe three exceptional circumstances under which soft parts may be fossilised.
\item Distinguish between a megafossil and a microfossil, and state why microfossils are valuable in petroleum exploration.
\item A geologist finds a well-preserved insect in amber and a dinosaur footprint in sandstone. Classify each and explain the difference in preservation.
\end{enumerate}
\end{exercice}

\begin{corrige}[Exercise: Checking your understanding]
\begin{enumerate}
\item A body fossil is the preserved actual remain of an organism's body or part of it. Examples: mollusc shells (calcite/aragonite), vertebrate bones and teeth (calcium phosphate), arthropod exoskeletons (chitin, often calcified) — also wood, spores.
\item Teeth are made of enamel, the most highly mineralised biological material (dense apatite, ~96\% inorganic), so they resist chemical and mechanical breakdown. Pollen walls are made of sporopollenin, an extremely decay-resistant organic polymer, so they survive where other tissues rot away.
\item Mummification (rapid drying in arid conditions), freezing (permafrost halting bacterial decay, e.g. mammoths), and entombment in amber (resin sealing organisms from air, water and microbes).
\item A megafossil is visible to the naked eye (shells, bones, leaves); a microfossil requires a microscope (foraminifera, ostracods, pollen, nannofossils). Microfossils evolve rapidly and are abundant in marine sediments, so they date and correlate subsurface strata in boreholes, guiding the search for oil-bearing rocks.
\item Insect in amber $\rightarrow$ body fossil (soft parts preserved by exclusion of oxygen and microbes). Dinosaur footprint $\rightarrow$ trace fossil (records activity, not the body). The insect preserves the organism itself; the footprint preserves only its behaviour.
\end{enumerate}
\end{corrige}

\begin{aretenir}[Key points]
\begin{itemize}
\item Fossils are classified into \cle{body fossils}, \cle{trace fossils} and \cle{chemical fossils}.
\item Body fossils are actual remains of organisms; hard parts (calcite, aragonite, silica, calcium phosphate, chitin, sporopollenin) are commonly preserved, soft parts only exceptionally (mummification, freezing, amber).
\item Some body fossils retain unaltered original material: tooth enamel, shells in recent deposits, insects in amber.
\item Body fossils range from megafossils (ammonites, trilobites, bones, leaves) to microfossils (foraminifera, ostracods, nannofossils, conodonts, pollen), the latter being essential tools in oil exploration (e.g. Cameroon's Rio del Rey and Douala basins).
\item Most fossils are shells and plant remains — not bones.
\item Distinguish \emph{type of fossil} (body/trace/chemical) from \emph{mode of fossilization} (unaltered/replacement/mould-cast/etc.).
\end{itemize}
\end{aretenir}

\coursec{Types of Fossils II: Trace Fossils and Chemical Fossils}

In the previous section we studied \cle{body fossils}: the preserved remains of the organisms themselves --- shells, bones, teeth, leaves. But imagine walking along the muddy shore of the Wouri estuary at low tide. You may never see the crab that lives there, yet the mud is covered with its footprints and pierced by the holes in which it hides. Long after the crab has died and its soft body has decayed, those marks in the sediment could be buried, hardened into rock, and discovered millions of years later. The rock would contain no part of the crab's body, yet it would tell us that the crab existed, how it moved, and where it lived. Such evidence of \emph{activity} rather than of \emph{anatomy} constitutes the second great category of fossils: \cle{trace fossils}. At an even more fundamental level, life can leave behind nothing but its chemistry --- organic molecules locked in the rock. These are \cle{chemical fossils}. This section explores both.

\courssub{Trace Fossils (Ichnofossils)}

\textbf{Intuitive idea.}\par
A detective does not need to catch the thief to prove a theft: fingerprints, footprints and disturbed objects are enough. In the same way, an organism moving, feeding, resting or dwelling in sediment leaves disturbances that can be preserved. The study of these traces is called \cle{ichnology}, and the traces themselves are named by their own system of \emph{ichnogenera} and \emph{ichnospecies}, independent of the biological classification of their makers.

\begin{definition}[Trace fossil (ichnofossil)]
A \cle{trace fossil} (or \cle{ichnofossil}) is a sedimentary structure produced by the \textbf{activity} of a living organism --- its locomotion, feeding, dwelling or resting --- and preserved in the rock record. It records \emph{behaviour}, not the body of the organism. The same trace may be produced by several different organisms, and one organism may produce several different traces.
\end{definition}

\begin{attention}[Trace $\neq$ body]
A common examination error is to assume that a trace fossil identifies its maker. A \emph{Cruziana} trail, for example, is usually attributed to trilobites, but similar trails can be made by other arthropods. Always state: the trace records \textbf{behaviour}; the identity of the producer is an \emph{interpretation}, not a certainty.
\end{attention}

\courssub{The Main Types of Trace Fossils}

\textbf{1. Tracks and trails (locomotion traces).}\par
\cle{Tracks} are discrete imprints left by feet or appendages --- a single dinosaur footprint is a track, and a series of them forms a \cle{trackway}. \cle{Trails} are continuous grooves or ridges made by an animal crawling, sliding or dragging part of its body, as when a worm or a trilobite ploughs through soft mud. Trackways are remarkably informative: from the spacing, size and depth of successive footprints of a dinosaur, palaeontologists can estimate the animal's hip height, its speed, and whether it walked alone or in a herd.

\textbf{2. Burrows (dwelling and feeding structures).}\par
\cle{Burrows} are tunnels excavated in soft sediment by organisms such as worms, shrimps and bivalves. Simple vertical tubes (e.g.\ the ichnogenus \emph{Skolithos}) are typically \emph{dwelling} burrows of suspension feeders living in high-energy shallow water, where a deep vertical refuge protects the animal from waves and currents. Horizontal or branching, meandering burrows are usually \emph{feeding} structures, in which the animal systematically mines the sediment for organic matter. Burrows may be preserved as \emph{full relief} (the filled tube itself) when the sediment that later filled the burrow differs from the surrounding sediment.

\textbf{3. Borings.}\par
\cle{Borings} are holes drilled or etched into \emph{hard} substrates --- rock, shells or wood --- by organisms such as boring sponges, bivalves and certain gastropods. Some predatory snails drill a neat circular hole through a bivalve shell to reach the flesh inside: such a bored shell is direct fossil evidence of predation.

\textbf{4. Coprolites (fossil faeces).}\par
\begin{definition}[Coprolite]
A \cle{coprolite} is fossilised faecal material. Because faeces contain the indigestible residues of meals --- fish scales, bone fragments, shell pieces, plant fibres, pollen --- coprolites provide direct evidence of the \textbf{diet} of the producer and of the \textbf{food chains} of ancient ecosystems. Phosphatic composition (rich in calcium phosphate) helps faecal matter fossilise rapidly before it disintegrates.
\end{definition}

\textbf{5. Gastroliths.}\par
\cle{Gastroliths} are polished ``stomach stones'' found within the body cavity of some fossil reptiles (and used today by birds and crocodiles). The animal swallowed stones to help grind food in a muscular gizzard. Their smooth, rounded, polished surfaces distinguish them from ordinary pebbles.

\textbf{6. Nests and eggs.}\par
Fossil \cle{eggs} and \cle{nesting structures} --- most famously dinosaur eggs arranged in clutches in excavated hollows --- record reproductive behaviour. The arrangement of eggs, the thickness and porosity of the shell, and the presence of embryos inside give information on parental care and development.

\begin{exemplebox}[Reading a dinosaur trackway]
A trackway shows footprints \SI{1.2}{m} apart (stride length), each \SI{0.6}{m} long. Empirical relations link hip height to roughly four times foot length (here about \SI{2.4}{m}), and stride length to speed. The regular, narrow spacing with no tail-drag mark indicates a large biped walking steadily, not running. If several parallel trackways of the same type occur together, \emph{gregarious} (herd) behaviour is inferred --- information no skeleton alone could give.
\end{exemplebox}

\begin{popfigure}
\begin{tikzpicture}[scale=0.95, font=\small]
  % --- Panel 1: footprint trackway ---
  \fill[popblueL] (0,0) rectangle (4.6,2.6);
  \foreach \x/\y/\a in {0.7/0.5/-15, 1.5/1.1/-10, 2.3/0.6/-15, 3.1/1.2/-10, 3.9/0.7/-15}{
    \begin{scope}[shift={(\x,\y)}, rotate=\a]
      \fill[popdark] (0,0) ellipse (0.22 and 0.32);
      \fill[popdark] (-0.18,0.38) circle (0.09);
      \fill[popdark] (0,0.44) circle (0.09);
      \fill[popdark] (0.18,0.38) circle (0.09);
    \end{scope}}
  \draw[->, poporange, thick] (0.4,2.3) -- (4.2,2.3);
  \node[popdark, anchor=west] at (0.1,2.85) {\textbf{Trackway} (locomotion)};
  % --- Panel 2: vertical burrow (cross-section) ---
  \fill[popblueL] (5.4,0) rectangle (8.2,2.6);
  \draw[popline, thick] (5.4,2.2) -- (8.2,2.2);
  \node[popmuted, anchor=west] at (5.45,2.42) {sediment surface};
  \fill[popgold] (6.6,0.3) -- (6.6,2.2) .. controls (6.6,2.35) and (7.0,2.35) .. (7.0,2.2) -- (7.0,0.3) .. controls (7.0,0.15) and (6.6,0.15) .. (6.6,0.3);
  \draw[popdark] (6.6,0.3) -- (6.6,2.2) (7.0,0.3) -- (7.0,2.2);
  \draw[->, popdark] (7.9,1.2) -- (7.15,1.2);
  \node[popdark, anchor=west, align=left] at (7.0,0.55) {sand fill};
  \node[popdark, anchor=west] at (5.45,2.85) {\textbf{Vertical burrow} (dwelling)};
  % --- Panel 3: feeding trail ---
  \fill[popblueL] (9.0,0) rectangle (13.6,2.6);
  \draw[popgreen, very thick] (9.3,0.5) .. controls (9.8,1.6) and (10.2,0.4) .. (10.7,1.3)
    .. controls (11.2,2.2) and (11.6,0.6) .. (12.1,1.5)
    .. controls (12.6,2.3) and (13.0,0.9) .. (13.3,1.7);
  \draw[->, popgreen, very thick] (13.3,1.7) -- (13.45,1.95);
  \node[popdark, anchor=west] at (9.05,2.85) {\textbf{Feeding trail} (grazing)};
  \node[popmuted, anchor=west] at (9.05,0.15) {meandering path avoids re-crossing};
\end{tikzpicture}
\legende{Three common trace fossils: a footprint trackway (locomotion), a vertical dwelling burrow seen in cross-section with contrasting sand fill, and a meandering feeding trail on a bedding surface.}
\end{popfigure}

\courssub{What Trace Fossils Tell Us}

Trace fossils are among the most powerful tools of the palaeoenvironmentalist:

\begin{itemize}
  \item \textbf{Behaviour and locomotion:} gait, speed, solitary versus herd behaviour, predation (borings), feeding strategies.
  \item \textbf{Water depth and energy:} vertical dwelling burrows (\emph{Skolithos}) dominate high-energy shallow sands; complex horizontal feeding trails characterise quiet, deeper, food-rich muds. Assemblages of traces recurring in particular environments define \emph{ichnofacies}.
  \item \textbf{Oxygen levels:} abundant, large, diverse burrows indicate well-oxygenated bottom water and sediment; anoxic muds are finely laminated and \emph{unburrowed}, because no burrowers could live in them.
  \item \textbf{Soft-bodied life:} worms and jellyfish almost never fossilise as bodies, but their trails and burrows are common --- traces extend the record of life far beyond organisms with hard parts.
\end{itemize}

\begin{propriete}[Trace fossils are facies fossils]
Trace fossils are typically \cle{facies fossils}: because a given behaviour (dwelling, grazing) is adapted to a particular \emph{environment} and can persist with little change for hundreds of millions of years, traces are generally \textbf{poor index fossils for dating} but \textbf{excellent indicators of environment} (water depth, energy, oxygenation, substrate consistency). \emph{(The reasoning --- behaviour is environment-controlled and evolves slowly --- should be reproducible in an examination answer.)}
\end{propriete}

\begin{aretenir}[Advantages of trace fossils]
\begin{itemize}
  \item They \textbf{cannot be transported}: a burrow or trackway is fixed in the sediment where the animal lived, so it records the \emph{in situ} environment (unlike shells, which currents may carry far away).
  \item They record \textbf{soft-bodied organisms} otherwise invisible in the fossil record.
  \item They record \textbf{behaviour directly} --- something body fossils can only suggest.
  \item They can be abundant where body fossils are completely absent (e.g.\ in some sandstones).
\end{itemize}
\end{aretenir}

\courssub{Stromatolites: Organo-Sedimentary Structures}

\cle{Stromatolites} are layered, dome-shaped or columnar structures built by communities of micro-organisms, especially \emph{cyanobacteria} (blue-green algae). Living in shallow sunlit water, the microbes form sticky mats that trap and bind fine sediment; the mat then grows up through the new sediment layer, and the cycle repeats, producing fine \emph{lamination}. Stromatolites are among the oldest evidence of life on Earth and were major reef-builders in Precambrian seas.

\begin{popfigure}
\begin{tikzpicture}[scale=1.0, font=\small]
  % substrate
  \fill[popblueL] (0,0) rectangle (12,0.7);
  \node[popmuted] at (6,0.35) {sea-floor sediment};
  % stromatolite domes
  \foreach \xc/\s in {2/1.0, 5.2/1.3, 8.8/0.9, 11/1.1}{
    \begin{scope}[shift={(\xc,0.7)}, scale=\s]
      \fill[popgreenL] (-1.2,0) .. controls (-1.2,1.4) and (-0.7,1.9) .. (0,1.9)
        .. controls (0.7,1.9) and (1.2,1.4) .. (1.2,0) -- cycle;
      \draw[popgreen, thick] (-1.0,0.25) .. controls (-1.0,1.2) and (-0.55,1.6) .. (0,1.6)
        .. controls (0.55,1.6) and (1.0,1.2) .. (1.0,0.25);
      \draw[popgreen, thick] (-0.75,0.5) .. controls (-0.75,1.1) and (-0.4,1.3) .. (0,1.3)
        .. controls (0.4,1.3) and (0.75,1.1) .. (0.75,0.5);
      \draw[popgreen, thick] (-0.45,0.7) .. controls (-0.45,1.0) and (-0.2,1.05) .. (0,1.05)
        .. controls (0.2,1.05) and (0.45,1.0) .. (0.45,0.7);
    \end{scope}}
  % labels
  \draw[->, popdark] (6.6,4.4) -- (5.5,3.1);
  \node[popdark, align=center] at (7.6,4.6) {convex algal laminae};
  \draw[->, popdark] (1.2,3.4) -- (1.9,2.2);
  \node[popdark, align=center] at (1.0,3.7) {living microbial\\mat on top};
  \draw[->, popdark] (10.6,3.2) -- (9.9,1.5);
  \node[popdark, align=center] at (11.3,3.5) {trapped sediment\\between mats};
\end{tikzpicture}
\legende{Structure of stromatolites: successive sticky microbial mats trap sediment and grow upward, building laminated domes. The lamination itself --- not any preserved cell --- is the fossil evidence.}
\end{popfigure}

\begin{attention}[Stromatolites are NOT body fossils]
A frequent mistake is to classify stromatolites as body fossils of algae. In most stromatolites \emph{no cells of the micro-organisms are preserved}. What is fossilised is the \textbf{sedimentary structure built by their activity} --- the lamination. Stromatolites are therefore \cle{organo-sedimentary structures}, functionally comparable to trace fossils (evidence of life activity, not of the body). In an examination, write: ``stromatolites are biogenic sedimentary structures produced by microbial mats.''
\end{attention}

\courssub{Chemical Fossils (Biomarkers)}

\textbf{Intuitive idea.}\par
Even when an organism decays completely and leaves neither body nor trace, its \emph{molecules} may survive. Certain tough organic compounds --- or their stable alteration products --- persist in sedimentary rocks for hundreds of millions of years. Detecting them in a rock is like detecting a chemical fingerprint of past life.

\begin{definition}[Chemical fossil (biomarker)]
A \cle{chemical fossil} or \cle{biomarker} (biological marker) is an organic molecule preserved in a rock or sediment whose structure or composition demonstrates a biological origin. Biomarkers are derived from the lipids, pigments and other compounds of once-living organisms and survive even when all morphological evidence has been destroyed.
\end{definition}

Examples include hydrocarbons derived from chlorophyll and from membrane lipids, and the molecular signatures of specific bacterial groups. Biomarkers are used to:
\begin{itemize}
  \item demonstrate the \textbf{presence of life} in ancient or metamorphosed rocks where no fossils survive (including the search for early life);
  \item identify the \textbf{types of organisms} that contributed organic matter to a sediment (e.g.\ algae versus land plants);
  \item trace the \textbf{source rocks of petroleum}: the biomarkers in crude oil link it to the organic-rich mudrocks from which it was generated --- of direct economic importance for petroleum exploration in sedimentary basins such as the Rio del Rey and Douala basins of Cameroon.
\end{itemize}

\begin{savaistu}[Did you know?]
The oldest convincing evidence of life on Earth combines both kinds of ``invisible'' fossils studied here: stromatolite-like laminated structures \emph{and} organic molecules with biological signatures, found in rocks more than three billion years old. Long before any animal had a shell to fossilise, life was already writing its diary in chemistry and sediment.
\end{savaistu}

\courssub{Comparison of Body, Trace and Chemical Fossils}

\begin{center}
\begin{tabularx}{\linewidth}{|l|X|X|X|}
\hline
\rowcolor{popblueL} \textbf{Criterion} & \textbf{Body fossils} & \textbf{Trace fossils} & \textbf{Chemical fossils} \\
\hline
Nature of evidence & Remains of the organism itself (shell, bone, leaf) & Structures made by activity (tracks, burrows, coprolites) & Organic molecules (biomarkers) in the rock \\
\hline
Main information & Anatomy, taxonomy, evolution & Behaviour, locomotion, environment & Presence and type of life; organic sources \\
\hline
Transport possible? & Yes (may be displaced after death) & No --- always \emph{in situ} & Dispersed within the sediment \\
\hline
Dating value & Good (index fossils) & Poor (facies fossils, long-ranging) & Poor; indicates life, not precise age \\
\hline
Environmental value & Moderate & Excellent (depth, energy, oxygen) & Good (type of ecosystem, source rock) \\
\hline
Examples & Ammonite shell, trilobite, bone & \emph{Cruziana} trail, \emph{Skolithos} burrow, coprolite, stromatolite & Chlorophyll-derived hydrocarbons in oil source rocks \\
\hline
\end{tabularx}
\end{center}

\begin{exoresolu}[Identifying fossil types]
A student examining a slab of sandstone from a shallow-marine formation observes: (a) numerous vertical tubes filled with sand; (b) on one bedding surface, a double row of small imprints in a regular series; (c) within the rock, chemical analysis reveals hydrocarbons derived from chlorophyll. Classify each observation and state one piece of information each provides.
\tcblower
\textbf{Solution.}
\begin{itemize}
  \item (a) The vertical tubes are \textbf{trace fossils} --- dwelling burrows (of \emph{Skolithos} type). They indicate a high-energy, shallow-water, well-oxygenated sandy environment, and they are \emph{in situ} evidence (burrows cannot be transported).
  \item (b) The regular series of imprints is a \textbf{trackway} (locomotion trace fossil). The spacing and size of the imprints give information on the producer's mode of locomotion, size and possibly speed --- evidence of \emph{behaviour}, not anatomy.
  \item (c) The hydrocarbons derived from chlorophyll are \textbf{chemical fossils (biomarkers)}. They prove that photosynthetic organisms contributed organic matter to the sediment, even though no body or trace of them is preserved.
\end{itemize}
\emph{Key point:} none of the three observations contains any part of an organism's body, yet together they document the life, behaviour and environment of an ancient sea floor.
\end{exoresolu}

\begin{exercice}[Applying the concepts]
\begin{enumerate}
  \item Define a trace fossil and give three examples, stating the behaviour each records.
  \item Explain why trace fossils are described as \emph{facies fossils} and why this makes them poor index fossils but good environmental indicators.
  \item A mudstone is perfectly laminated and contains no burrows. What does this suggest about bottom-water oxygen when the mud was deposited? Justify your answer.
  \item Why are stromatolites classified as organo-sedimentary structures rather than body fossils?
  \item Give two practical uses of biomarkers in petroleum geology.
\end{enumerate}
\end{exercice}

\begin{corrige}[Exercise 1]
\begin{enumerate}
  \item A trace fossil is a sedimentary structure produced by the activity of a living organism and preserved in rock, recording behaviour rather than the body. Examples: footprints/trackways (locomotion), burrows (dwelling or feeding), coprolites (diet/feeding), borings (predation or attachment), trails (crawling).
  \item A given behaviour is adapted to a particular environment (substrate, energy, oxygen) and changes very slowly through geological time, so the same trace occurs in the same facies over enormous time ranges. Hence traces characterise \emph{environment} (facies) rather than \emph{age}: poor index fossils, excellent palaeoenvironmental indicators.
  \item Bottom waters were poorly oxygenated or anoxic: burrowing organisms could not live in the sediment, so lamination was never disturbed (bioturbated). Well-oxygenated muds are characteristically churned by burrowers.
  \item Because the micro-organisms' bodies are usually not preserved; only the laminated sedimentary structure built by the activity of microbial mats is fossilised. It is evidence of life \emph{activity}, like a trace fossil --- an organo-sedimentary (biogenic) structure.
  \item Biomarkers identify the source rocks from which petroleum was generated and reveal the type of organic matter (e.g.\ algal versus land-plant) that produced the oil, guiding exploration in basins such as those of coastal Cameroon.
\end{enumerate}
\end{corrige}

\begin{aretenir}[Summary]
\begin{itemize}
  \item \cle{Trace fossils} record \textbf{activity}: tracks and trails (locomotion), burrows (dwelling/feeding), borings, coprolites (diet), gastroliths, nests and eggs.
  \item They reveal behaviour, water depth, energy and oxygenation; they are \emph{in situ} and record soft-bodied life.
  \item Trace fossils are \cle{facies fossils}: environment indicators, not age indicators.
  \item \cle{Stromatolites} are laminated organo-sedimentary structures built by microbial mats --- not body fossils.
  \item \cle{Chemical fossils (biomarkers)} are preserved organic molecules proving past life and tracing petroleum source rocks.
\end{itemize}
\end{aretenir}

\coursec{Conditions for Fossilization}

In the previous sections we explored the rich variety of body fossils, trace fossils and chemical fossils — evidence that life leaves behind not only its bodily remains but also its activities and molecular signatures. Yet the vast majority of organisms that have ever lived left \emph{no} fossil record at all. Fossilization is not the rule; it is a remarkable exception. Understanding \emph{why} some organisms become fossils while others vanish without a trace is the key to interpreting the fossil record correctly. This section examines the stringent conditions that must be met for an organism to enter the geological archive.

\courssub{Fossilization: A Rare and Exceptional Event}

When an organism dies, its organic matter is immediately subjected to a host of destructive processes. Scavengers and predators dismember carcasses; bacteria and fungi decompose soft tissues; currents, waves and wind scatter skeletal elements; and chemical weathering dissolves or oxidizes mineralized hard parts. In most environments, a dead organism is completely recycled back into the ecosystem within days, weeks, or at most a few years.

\begin{definition}[Fossilization]
Fossilization is the set of physical, chemical and biological processes that transform the remains, traces or products of once-living organisms into fossils preserved in the geological record. It involves the transition of organic material from the biosphere into the lithosphere.
\end{definition}

\begin{aretenir}[The Odds Against Fossilization]
\begin{itemize}
    \item Only a very small fraction of all the species that have ever lived are known from fossils; the fossil record samples past life extremely incompletely.
    \item Soft-bodied organisms (worms, jellyfish, most plant tissues) fossilize only under extraordinary circumstances.
    \item Even organisms with hard parts (shells, bones, wood) are usually destroyed before burial.
    \item The fossil record is therefore a highly biased sample of past life, filtered by the conditions of preservation.
\end{itemize}
\end{aretenir}

\courssub{The Five Essential Conditions for Fossilization}

For fossilization to occur, a combination of favourable circumstances must coincide. We can summarize them as a checklist of five conditions. The more conditions are met, the higher the probability of preservation.

\begin{methode}[Checklist: Five Conditions for Fossilization]
\begin{enumerate}
    \item \textbf{Possession of hard parts:} The organism should have mineralized tissues (shells, bones, teeth, woody tissue) or resistant organic compounds (chitin, sporopollenin, lignin). Soft tissues decay too rapidly in ordinary settings.
    \item \textbf{Rapid burial:} The remains must be covered by sediment quickly — before scavengers, currents or decay destroy them. This requires high sedimentation rates or sudden events (storms, floods, volcanic ash falls).
    \item \textbf{Anaerobic (oxygen-poor) conditions:} Burial in an oxygen-deficient environment (anoxic bottom waters, deep mud, tar, amber, ice) drastically slows bacterial decay and oxidation.
    \item \textbf{Absence of disturbance:} The depositional setting must be relatively quiet (low energy) to prevent physical breakage, abrasion and winnowing of remains. Bioturbation (burrowing by organisms) must be minimal.
    \item \textbf{Favourable chemistry for early mineralization:} Pore waters must be saturated with respect to minerals (calcite, silica, pyrite, phosphate) that can precipitate in pore spaces or replace organic material before compaction and dissolution occur. Later diagenesis and metamorphism must not destroy the fossil.
\end{enumerate}
\end{methode}

Let us examine each condition in detail.

\courssub{Condition 1: Possession of Hard Parts}

Organisms with mineralized skeletons (calcium carbonate, calcium phosphate, silica) or resistant organic walls have an enormous advantage.

\begin{itemize}
    \item \textbf{Mineralized hard parts:} shells of molluscs, brachiopods, corals and foraminifera; bones and teeth of vertebrates; tests of echinoderms; carapaces of arthropods (partially mineralized); frustules of diatoms and skeletons of radiolarians (silica).
    \item \textbf{Resistant organic parts:} woody tissue (lignin) in plants; spores and pollen (sporopollenin), among the most resistant organic materials known; chitin in arthropod cuticles; the periderm of graptolites.
\end{itemize}

\begin{propriete}[Hard Parts and Fossilization Probability]
Organisms possessing mineralized hard parts have a fossilization potential far higher than that of soft-bodied organisms. This is the single most important intrinsic factor. However, hard parts alone are insufficient without rapid burial.
\end{propriete}

\begin{attention}[Common Mistake]
Do not assume that because an organism has hard parts it \emph{will} fossilize. A clam shell lying on a beach will be broken by waves, bored by organisms and dissolved by slightly acidic rain within years. Hard parts are a \emph{necessary but not sufficient} condition.
\end{attention}

\courssub{Condition 2: Rapid Burial}

Rapid burial is the critical extrinsic factor. It physically shields the remains from:
\begin{itemize}
    \item scavengers and mechanical disarticulation;
    \item oxygenated water that fuels aerobic decay;
    \item currents and waves that transport and abrade fragments.
\end{itemize}

Mechanisms of rapid burial include:
\begin{itemize}
    \item \textbf{Storm deposits (tempestites):} high-energy events that dump thick layers of sediment in hours (e.g., hurricane deposits on the continental shelf).
    \item \textbf{Flood events:} river floods carrying massive sediment loads into deltas, lakes and coastal seas.
    \item \textbf{Volcanic ash falls:} near-instantaneous blanketing of landscapes (e.g., the Laetoli footprints in Tanzania, preserved beneath volcanic ash).
    \item \textbf{Turbidity currents:} underwater avalanches that transport sediment and organisms into deep basins.
    \item \textbf{Mass wasting:} landslides and mudflows.
\end{itemize}

\begin{exemplebox}[A Cameroonian Analogue for Catastrophic Burial: Lake Nyos]
The 1986 limnic eruption at Lake Nyos (Northwest Region) released a dense cloud of carbon dioxide that suffocated people and animals in the surrounding valleys. Carcasses were subsequently covered by landslide debris and by sediment washed into the valleys by heavy rains. Although this tragedy is far too recent to have produced fossils, it is a useful \emph{modern analogue} showing how a catastrophic event can kill and bury an entire community at once — the first step towards a fossilized ``death assemblage''.
\end{exemplebox}

\courssub{Condition 3: Anaerobic (Oxygen-Poor) Conditions}

Aerobic bacteria require oxygen to decompose organic matter efficiently. In anoxic environments, decay is carried out only by much slower anaerobic bacteria (sulfate reducers, methanogens). This gives time for mineralization to begin.

Environments that are commonly anoxic:
\begin{itemize}
    \item deep basins with restricted water circulation (e.g., the Black Sea, some fjords);
    \item permanently stratified (meromictic) lakes, where bottom waters never mix with surface waters;
    \item lagoons and estuaries with high organic input and restricted exchange with the open sea;
    \item sediment pore waters just below the sediment--water interface (below the ``redox boundary'');
    \item tar pits (asphalt seeps), amber (fossilized tree resin) and ice (permafrost, glaciers) — extreme media that exclude oxygen and preserve soft tissues exceptionally well.
\end{itemize}

\begin{propriete}[Anoxia and Soft-Tissue Preservation]
Anoxic conditions are a prerequisite for the preservation of soft tissues (muscles, guts, skin, feathers, leaves) and for the formation of \emph{Konservat-Lagerstätten} — deposits of exceptional preservation such as the Burgess Shale (Canada), the Solnhofen Limestone (Germany) and the Messel Pit (Germany).
\end{propriete}

\courssub{Condition 4: Absence of Disturbance (Quiet Depositional Environment)}

Even if buried, fossils can be destroyed by:
\begin{itemize}
    \item \textbf{Bioturbation:} burrowing organisms (worms, clams, crustaceans) churn the sediment, fragmenting shells, destroying delicate structures and mixing fossil assemblages of different ages (time averaging). Intensely bioturbated sediments, such as well-oxygenated shallow shelves, show poor preservation of articulated fossils.
    \item \textbf{Physical reworking:} waves, currents and tides winnow fine sediment, concentrate shells into lag deposits and abrade surfaces.
    \item \textbf{Compaction and tectonism:} later burial under kilometres of sediment flattens fossils; mountain building fractures and may metamorphose them.
\end{itemize}

Quiet environments favouring preservation:
\begin{itemize}
    \item deep marine settings (below storm wave base);
    \item lake bottoms, especially deep anoxic lakes;
    \item low-energy lagoons and tidal flats (if not heavily bioturbated);
    \item floodplains and swamps (overbank fines);
    \item caves and fissures.
\end{itemize}

\courssub{Condition 5: Favourable Chemistry for Early Mineralization}

The final condition concerns the geochemistry of the pore waters during early diagenesis. For a fossil to survive the immense spans of geological time, its original material must be stabilized.

\begin{itemize}
    \item \textbf{Permineralization (petrifaction):} mineral-rich groundwater precipitates minerals (calcite, silica, pyrite, hematite) in the pore spaces of bone, wood or shell, preserving internal microstructure.
    \item \textbf{Replacement:} the original material dissolves and is simultaneously replaced by a new mineral (e.g., silica replacing a calcite shell, or pyrite replacing organic matter), often replicating fine detail.
    \item \textbf{Recrystallization:} unstable mineral phases (aragonite, high-Mg calcite) invert to stable low-Mg calcite. This may preserve the overall shape but destroys microstructure.
    \item \textbf{Dissolution:} if pore waters are undersaturated (e.g., acidic due to organic decay or meteoric water), shells dissolve completely, leaving only moulds.
\end{itemize}

\begin{propriete}[The Early Diagenetic Window]
The ``window of opportunity'' for fossilization is early diagenesis — the first metres to tens of metres of burial. Once the sediment is deeply buried and lithified, the fossil is relatively safe, unless later metamorphism, tectonic deformation or weathering intervenes.
\end{propriete}

\courssub{Environments Favourable to Fossilization}

Combining the five conditions, we can rank depositional environments by their fossilization potential.

\begin{center}
\begin{tabularx}{\linewidth}{|l|X|X|}\hline
\rowcolor{popblueL}\textbf{Environment} & \textbf{Favourable Factors} & \textbf{Typical Fossils} \\ \hline
\textbf{Marine shelf (below fair-weather wave base)} & Rapid burial by storms; anoxic bottom waters possible; hard parts abundant & Molluscs, brachiopods, echinoderms, bryozoans, foraminifera \\ \hline
\textbf{Deep marine (turbidites, pelagic ooze)} & Very quiet; anoxia possible; rapid burial by turbidity currents & Microfossils (foraminifera, radiolarians, diatoms), fish teeth, rare articulated vertebrates \\ \hline
\textbf{Lake bottoms (deep, stratified)} & Anoxic; quiet; rapid burial by fine laminae & Fish, insects, plants, ostracods, diatoms (e.g., crater lakes of the Cameroon Volcanic Line such as Lake Barombi Mbo) \\ \hline
\textbf{Deltas, floodplains, swamps} & Very rapid burial; anoxic waterlogged soils; high sedimentation & Plants (coal), vertebrates, freshwater molluscs, footprints \\ \hline
\textbf{Lagoons, restricted bays} & Anoxic; quiet; high salinity excludes bioturbators & Soft-bodied organisms, articulated arthropods, fish (Solnhofen-type) \\ \hline
\textbf{Tar pits (asphalt seeps)} & Anaerobic; sticky trap; rapid entombment & Vertebrates, insects, plants (e.g., La Brea Tar Pits, USA) \\ \hline
\textbf{Amber (fossil resin)} & Anaerobic; chemical preservation; rapid entombment & Insects, spiders, feathers, flowers, microorganisms \\ \hline
\textbf{Ice / permafrost} & Freezing arrests decay; anaerobic & Mammoths, woolly rhinoceroses, ancient DNA \\ \hline
\end{tabularx}
\end{center}

\courssub{Why Terrestrial Environments Fossilize Poorly}

The fossil record is overwhelmingly marine. Terrestrial environments are notoriously poor at preserving fossils for several reasons:

\begin{enumerate}
    \item \textbf{Erosion dominates:} continents are net erosional surfaces. Sediment is transported \emph{to} the oceans rather than accumulated on land. Most terrestrial sediments are temporary (alluvial fans, river channels) and are later re-eroded.
    \item \textbf{Oxidation:} well-drained soils are oxygen-rich. Organic matter decays completely, and iron minerals oxidize to red or brown hematite and goethite (red beds), destroying organic fossils.
    \item \textbf{Acidic soils:} humic acids in forest soils dissolve calcium carbonate shells and bones rapidly. Calcareous fossils are rare in non-calcareous terrestrial rocks.
    \item \textbf{Bioturbation:} soil organisms (earthworms, termites, roots) thoroughly mix the upper metres of sediment, destroying primary structures and fragmenting remains.
    \item \textbf{Scavenging and disarticulation:} terrestrial vertebrate carcasses are efficiently dismembered by scavengers (hyenas, vultures, insects) and scattered over large areas.
\end{enumerate}

\begin{aretenir}[Terrestrial Fossilization Windows]
Terrestrial fossils are mostly found in:
\begin{itemize}
    \item \textbf{lacustrine deposits} (lake muds, especially of deep anoxic lakes);
    \item \textbf{fluvial overbank fines} (floodplain mudstones, palaeosols with calcrete);
    \item \textbf{cave and fissure fills} (sediment washed into caves, e.g., hominid sites in East and South Africa);
    \item \textbf{volcanic ash falls} (instant burial, e.g., Laetoli in Tanzania);
    \item \textbf{tar pits, amber and permafrost} (exceptional preservation).
\end{itemize}
\end{aretenir}

\courssub{Illustrations}

\begin{popfigure}
\centering
\begin{tikzpicture}[
    node distance=1.2cm and 1.6cm,
    box/.style={draw=popdark, thick, rounded corners, fill=popblueL, text width=2.8cm, align=center, minimum height=1.2cm},
    arrow/.style={->, thick, popdark},
    start/.style={box, fill=popgreenL},
    end/.style={box, fill=poporange},
]
\node[start] (death) {Death of organism};
\node[box, right=of death] (decay) {Decay, scavenging, disarticulation, dissolution};
\node[box, below=of decay] (burial) {Rapid burial by sediment};
\node[box, left=of burial] (accum) {Sediment accumulation and compaction};
\node[box, left=of accum] (diagen) {Early diagenesis: permineralization, replacement};
\node[end, left=of diagen] (fossil) {Fossil preserved in sedimentary rock};

\draw[arrow] (death) -- node[above, font=\small, align=center] {Most organisms: destruction} (decay);
\draw[arrow, popgreen] (death.south) |- node[above, font=\small, pos=0.75] {Rare: rapid burial} (burial.west);
\draw[arrow] (burial) -- (accum);
\draw[arrow] (accum) -- (diagen);
\draw[arrow] (diagen) -- (fossil);

\node[box, fill=poppink, right=of burial, text width=3cm] (destroyed) {Complete destruction: no fossil record};
\draw[arrow, poppink] (decay) -- (destroyed);
\end{tikzpicture}
\legende{Flow diagram showing the pathway from death to fossilization. The vast majority of organisms follow the path to complete destruction. Only rapid burial diverts remains into the fossilization pathway.}
\end{popfigure}

\begin{popfigure}
\centering
\begin{tikzpicture}[scale=0.9]
% Sea level
\draw[popblue, very thick] (-5,2) -- (5,2) node[right] {Sea level};
% Water column
\fill[popblueL, opacity=0.3] (-5,2) rectangle (5,-4);
% Sediment layers
\fill[popmuted!50] (-5,-4) rectangle (5,-5);
\fill[popmuted!30] (-5,-5) rectangle (5,-6);
% Labels for environments
\node[align=center, text width=3.2cm] at (-3.5,1) {\textbf{High-energy shoreface}\\\emph{Unfavourable:} waves, currents, oxidation, bioturbation};
\node[align=center, text width=3.2cm] at (3.5,1) {\textbf{Offshore shelf}\\\emph{Favourable:} below wave base, storm burial, anoxia possible};
\node[align=center, text width=3.2cm] at (-3.5,-3) {\textbf{Shoreface sands}\\Shell hash: fragmented, abraded, time-averaged};
\node[align=center, text width=3.2cm] at (3.5,-3) {\textbf{Offshore muds}\\Articulated shells, soft tissues possible, low bioturbation};
% Arrows for sedimentation
\draw[->, thick, poporange] (-3.5,0.3) -- (-3.5,-2.3) node[midway, left, font=\small, align=center] {Storm\\layers};
\draw[->, thick, poporange] (3.5,0.3) -- (3.5,-2.3) node[midway, right, font=\small, align=center] {Settling +\\storms};
% Living organisms (surface)
\foreach \x/\y in {-4/1.5, -3/1.2, -2/1.8} {
    \draw[popdark] (\x,\y) ellipse (0.15 and 0.1);
}
\foreach \x/\y in {2/1.5, 3/1.2, 4/1.8, 2.5/0.5, 3.5/0.2} {
    \draw[popgreen] (\x,\y) ellipse (0.15 and 0.1);
}
% Buried fossils
\foreach \x/\y in {-4/-4.3, -3/-4.6, -2/-4.4} {
    \draw[popdark, dashed] (\x,\y) ellipse (0.15 and 0.1);
}
\foreach \x/\y in {2/-4.3, 3/-4.6, 4/-4.4, 2.5/-4.8, 3.5/-5.1} {
    \draw[popgreen] (\x,\y) ellipse (0.15 and 0.1);
}
% Legend
\node[anchor=south west, fill=white, draw=popdark, rounded corners, align=left, font=\small] at (-4.9,-5.9) {\textcolor{popdark}{\rule{0.5cm}{0.15cm}} Fragmented, abraded shells\\ \textcolor{popgreen}{\rule{0.5cm}{0.15cm}} Well-preserved, articulated fossils};
\end{tikzpicture}
\legende{Cross-section of a marine shelf showing contrasting fossilization potential. The high-energy shoreface (left) destroys and fragments remains. The quieter offshore shelf (right), below fair-weather wave base, receives storm deposits and can develop anoxic bottom waters, favouring exceptional preservation.}
\end{popfigure}

\courssub{Worked Example: Assessing Fossilization Potential}

\begin{exoresolu}[Comparing Two Depositional Settings]
\textbf{Statement:} Two fossil assemblages are found in the same geological formation:
\begin{enumerate}
    \item Assemblage A: disarticulated, abraded bivalve shells concentrated in a coarse sandstone with cross-bedding. The shells are a mixture of species that lived in different habitats.
    \item Assemblage B: articulated bivalves (both valves closed), delicate gastropods with intact apertures, and rare crustacean cuticles preserved in a fine-grained, laminated shale containing pyrite framboids.
\end{enumerate}
Analyse the fossilization conditions for each assemblage and deduce the depositional environment.
\tcblower
\textbf{Detailed solution:}

\textbf{Assemblage A (sandstone):}
\begin{itemize}
    \item \textbf{Hard parts:} present (bivalve shells) — Condition 1 met.
    \item \textbf{Rapid burial:} cross-bedding indicates migrating ripples or dunes — burial was episodic but not rapid enough to prevent reworking. The shells are abraded and disarticulated, implying prolonged exposure on the seafloor.
    \item \textbf{Anaerobic conditions:} coarse sandstone implies well-oxygenated pore waters — Condition 3 \emph{not} met.
    \item \textbf{Absence of disturbance:} cross-bedding and abrasion show strong current activity — Condition 4 \emph{not} met. The mixing of habitats indicates transport and time averaging.
    \item \textbf{Chemistry:} oxidizing conditions — original aragonite was probably dissolved or recrystallized.
    \item \textbf{Conclusion:} low fossilization potential. This is a \emph{lag deposit} (a winnowed concentration) in a high-energy shoreface or tidal channel. The assemblage is \emph{allochthonous} (transported) and time-averaged.
\end{itemize}

\textbf{Assemblage B (laminated shale):}
\begin{itemize}
    \item \textbf{Hard parts:} present (shells, cuticles) — Condition 1 met.
    \item \textbf{Rapid burial:} laminations indicate quiet settling from suspension; articulated bivalves with closed valves imply burial before the adductor muscles decayed (days to weeks) — Condition 2 met.
    \item \textbf{Anaerobic conditions:} pyrite framboids (\ce{FeS2}) are a classic indicator of anoxic, sulfidic pore waters produced by sulfate-reducing bacteria — Condition 3 met. This also explains the preservation of crustacean cuticle (chitin).
    \item \textbf{Absence of disturbance:} perfect laminations mean no bioturbation and no currents — Condition 4 met.
    \item \textbf{Chemistry:} pyritization of soft tissues and probable early carbonate cement — Condition 5 met.
    \item \textbf{Conclusion:} high fossilization potential. This is an \emph{autochthonous} (in situ) assemblage in a low-energy, anoxic offshore shelf or restricted basin — preservation of \emph{Konservat-Lagerstätte} type.
\end{itemize}

\textbf{Key lesson:} the sedimentology (sandstone versus shale, cross-bedding versus lamination, presence of pyrite) directly reveals which of the five conditions were satisfied.
\end{exoresolu}

\courssub{Summary and Key Points}

\begin{aretenir}[Essential Summary: Conditions for Fossilization]
\begin{enumerate}
    \item Fossilization is \textbf{exceptional}; the vast majority of organisms decay completely.
    \item \textbf{Five conditions} must align: (1) hard parts, (2) rapid burial, (3) anaerobic conditions, (4) quiet environment (no disturbance), (5) favourable early diagenetic chemistry.
    \item \textbf{Marine shelves and deep basins} are the most favourable common environments; \textbf{terrestrial settings} are poor because of erosion, oxidation and acidic soils.
    \item \textbf{Lagerstätten} (deposits of exceptional preservation) form when \emph{all} conditions are met, especially anoxia combined with rapid burial (e.g., volcanic ash falls, anoxic lagoons).
    \item The \textbf{sedimentary context} (rock type, sedimentary structures, geochemistry) is the primary evidence for reconstructing fossilization conditions.
\end{enumerate}
\end{aretenir}

\begin{attention}[Examination Traps]
\begin{itemize}
    \item \textbf{Trap:} ``Fossilization requires hard parts.'' \textbf{Correction:} hard parts increase the probability, but soft-bodied organisms fossilize in anoxic, rapid-burial settings (e.g., the Burgess Shale). It is the \emph{combination} of conditions that matters.
    \item \textbf{Trap:} ``Rapid burial alone causes fossilization.'' \textbf{Correction:} rapid burial in oxygenated, bioturbated sand still leads to destruction. Anaerobiosis is equally critical.
    \item \textbf{Trap:} confusing \emph{time averaging} (mixing of generations over centuries or millennia in a single bed) with a single catastrophic burial event. Use articulation, abrasion and sedimentary structures to distinguish them.
    \item \textbf{Trap:} assuming terrestrial fossils are as common as marine fossils. Always justify terrestrial fossil occurrences with specific environments (lakes, caves, ash falls).
\end{itemize}
\end{attention}

\begin{exercice}[Self-Assessment]
\begin{enumerate}
    \item A dinosaur skeleton is found in a sandstone channel fill with cross-bedding. The bones are disarticulated, show tooth marks and are mixed with rounded pebbles. Which of the five conditions were \emph{not} met? What does this tell you about the depositional history?
    \item Explain why pollen and spores (sporopollenin) are abundant in the fossil record even in oxidized terrestrial red beds where bones are absent.
    \item A shale unit contains abundant pyritized ammonites and carbonized plant leaves. The shale is laminated with no bioturbation. List the five conditions and state the evidence from the description that each was satisfied.
    \item Why are the volcanic ash layers of the Cameroon Volcanic Line (e.g., around Mount Cameroon) potential sites for exceptional fossil preservation? Which conditions do they satisfy uniquely?
\end{enumerate}
\end{exercice}

\begin{corrige}[Exercise 1]
\begin{enumerate}
    \item \textbf{Conditions not met:} (2) rapid burial — the bones were exposed long enough for scavenging and disarticulation; (3) anaerobic conditions — channel sandstones are well oxygenated; (4) absence of disturbance — cross-bedding and pebbles show high-energy currents. \textbf{History:} the carcass decomposed on a floodplain; the bones were then transported and concentrated in a river-channel lag deposit (time-averaged, allochthonous).
    \item Sporopollenin is chemically extremely resistant to oxidation and acid hydrolysis. It survives in oxidizing, acidic soils where calcium phosphate (bone) and calcium carbonate (shell) dissolve. Condition 1 (resistant organic wall) is met even when Conditions 3 and 5 fail for mineralized tissues.
    \item (1) Hard parts: ammonite shells (aragonite) and plant leaves (carbon). (2) Rapid burial: articulated ammonites and delicate leaves preserved flat. (3) Anaerobic conditions: pyrite (\ce{FeS2}) requires sulfate reduction under anoxia. (4) Absence of disturbance: laminations and no bioturbation. (5) Favourable chemistry: pyritization of ammonites and carbonization of leaves (early diagenetic mineralization).
    \item Volcanic ash falls provide \textbf{instantaneous burial} (Condition 2 par excellence), sealing organisms in place. Fresh ash is fine-grained and initially sterile, creating anoxic microenvironments (Condition 3) and preventing bioturbation (Condition 4). As glass shards alter to clay, they provide reactive surfaces for early mineralization (Condition 5). This is why deposits such as Laetoli (Tanzania) preserve footprints and delicate remains in ash.
\end{enumerate}
\end{corrige}

\coursec{Modes of Fossilization I: Unaltered Remains and Replacement Processes}

In the previous section we established that fossilization is an exceptional event: an organism must possess resistant hard parts, be buried rapidly by sediment, and escape scavengers, bacteria and chemical destruction. But burial is only the beginning of the story. Once enclosed in sediment, the remains of the organism enter a long chemical dialogue with the surrounding rock and the water circulating through it. Depending on the nature of that dialogue, the fossil that a palaeontologist finally extracts may be the \emph{original} tissue of the organism, astonishingly intact after thousands of years, or it may be a mineral replica in which not a single atom of the original creature remains. The different pathways by which buried remains are transformed into fossils are called the \cle{modes of fossilization}. In this section we study the modes in which the remains are preserved unaltered, and the three great mineralogical processes --- permineralization, replacement and recrystallization --- that convert organic hard parts into stone while conserving their form. Moulds, casts, carbonization and imprints will be treated in the following section.

\courssub{Overview of the Modes of Fossilization}

The modes of fossilization are traditionally grouped into two families: those that preserve the organism \textbf{unaltered} (or nearly so), and those in which the remains are \textbf{altered} by mineralogical or chemical change. The table below summarises the whole family; the shaded rows are those studied in this section.

\begin{center}
\begin{tabularx}{\linewidth}{|l|X|X|}
\hline
\rowcolor{popblueL} \textbf{Mode} & \textbf{Process} & \textbf{Typical example} \\
\hline
\rowcolor{popblueL} Freezing & Whole body preserved in permafrost & Woolly mammoths of Siberia \\
\hline
\rowcolor{popblueL} Desiccation & Drying (mummification) in arid caves or deserts & Mummified ground sloths \\
\hline
\rowcolor{popblueL} Amber entrapment & Insects trapped in hardened plant resin & Insects in Baltic amber \\
\hline
\rowcolor{popblueL} Tar pits and peat bogs & Asphalt or acidic bog water preserves bones and bodies & Bones of La Brea tar pits; bog bodies \\
\hline
\rowcolor{popblueL} Permineralization & Minerals deposited in pores of bone or wood & Petrified wood, dinosaur bone \\
\hline
\rowcolor{popblueL} Replacement & Original material dissolved and substituted by mineral & Silicified shells, pyritized ammonites \\
\hline
\rowcolor{popblueL} Recrystallization & Crystal structure changes, composition unchanged & Aragonite shells converted to calcite \\
\hline
Carbonization & Volatiles lost, leaving a carbon film & Leaves, graptolites, soft-bodied animals \\
\hline
Moulds and casts & Impression of the organism, later filled by sediment & Internal moulds of bivalves \\
\hline
Imprints and tracks & Traces of activity or surface impressions & Footprints, burrows, leaf imprints \\
\hline
\end{tabularx}
\end{center}

\courssub{Unaltered (Original) Preservation}

Imagine finding an animal that died tens of thousands of years ago and yet still has its skin, its hair and even the remains of its last meal in its stomach. This is not science fiction: under exceptional circumstances, the normal agents of decay --- oxygen, bacteria, fungi and scavengers --- are completely excluded, and the soft tissues survive. Such fossils are said to be preserved \cle{unaltered} or in their \cle{original} state. They are extremely rare, but they are scientifically precious because they reveal the anatomy, colour and diet of extinct animals in a way that mineralized fossils never can.

\textbf{Freezing.} In the permanently frozen ground (\cle{permafrost}) of Siberia and Alaska, the carcasses of woolly mammoths and woolly rhinoceroses that died during the last Ice Age have been kept in a natural deep-freeze. At temperatures permanently below $0\ ^{\circ}\mathrm{C}$, bacterial activity is virtually stopped, and flesh, hair and hide survive. Some mammoth carcasses are so fresh that their flesh has been eaten by dogs, and their stomach contents reveal the grasses and sedges of the vanished tundra-steppe.

\textbf{Desiccation (mummification).} In very dry caves and deserts, a carcass may dry out faster than bacteria can decompose it. The result is a natural mummy: skin, tendons and dried muscles survive. Mummified remains of ground sloths, complete with skin and dung, have been recovered from dry caves in North and South America.

\textbf{Entrapment in amber.} \cle{Amber} is the fossilized resin of ancient trees. Insects, spiders and small lizards that became stuck in the sticky resin were sealed away from air and bacteria, and the resin hardened into a transparent golden solid. Insects in amber are often preserved with their finest details --- wing veins, hairs, even internal organs --- intact. Amber is therefore a window on small, delicate organisms that would otherwise leave no fossil record at all.

\textbf{Tar pits and peat bogs.} At natural asphalt seeps such as the famous La Brea tar pits of California, animals that wandered onto the sticky surface became trapped; their bones sank into the asphalt, which excluded oxygen and bacteria. Thousands of skeletons of sabre-toothed cats, dire wolves and mammoths have been recovered. In the acidic, oxygen-poor waters of peat bogs in northern Europe, even human bodies have been preserved with skin and hair intact, the acids tanning the skin like leather.

\begin{attention}[Unaltered does not mean eternal]
Unaltered preservation is only possible where decay is permanently blocked: continuous freezing, permanent dryness, or an airtight chemical seal. Such conditions are rare and geologically short-lived --- a thaw, a flood or a change of climate destroys the fossil. Consequently, unaltered remains are almost always geologically \emph{young} (Quaternary), and they can never tell us about the deep history of life.
\end{attention}

\courssub{Permineralization (Petrification)}

Pick up a piece of petrified wood: it looks like wood, it shows the growth rings and even the vessels of the original tree, yet it is heavy and rings like stone when struck. How can wood become stone without losing its architecture? The answer is \cle{permineralization}.

Bone and wood are \emph{porous} materials. Bone contains microscopic cavities that housed blood vessels and bone cells; wood is a lattice of hollow cells and vessels. When such material is buried in sediment, groundwater rich in dissolved minerals --- commonly silica ($\mathrm{SiO_2}$), calcite ($\mathrm{CaCO_3}$) or pyrite ($\mathrm{FeS_2}$) --- percolates through it. As the water moves through the pores, the dissolved minerals precipitate and crystallise \emph{inside the empty spaces}, gradually filling every cavity with solid mineral. The original organic framework may remain, at least in part, so that the internal anatomy is conserved in exquisite detail.

\begin{definition}[Permineralization]
\cle{Permineralization} (also called \cle{petrification}) is the mode of fossilization in which mineral-rich groundwater infiltrates the pores and cavities of a porous hard part (bone, wood, shell) and deposits minerals --- typically silica, calcite or pyrite --- within those empty spaces, without removing the original organic material. The fossil is therefore a composite: original tissue plus infilling mineral.
\end{definition}

\begin{popfigure}
\begin{tikzpicture}[font=\small]
% Step 1: porous bone
\draw[thick, popink, rounded corners=2pt] (0,0) rectangle (2.6,1.2);
\foreach \x/\y in {0.4/0.3, 0.9/0.75, 1.4/0.35, 1.9/0.8, 2.2/0.4, 0.6/0.85, 1.7/0.55}
  \draw[popink, fill=white] (\x,\y) circle (0.09);
\node[align=center, text width=2.6cm] at (1.3,-0.55) {1. Porous bone buried in sediment};
% arrow 1
\draw[->, very thick, poporange] (2.8,0.6) -- (3.6,0.6);
% Step 2: infiltration
\draw[thick, popink, rounded corners=2pt] (3.8,0) rectangle (6.4,1.2);
\foreach \x/\y in {4.2/0.3, 4.7/0.75, 5.2/0.35, 5.7/0.8, 6.0/0.4, 4.4/0.85, 5.5/0.55}
  \draw[popink, fill=white] (\x,\y) circle (0.09);
\draw[->, thick, popblue] (4.2,1.7) -- (4.2,1.25);
\draw[->, thick, popblue] (5.1,1.7) -- (5.1,1.25);
\draw[->, thick, popblue] (6.0,1.7) -- (6.0,1.25);
\node[popblue, align=center, text width=2.8cm] at (5.1,2.0) {mineral-rich groundwater};
\node[align=center, text width=2.6cm] at (5.1,-0.55) {2. Water infiltrates the pores};
% arrow 2
\draw[->, very thick, poporange] (6.6,0.6) -- (7.4,0.6);
% Step 3: mineral deposition
\draw[thick, popink, rounded corners=2pt] (7.6,0) rectangle (10.2,1.2);
\foreach \x/\y in {8.0/0.3, 8.5/0.75, 9.0/0.35, 9.5/0.8, 9.8/0.4, 8.2/0.85, 9.3/0.55}
  \draw[popink, fill=popgold] (\x,\y) circle (0.09);
\node[align=center, text width=2.6cm] at (8.9,-0.55) {3. Minerals crystallise in the pores};
\end{tikzpicture}
\legende{The three stages of permineralization of a bone: the empty pores (white circles) are progressively filled with mineral (gold circles) precipitated from circulating groundwater, while the original bony framework is retained.}
\end{popfigure}

\begin{propriete}[Petrified wood preserves annual rings]
In petrified wood, silica fills the hollow cell cavities and may also impregnate the cell walls. Because the mineral faithfully occupies the microscopic architecture of the wood, the cellular structure --- vessels, fibres and the \cle{annual growth rings} --- is preserved and can be studied in thin section under the microscope. Petrified wood can therefore be identified to the level of the original tree family, and its rings record the climate of the forest in which it grew. (The explanation of this property is examinable as an illustration of permineralization.)
\end{propriete}

\begin{exemplebox}[Dinosaur bone]
Most dinosaur bones are permineralized. The hard mineral matter of the living bone (calcium phosphate) is still present, but the marrow cavity and the microscopic canals that once carried blood vessels are now filled with calcite or silica. Cut and polished, such a bone still shows the spongy internal texture of living bone --- proof that the pores were filled, not the bone dissolved away.
\end{exemplebox}

\courssub{Replacement}

Permineralization \emph{adds} mineral to a fossil. \cle{Replacement} goes further: it \emph{substitutes} the original substance entirely.

\begin{definition}[Replacement]
\cle{Replacement} is the mode of fossilization in which the original hard material of the organism (shell, bone or wood) is dissolved by percolating water and simultaneously substituted, molecule by molecule, by a new mineral such as silica ($\mathrm{SiO_2}$), calcite ($\mathrm{CaCO_3}$), pyrite ($\mathrm{FeS_2}$) or iron oxides. Because the substitution proceeds gradually at a microscopic front, the external form --- and sometimes the fine internal structure --- of the original is faithfully reproduced in the new mineral.
\end{definition}

The key idea is that dissolution and precipitation occur \emph{at the same time and at the same place}. If the shell simply dissolved away, only a hole (a mould) would remain. Instead, each tiny volume of original material is replaced by mineral as it dissolves, so the shape is never lost. Well-known examples include:

\begin{itemize}
\item \textbf{Silicified shells and wood:} original calcite shells or woody tissue replaced by silica; fine ornamentation is preserved.
\item \textbf{Pyritized fossils:} ammonites and other shells replaced by iron pyrite (``fool's gold''), giving them a brilliant metallic lustre.
\item \textbf{Opalized fossils:} in parts of Australia, shells and even bones replaced by precious opal.
\end{itemize}

\courssub{Recrystallization}

Some fossils are chemically unchanged but structurally transformed. Many mollusc shells are originally built of \cle{aragonite}, an unstable crystalline form of calcium carbonate. Over geological time, aragonite spontaneously reorganises into \cle{calcite}, the stable form of the same compound, $\mathrm{CaCO_3}$. The composition is identical; only the arrangement of the atoms in the crystal lattice has changed.

\begin{definition}[Recrystallization]
\cle{Recrystallization} is the mode of fossilization in which the mineral of a hard part changes its crystal structure or crystal size \emph{without any change of chemical composition}. The classic example is the conversion of aragonite to calcite in shells: $\mathrm{CaCO_3}$ (aragonite) $\longrightarrow$ $\mathrm{CaCO_3}$ (calcite).
\end{definition}

Because the new calcite crystals are usually larger than the original aragonite crystallites, recrystallization often \emph{blurs or destroys the finest internal microstructure} of the shell, even though the external shape remains perfect. This is an important diagnostic point: a shell that looks well preserved but whose microscopic wall structure is coarse and mosaic-like has been recrystallized.

\begin{popfigure}
\begin{tikzpicture}[font=\small]
% Unaltered shell
\begin{scope}
\draw[thick, popink, fill=popblueL] (0,0) arc[start angle=180, end angle=0, radius=0.9] -- cycle;
\draw[popink] (0.15,0.15) arc[start angle=180, end angle=0, radius=0.75];
\draw[popink] (0.35,0.3) arc[start angle=180, end angle=0, radius=0.55];
\node[align=center, text width=2.6cm] at (0.9,-0.7) {Unaltered shell\\ original aragonite, fine structure};
\end{scope}
% Recrystallized shell
\begin{scope}[xshift=4.2cm]
\draw[thick, popink, fill=popgreenL] (0,0) arc[start angle=180, end angle=0, radius=0.9] -- cycle;
\draw[popgreen] (0.2,0.1) -- (0.7,0.55) -- (1.1,0.2) -- (1.5,0.6);
\draw[popgreen] (0.4,0.05) -- (0.9,0.7);
\draw[popgreen] (1.2,0.05) -- (1.6,0.35);
\node[align=center, text width=2.8cm] at (0.9,-0.7) {Recrystallized\\ calcite, coarse crystals, same composition};
\end{scope}
% Replaced shell
\begin{scope}[xshift=8.4cm]
\draw[thick, popink, fill=popgold] (0,0) arc[start angle=180, end angle=0, radius=0.9] -- cycle;
\draw[popink] (0.15,0.15) arc[start angle=180, end angle=0, radius=0.75];
\draw[popink] (0.35,0.3) arc[start angle=180, end angle=0, radius=0.55];
\node[align=center, text width=2.8cm] at (0.9,-0.7) {Replaced\\ silica or pyrite, new composition, form kept};
\end{scope}
\end{tikzpicture}
\legende{Comparison strip: an unaltered shell keeps its original mineral and microstructure; a recrystallized shell keeps its composition ($\mathrm{CaCO_3}$) but its crystals have grown coarse; a replaced shell keeps its shape but is now made of a different mineral.}
\end{popfigure}

\courssub{Distinguishing the Three Processes}

The three processes are frequently confused in examinations. The decisive questions to ask about any altered fossil are: \emph{Was anything added? Was the original material removed? Did the composition change?}

\begin{center}
\begin{tabularx}{\linewidth}{|l|X|X|X|}
\hline
\rowcolor{popblueL} \textbf{Criterion} & \textbf{Permineralization} & \textbf{Replacement} & \textbf{Recrystallization} \\
\hline
Original material & Retained (pores filled) & Dissolved and substituted & Retained \\
\hline
Chemical composition & Composite (organic + mineral) & Changed (new mineral) & Unchanged \\
\hline
Internal structure & Excellently preserved & Preserved if substitution is gradual & Often blurred (coarser crystals) \\
\hline
Typical example & Petrified wood, bone & Silicified or pyritized shells & Aragonite shell $\to$ calcite \\
\hline
\end{tabularx}
\end{center}

\begin{attention}[Common mistake: permineralization versus replacement]
Candidates frequently write that in permineralization ``the bone is replaced by silica''. This is wrong. In \cle{permineralization} the mineral only \emph{fills the empty pores}; the original bone or wood substance remains in place. In \cle{replacement} the original substance itself is \emph{dissolved away and substituted} by the new mineral. A permineralized fossil still contains original material; a replaced fossil contains none. Do not confuse either process with recrystallization, in which the composition does not change at all.
\end{attention}

\begin{exoresolu}[Identifying the mode of fossilization]
A palaeontologist collects three fossils from a sedimentary series: (a) a log of wood whose cells are filled with silica but whose carbonised cell walls are still present, showing clear growth rings; (b) an ammonite shell made entirely of iron pyrite; (c) a bivalve shell made of calcite, although living relatives of this bivalve have aragonite shells. For each fossil, state the mode of fossilization and justify your answer.
\tcblower
\textbf{Solution.}
\begin{itemize}
\item (a) \textbf{Permineralization.} The silica occupies only the cell cavities; the original (carbonised) cell walls are still present, and the growth rings prove that the original architecture was retained. Nothing was substituted --- mineral was added to empty spaces.
\item (b) \textbf{Replacement.} The original shell (which would have been $\mathrm{CaCO_3}$) has been completely dissolved and substituted molecule by molecule by pyrite ($\mathrm{FeS_2}$). The composition has changed, but the external form of the ammonite is faithfully preserved.
\item (c) \textbf{Recrystallization.} The shell is still calcium carbonate, $\mathrm{CaCO_3}$ --- the composition is unchanged --- but the unstable aragonite of the living animal has been converted into the stable polymorph calcite during burial.
\end{itemize}
\end{exoresolu}

\begin{exercice}[Applying the criteria]
A fossil coral is found in a limestone. Chemical analysis shows that the skeleton consists of silica, yet the finest details of the coral's septa are perfectly reproduced. (1) Name the mode of fossilization. (2) Explain, in three or four sentences, the sequence of events that produced this fossil. (3) State how the fossil would differ if it had undergone recrystallization instead.
\end{exercice}

\begin{corrige}[Exercise 1]
(1) \textbf{Replacement} (silicification). (2) After burial, groundwater rich in dissolved silica percolated through the coral skeleton. The original calcium carbonate was dissolved little by little, and at the same microscopic front, silica was precipitated in its place, molecule by molecule. Because dissolution and precipitation were simultaneous and gradual, the form and fine structure of the septa were faithfully reproduced in silica. (3) Under recrystallization the composition would have remained $\mathrm{CaCO_3}$ (aragonite converted to calcite), and the fine structure of the septa would probably have been blurred by the growth of coarser calcite crystals.
\end{corrige}

\begin{aretenir}[Key points]
\begin{itemize}
\item Unaltered preservation (freezing, desiccation, amber, tar pits, peat bogs) conserves soft tissues but is rare and restricted to geologically young fossils.
\item \cle{Permineralization}: minerals precipitate \emph{inside the pores} of bone or wood; original material is retained; internal anatomy superbly preserved (petrified wood and its annual rings).
\item \cle{Replacement}: original material dissolved and substituted molecule by molecule by silica, calcite, pyrite or iron oxides; composition changes, form is preserved.
\item \cle{Recrystallization}: crystal structure changes without change of composition (aragonite $\to$ calcite); fine microstructure is often lost.
\item The diagnostic questions are: was material added (permineralization), substituted (replacement), or merely reorganised (recrystallization)?
\end{itemize}
\end{aretenir}

\begin{savaistu}[Fossil wood in Cameroon]
Silicified (petrified) wood occurs in the sedimentary basins of northern Cameroon, where ancient trunks were buried by river sediments and later permineralized and replaced by silica carried in groundwater. Such fossils testify that regions which are today dry savanna once supported dense forests crossed by large rivers --- a direct illustration of how modes of fossilization allow us to reconstruct vanished environments.
\end{savaistu}

\coursec{Modes of Fossilization II: Moulds, Casts, Carbonization and Imprints}

In the previous section we studied the modes of fossilization in which at least part of the \emph{original material} of the organism survives: unaltered hard parts, frozen or mummified soft tissues, and replacement processes such as permineralization, silicification and pyritization. We now turn to a second great family of modes of fossilization, in which the fossil is essentially a \cle{record of the shape or chemistry} of the organism rather than the organism itself. When you pick up an ammonite ``shell'' from a limestone in the field, you are very often holding \emph{not} the original shell at all, but a natural stone replica of it. Understanding how such replicas — \cle{moulds} and \cle{casts} — form, and how soft organisms can survive as \cle{carbon films} and \cle{imprints}, is central to palaeontology and is frequently examined at the GCE Advanced Level.

\courssub{Moulds: Negative Records of Organisms}

\begin{definition}[Mould]
A \cle{mould} is a \textbf{negative impression} of an organism preserved in sediment or rock. It reproduces the shape of the organism in \textbf{hollow} (negative relief), like a jelly mould in a kitchen. Two types are distinguished:
\begin{itemize}
\item An \cle{external mould} is the negative imprint of the \textbf{outside} surface of an organism, formed when sediment hardens around the organism and the organism is later dissolved or decayed away.
\item An \cle{internal mould} (or \cle{steinkern}, from the German for ``stone kernel'') is the lithified sediment that filled the \textbf{inside} of a hollow shell, preserving the shape of the shell's interior.
\end{itemize}
\end{definition}

To build intuition, press a coin into soft clay and remove it: the hollow left behind is an \emph{external mould} of the coin. Now imagine a snail shell filled with wet sand: if the sand hardens and the shell dissolves, the solid sand core that remains is an \emph{internal mould} — it shows the inside of the shell, not the outside.

\begin{methode}[How an external and internal mould form]
\begin{enumerate}
\item A shelled organism (e.g. a bivalve or gastropod) dies and is \textbf{rapidly buried} by fine sediment. Sediment also enters and fills the hollow interior of the shell.
\item The sediment is \textbf{lithified} (compacted and cemented) into rock while the shell is still present.
\item Groundwater percolates through the rock and \textbf{dissolves the shell} (shells of calcium carbonate dissolve readily in slightly acidic water).
\item Two negative records remain: the \textbf{external mould} (the cavity in the surrounding rock showing the outside of the shell) and the \textbf{internal mould} or steinkern (the solid lithified filling showing the inside of the shell).
\end{enumerate}
\end{methode}

\begin{attention}[Moulds are hollow — do not confuse relief]
A mould shows \textbf{negative relief}: features that stuck out on the organism appear as depressions in the mould, and vice versa. If a specimen is a \emph{solid} three-dimensional object with positive relief, it is \textbf{not} a mould — it is a cast or a replaced shell. Examiners regularly test this distinction.
\end{attention}

\courssub{Casts: Positive Replicas}

\begin{definition}[Cast]
A \cle{cast} is a \textbf{positive, three-dimensional replica} of an organism, formed when a mould (external or internal) is filled by sediment or by minerals precipitated from solution, and this filling is later lithified. The cast reproduces the \textbf{external shape and size} of the original organism, but contains \textbf{none of its original material}.
\end{definition}

The relationship is exactly that of a plaster workshop: the mould is the negative form, and the cast is the positive object poured into it. Nature performs the same operation over geological time.

\begin{methode}[How a cast forms]
\begin{enumerate}
\item An external or internal mould exists in the rock (the original shell has dissolved).
\item Sediment is washed into the cavity, \emph{or} mineral matter (e.g. calcite, silica, iron oxides) precipitates from groundwater and fills it.
\item The filling \textbf{hardens} (lithification or cementation).
\item Erosion later splits the rock along the mould surface, revealing the solid cast — a natural statue of the organism in stone.
\end{enumerate}
\end{methode}

\begin{attention}[Common mistake: calling a cast an ``original shell'']
A cast \textbf{looks exactly like} the original shell, which is why students often describe it as ``the fossilised shell itself''. This is wrong. In a cast, the original shell material has been \textbf{completely dissolved}; the fossil is made of \emph{new} sediment or minerals that merely copy the shape. Internal details of shell structure (e.g. original micro-layers) are lost. Always state clearly: \emph{a cast reproduces form, not substance}.
\end{attention}

\begin{popfigure}
\begin{center}
\begin{tikzpicture}[scale=0.92, every node/.style={font=\small}]
% Stage 1: buried shell
\fill[popblueL] (0,0) rectangle (3.4,2.2);
\draw[thick, popink] (1.7,1.1) ellipse (0.9 and 0.55);
\draw[thick, popink] (1.7,1.1) ellipse (0.45 and 0.25);
\node[align=center] at (1.7,-0.5) {\textbf{1. Burial}\\ shell in sediment};
% Stage 2: dissolution -> moulds
\begin{scope}[xshift=4.6cm]
\fill[popblueL] (0,0) rectangle (3.4,2.2);
\draw[thick, dashed, popdark] (1.7,1.1) ellipse (0.9 and 0.55);
\fill[popgold!60] (1.7,1.1) ellipse (0.45 and 0.25);
\draw[thick, popdark] (1.7,1.1) ellipse (0.45 and 0.25);
\node[align=center] at (1.7,-0.5) {\textbf{2. Shell dissolves}\\ external mould (dashed cavity)\\ + internal mould (solid fill)};
\end{scope}
% Stage 3: cast
\begin{scope}[xshift=9.2cm]
\fill[popblueL] (0,0) rectangle (3.4,2.2);
\fill[poporange!70] (1.7,1.1) ellipse (0.9 and 0.55);
\draw[thick, popink] (1.7,1.1) ellipse (0.9 and 0.55);
\node[align=center] at (1.7,-0.5) {\textbf{3. Mould filled}\\ solid cast forms};
\end{scope}
% arrows
\draw[-{Stealth}, very thick, poppurple] (3.5,1.1) -- (4.5,1.1);
\draw[-{Stealth}, very thick, poppurple] (8.1,1.1) -- (9.1,1.1);
\end{tikzpicture}
\end{center}
\legende{Formation of moulds and a cast. (1) A shell is buried and its interior fills with sediment. (2) The shell dissolves, leaving an external mould (cavity) and an internal mould (lithified filling). (3) The external mould is later filled with sediment or minerals, producing a solid cast.}
\end{popfigure}

\courssub{Carbonization (Distillation)}

So far we have considered organisms with hard parts. What happens to a fern leaf, a fish, or an insect — organisms that are mostly soft tissue? Under the right conditions they can be preserved by \cle{carbonization}.

\begin{definition}[Carbonization]
\cle{Carbonization} (also called \emph{distillation}) is a mode of fossilization in which the soft tissues of an organism are buried and compressed under fine sediment; heat and pressure drive off the volatile components (water, oxygen, nitrogen, hydrogen), leaving behind a \textbf{thin film of carbon} that outlines the organism on the bedding surface.
\end{definition}

The process is chemically analogous to what happens when organic matter is heated without air: the ``volatile'' elements escape, and only the stable carbon skeleton remains. Because the residue is essentially pure carbon, the fossil appears as a \textbf{dark, silhouette-like film} — often showing astonishing detail such as leaf veins, insect wings or fish fins.

\begin{exemplebox}[Carbonized remains]
\begin{itemize}
\item \textbf{Fern leaves} preserved as black carbon films on dark shales, with every leaflet and vein visible.
\item \textbf{Fish} from finely laminated lake sediments, preserved as carbon outlines of the body and fins.
\item \textbf{Insects}, whose delicate wings leave detailed carbon films in fine-grained lake or lagoon deposits.
\end{itemize}
\end{exemplebox}

\begin{popfigure}
\begin{center}
\begin{tikzpicture}[scale=1.0]
% shale block
\fill[gray!25] (0,0) rectangle (7,3.4);
\draw[gray!60] (0,0.6) -- (7,0.6);
\draw[gray!60] (0,1.3) -- (7,1.3);
\draw[gray!60] (0,2.1) -- (7,2.1);
\draw[gray!60] (0,2.8) -- (7,2.8);
% carbon film fern
\draw[line width=2.2pt, popink] (3.5,0.7) -- (3.5,2.9);
\foreach \y/\l in {0.95/0.9, 1.25/1.15, 1.55/1.3, 1.85/1.3, 2.15/1.1, 2.45/0.85, 2.7/0.55}{
  \draw[line width=1.6pt, popink] (3.5,\y) -- (3.5-\l,\y+0.28);
  \draw[line width=1.6pt, popink] (3.5,\y) -- (3.5+\l,\y+0.28);
}
\node[align=center, font=\small] at (5.6,1.7) {thin carbon film\\ (leaf outline)};
\draw[-{Stealth}, poppurple] (4.9,1.7) -- (4.15,1.85);
\node[align=center, font=\small] at (1.5,3.0) {fine shale laminae};
\draw[-{Stealth}, poppurple] (1.5,2.75) -- (1.5,2.2);
\end{tikzpicture}
\end{center}
\legende{A carbonized fern leaf on shale. Volatile elements have been driven off by compression, leaving a thin dark carbon film that faithfully preserves the outline of the stem and leaflets between the fine laminae.}
\end{popfigure}

\courssub{Imprints and Compressions}

Closely related to carbonization are \cle{imprints} and \cle{compressions}, typical of delicate organisms buried in very fine-grained, thinly laminated sediments such as lake muds, volcanic ash or lagoonal silts.

\begin{definition}[Imprint and compression]
An \cle{imprint} is a shallow surface impression left by a delicate organism or plant part pressed into fine sediment before lithification — even if no organic material remains. A \cle{compression} is the \textbf{flattened} remain of an organism squeezed between sediment laminae during compaction; it may retain a carbon film (linking it to carbonization) or only the flattened outline.
\end{definition}

The distinction is subtle but examinable: in a \emph{compression}, some flattened organic residue (often carbon) is usually still present between the laminae; in a pure \emph{imprint}, only the surface marking remains, like a footprint of the leaf on the rock. Both require \textbf{rapid burial in quiet water} so that fine laminae can drape over the organism without disturbing it — conditions found, for example, on the beds of crater lakes such as those of the Cameroon Volcanic Line, where fine ash and mud accumulate undisturbed.

\courssub{Tracks and Burrows: Fossilized Activity}

A final mode of fossilization records not the organism's body but its \textbf{activity}. When an animal walks across soft mud, its footprints may be buried by the next layer of sediment and preserved as \cle{tracks}; when a worm burrows through the sediment, the burrow may be filled with different sediment and preserved as a \cle{burrow}. These are the \cle{trace fossils} (ichnofossils) introduced in Section 3 — here we emphasise that their \emph{mode of fossilization} is essentially that of an imprint or mould of behaviour: the sediment records the disturbance, and lithification fixes it permanently. Dinosaur trackways, trilobite crawling traces and worm burrows are classic examples.

\courssub{Worked Example and Summary}

\begin{exoresolu}[Identifying modes of fossilization]
A student collects three specimens from a shale quarry: (a) a solid stone object with the exact external shape of an ammonite, made of crystalline calcite; (b) a hollow cavity in the rock showing the ribbed outside surface of a bivalve shell; (c) a black film on a bedding plane showing the outline of a fish. For each specimen, state the mode of fossilization and justify your answer.
\tcblower
\textbf{(a) Cast.} The specimen is solid, three-dimensional, and made of new mineral matter (calcite) that filled a mould after the original shell dissolved. It reproduces the external form but contains no original shell material.

\textbf{(b) External mould.} The specimen is a \emph{cavity} (negative relief) reproducing the outside surface of the shell. The shell itself dissolved away, leaving the imprint in the surrounding rock.

\textbf{(c) Carbonization.} The fish is preserved as a thin carbon film: heat and pressure drove off the volatile elements from the soft tissues, leaving only carbon outlining the body and fins on the bedding plane.
\end{exoresolu}

\begin{aretenir}[Summary of the modes of fossilization in this section]
\begin{center}
\begin{tabularx}{\linewidth}{|l|X|l|l|}
\hline
\rowcolor{popblueL} \textbf{Mode} & \textbf{Process} & \textbf{Type of organism} & \textbf{Example} \\
\hline
External mould & Organism buried; sediment hardens around it; organism dissolves, leaving a negative imprint of the outside & Shelled organisms & Bivalve or ammonite impression in shale \\
\hline
Internal mould (steinkern) & Sediment fills the inside of a hollow shell and lithifies; shell later dissolves & Hollow-shelled organisms & Gastropod steinkern \\
\hline
Cast & A mould is filled with sediment or minerals which harden, giving a positive 3D replica & Shelled organisms & Calcite cast of an ammonite \\
\hline
Carbonization & Compression drives off volatiles; a thin carbon film of the soft tissues remains & Soft-bodied or delicate organisms & Fern leaf, fish, insect films \\
\hline
Imprint / compression & Delicate remains flattened between fine laminae; outline (with or without carbon) preserved & Plants, insects & Leaves and insects in laminated lake shale \\
\hline
Tracks and burrows & Sediment records the activity of an animal; lithification fixes the trace & Mobile or burrowing animals & Dinosaur tracks, worm burrows \\
\hline
\end{tabularx}
\end{center}
\end{aretenir}

\begin{methode}[Distinguishing a mould from a cast on a specimen]
\begin{enumerate}
\item \textbf{Check the relief.} Run your finger over the fossil. If the features are \emph{hollow} (negative relief — ribs appear as grooves), it is a \textbf{mould}. If they stand out (positive relief), it is a \textbf{cast} or replaced shell.
\item \textbf{Check the material.} If the fossil is made of different sediment or mineral from the original shell (e.g. calcite or sandstone in the shape of a shell), it is a \textbf{cast}.
\item \textbf{Check the geometry.} A cast is a solid 3D object that fits into the mould like a key into a keyhole; a mould is the cavity itself.
\item \textbf{Conclude and justify} using the vocabulary: negative vs positive relief, original material absent, filling of a cavity.
\end{enumerate}
\end{methode}

\begin{exercice}[Moulds, casts and carbon films]
\begin{enumerate}
\item Define an internal mould (steinkern) and explain, in four steps, how it forms.
\item A fossil shell is made entirely of silica yet shows perfect external ornamentation. Explain why this could be either a replaced shell or a cast, and state what additional observation would help you decide.
\item Explain why carbonization is especially important for the fossil record of soft-bodied organisms, naming two types of organisms commonly preserved this way.
\end{enumerate}
\end{exercice}

\begin{corrige}[Exercise 1]
\textbf{1.} An internal mould (steinkern) is the lithified sediment that filled the inside of a hollow shell, preserving the shape of the shell's interior. Formation: (i) the shelled organism dies and is buried, and sediment enters and fills the shell's cavity; (ii) the sediment inside hardens during lithification; (iii) groundwater dissolves the shell wall; (iv) the solid internal filling remains as the steinkern.

\textbf{2.} A replaced shell keeps the original shell \emph{structure and position}, with silica substituting the original carbonate atom by atom; a cast is silica that merely filled a mould after the shell dissolved, so fine internal shell micro-structure is absent. Examining a broken or thin section for preserved original shell micro-layers would distinguish replacement (layers present) from a cast (layers absent).

\textbf{3.} Soft-bodied organisms lack hard parts and normally decay completely, so carbonization is often their \emph{only} chance of entering the fossil record: compression drives off volatiles and leaves a carbon film recording their outline. Common examples: leaves (e.g. ferns) and insects (also fish).
\end{corrige}

\begin{savaistu}[Natural casts in your kitchen]
The mould--cast relationship is exactly the one used to make \emph{beignet} moulds or plaster figurines: a negative form is filled to produce a positive copy. Nature made such copies millions of years before humans did — every ammonite cast in a limestone is a statue sculpted by groundwater and time, not by the animal itself.
\end{savaistu}

\coursec{Preservation Potential and Occurrence of Fossils}

In the previous section we saw \emph{how} an organism can be transformed into a fossil: by preservation of unaltered remains, by permineralization and replacement, or by the formation of moulds, casts, carbon films and imprints. But a natural question now arises: \emph{why do some creatures fossilize abundantly while others almost never do?} Walk along any fossiliferous outcrop --- for example the Cretaceous limestones of the Douala Basin --- and you will find countless shells of molluscs, but you will search in vain for the jellyfish, worms and soft-bodied shrimps that certainly lived in the same sea. The answer lies in the concept of \cle{preservation potential}, and in the rules that govern \emph{where} fossils actually occur in the Earth's crust. These two ideas are the subject of this section.

\courssub{Preservation Potential}

\textbf{Intuition first.} Imagine two animals dying side by side on a beach near Limbe: a clam with a thick calcium carbonate shell, and a jellyfish made of more than $95\%$ water. After burial, scavengers and bacteria attack both. The jellyfish's tissues decompose within days; nothing solid remains to be mineralized. The clam's shell, however, resists decay, dissolution and transport, and can survive for millions of years. The two organisms had very different \emph{chances} of entering the fossil record.

\begin{definition}[Preservation potential]
The \cle{preservation potential} of an organism (or of one of its parts) is the probability that it will survive decay, destruction and transport long enough to be buried and fossilized, and thus appear in the fossil record. It depends on the nature of the organism's tissues and skeleton, its size, its abundance, its habitat, and the environment of deposition.
\end{definition}

\courssub{Factors Controlling Preservation Potential}

\textbf{1. Hard parts versus soft parts.} This is the dominant factor. \cle{Hard parts} --- shells, bones, teeth, tests, exoskeletons, wood --- resist biological decay and mechanical destruction, and therefore fossilize readily. \cle{Soft parts} --- skin, muscles, internal organs, leaves, jelly-like bodies --- are eaten by scavengers or destroyed by bacteria within days or weeks, and are preserved only under exceptional conditions (rapid burial, anoxia, early mineralization).

\textbf{2. Durability of skeletal mineralogy.} Not all hard parts are equal. The resistance of a skeleton depends on its mineral composition:

\begin{center}
\begin{tabularx}{\linewidth}{|l|X|l|}
\hline
\rowcolor{popblueL} \textbf{Skeletal material} & \textbf{Examples} & \textbf{Durability} \\
\hline
Calcite (low-Mg) & brachiopod shells, oyster shells, trilobite cuticle & Very high \\
\hline
Calcium phosphate & vertebrate bones and teeth, conodonts & Very high \\
\hline
Silica (\ce{SiO2}) & diatom and radiolarian tests, sponge spicules & High \\
\hline
Aragonite & many gastropod and bivalve shells, corals & Moderate (dissolves or recrystallizes to calcite) \\
\hline
Chitin / organic cuticle & insect exoskeletons, graptolite periderm & Low to moderate \\
\hline
Organic soft tissues & jellyfish, worms, muscles & Very low \\
\hline
\end{tabularx}
\end{center}

\textbf{3. Size and robustness.} Large, thick-shelled or thick-boned organisms withstand transport and burial compaction better than tiny, fragile ones --- although very small fossils with resistant skeletons (foraminifera, pollen, spores) are preserved in enormous numbers because each individual is produced in millions.

\textbf{4. Abundance.} A species represented by billions of individuals (e.g.\ planktonic foraminifera) is statistically far more likely to leave fossils than a rare species with few individuals (e.g.\ a large predator or a bird).

\textbf{5. Habitat.} \cle{Marine} organisms, especially those living on or in the sea floor (benthos), have a high preservation potential: they live exactly where sediment accumulates, so burial is immediate and gentle. \cle{Terrestrial} organisms have a low potential: on land, erosion dominates over deposition, carcasses are exposed to scavengers, weather and oxidation, and only those swept into lakes, rivers, swamps, tar seeps or caves stand a chance of fossilization.

\begin{exemplebox}[High and low preservation potential]
\textbf{High potential:} mollusc shells (bivalves, gastropods, ammonites), brachiopods, echinoderms (crinoids, echinoids), foraminifera, ostracods, vertebrate teeth and bones.\\
\textbf{Medium potential:} corals and aragonitic molluscs (aragonite may dissolve, leaving moulds), fish skeletons, plant leaves and wood.\\
\textbf{Low potential:} jellyfish, worms, soft-bodied arthropods, insects, birds (light hollow bones, rare, terrestrial), flowers.
\end{exemplebox}

\begin{popfigure}
\begin{center}
\begin{tikzpicture}[scale=0.92]
\draw[->,thick] (0,0) -- (10.6,0) node[below left=2pt and -30pt,align=center,text width=3.2cm]{Preservation potential};
\draw[->,thick] (0,0) -- (0,5.6);
\foreach \y/\lab in {1/Low,3/Medium,5/High}{
  \draw (0,\y) node[left]{\lab};
  \draw[gray!40,dashed] (0,\y) -- (10.4,\y);
}
% High bars
\fill[popgreen] (0.4,0) rectangle (1.6,5);
\node[rotate=90,white,align=center,text width=4.4cm] at (1.0,2.5){\textbf{Molluscs, brachiopods}};
\fill[popgreen] (1.8,0) rectangle (3.0,4.8);
\node[rotate=90,white,align=center,text width=4.2cm] at (2.4,2.4){\textbf{Echinoderms, foraminifera}};
\fill[popgreen] (3.2,0) rectangle (4.4,4.6);
\node[rotate=90,white,align=center,text width=4.0cm] at (3.8,2.3){\textbf{Vertebrate bones \& teeth}};
% Medium bars
\fill[popgold] (4.8,0) rectangle (6.0,3);
\node[rotate=90,white,align=center,text width=2.4cm] at (5.4,1.5){\textbf{Corals, aragonitic shells}};
\fill[popgold] (6.2,0) rectangle (7.4,2.8);
\node[rotate=90,white,align=center,text width=2.2cm] at (6.8,1.4){\textbf{Plants, fish}};
% Low bars
\fill[poppink] (7.8,0) rectangle (9.0,1);
\node[rotate=90,white,align=center,text width=0.8cm] at (8.4,0.5){\textbf{Insects, birds}};
\fill[poppink] (9.2,0) rectangle (10.4,0.6);
\node[rotate=90,white,align=center,text width=0.5cm,font=\small] at (9.8,0.3){\textbf{Jellyfish, worms}};
\end{tikzpicture}
\end{center}
\legende{Relative preservation potential of the main organism groups, controlled by skeletal mineralogy, robustness, abundance and habitat.}
\end{popfigure}

\begin{attention}[Common mistake]
Do not claim that ``only organisms with hard parts can be fossilized''. Soft-bodied organisms \emph{can} be preserved, but only under exceptional conditions (rapid burial in anoxic mud, carbonization, amber). The correct statement is that their preservation potential is \emph{very low}, not zero.
\end{attention}

\courssub{Occurrence of Fossils: Which Rocks Contain Them?}

\textbf{Observation.} A student collecting around Mount Cameroon finds beautiful crystals in basaltic lava, but never a shell. In the quarries of Figuil (North Region), however, the grey limestone is packed with marine shells. Why this contrast?

\begin{propriete}[Occurrence of fossils in rocks]
Fossils occur almost exclusively in \cle{sedimentary rocks}, because only sedimentary processes bury organic remains gently, rapidly and at low temperature. In \cle{igneous rocks}, the heat of magma (several hundred to over $1000\ ^{\circ}\mathrm{C}$) destroys all organic material, so fossils are never found in granite or basalt --- the only exceptions are remains trapped in \emph{volcanic ash} (a sedimentary deposit) or in \emph{amber}. In \cle{metamorphic rocks}, recrystallization under heat and pressure normally destroys fossils; only \emph{low-grade} metamorphic rocks (e.g.\ slates) may exceptionally retain deformed fossils such as graptolites.
\end{propriete}

\begin{attention}[Trap in the GCE examination]
A classic question asks: ``Why are fossils absent from igneous rocks?'' Many candidates answer vaguely ``because igneous rocks are hard''. The rigorous answer is: magma temperatures destroy organic tissues and skeletons, and crystallization from a melt provides no gentle burial process. Never write that fossils occur ``in all rock types''.
\end{attention}

\courssub{Distribution of Fossils within Strata}

Even within a sedimentary succession, fossils are not spread uniformly. They are concentrated in particular beds:

\begin{itemize}
\item \textbf{Fossiliferous horizons:} beds rich in fossils, alternating with \emph{barren} (fossil-poor) beds, reflecting changes in environment, sediment supply or water depth.
\item \textbf{Shell beds (coquinas):} layers where currents or storm waves concentrated shells into dense accumulations.
\item \textbf{Lagerstätten:} rare sites of truly exceptional preservation.
\end{itemize}

\begin{definition}[Lagerstätte]
A \cle{Lagerstätte} (plural \emph{Lagerstätten}, from German ``storage place'') is a sedimentary deposit exhibiting \emph{exceptional} fossil preservation, either in the \emph{quality} of preservation (soft tissues, complete articulated skeletons --- \emph{Konservat-Lagerstätte}) or in the \emph{quantity} of fossils concentrated (\emph{Konzentrat-Lagerstätte}).
\end{definition}

\begin{savaistu}[Famous Lagerstätten]
The \textbf{Burgess Shale} (British Columbia, Canada, Middle Cambrian) preserves soft-bodied marine animals --- worms, arthropods and strange forms with no living relatives --- as flattened films, revealing the explosion of animal life about $505$ million years ago. The \textbf{Solnhofen Limestone} (Bavaria, Germany, Late Jurassic) is a fine-grained lithographic limestone deposited in quiet, salty lagoons; it has yielded the famous bird \emph{Archaeopteryx} with feather impressions, as well as jellyfish and insects. Such sites show that the ``ordinary'' fossil record gives a strongly biased, hard-part-dominated picture of past life.
\end{savaistu}

\begin{popfigure}
\begin{center}
\begin{tikzpicture}[scale=1.0]
% Stratigraphic column
\fill[gray!25] (0,0) rectangle (3,1.2);
\fill[popblueL] (0,1.2) rectangle (3,2.0);
\fill[gray!25] (0,2.0) rectangle (3,3.2);
\fill[popgreenL] (0,3.2) rectangle (3,4.4);
\fill[gray!25] (0,4.4) rectangle (3,5.4);
\fill[popgold!50] (0,5.4) rectangle (3,6.2);
\draw[thick] (0,0) rectangle (3,6.2);
% Fossil symbols
\foreach \p in {(0.6,1.6),(1.4,1.5),(2.2,1.7),(1.0,1.8),(2.6,1.4)}{
  \draw[popblue,thick] \p circle (0.09);}
\foreach \p in {(0.5,3.7),(1.3,3.9),(2.1,3.6),(2.6,4.1),(0.9,4.2),(1.7,3.4)}{
  \draw[popgreen,thick] \p circle (0.09);}
\foreach \p in {(0.8,5.8),(1.8,5.7),(2.5,5.9)}{
  \draw[popdark,thick] \p circle (0.09);}
% Labels
\node[right,align=left,text width=5.6cm] at (3.2,5.8){\textbf{Shell bed (coquina)}: storm-concentrated shells --- very rich horizon};
\node[right,align=left,text width=5.6cm] at (3.2,4.9){Barren sandstone (no fossils)};
\node[right,align=left,text width=5.6cm] at (3.2,3.8){\textbf{Fossiliferous limestone}: abundant echinoids and bivalves};
\node[right,align=left,text width=5.6cm] at (3.2,2.6){Barren shale (unfossiliferous)};
\node[right,align=left,text width=5.6cm] at (3.2,1.6){\textbf{Fossiliferous marl}: ammonites, foraminifera};
\node[right,align=left,text width=5.6cm] at (3.2,0.6){Barren conglomerate (high energy, destructive)};
\draw[->,thick] (-0.4,0) -- (-0.4,6.2) node[above,align=center,text width=2cm]{Younger};
\end{tikzpicture}
\end{center}
\legende{Stratigraphic column showing fossiliferous beds alternating with barren beds: fossil occurrence is controlled by depositional environment and energy.}
\end{popfigure}

\courssub{Field Techniques: From Outcrop to Museum Drawer}

Finding a fossil is only the beginning. The GCE examiner expects you to outline, briefly and in order, the four stages of palaeontological fieldwork:

\begin{methode}[Prospecting, extraction, preparation, cataloguing]
\begin{enumerate}
\item \textbf{Prospecting:} locate promising fossiliferous sedimentary outcrops using geological maps; examine cliffs, quarries, road cuts and river banks; note the exact bed from which each fossil comes.
\item \textbf{Extraction:} remove fossils carefully with hammer, chisel and brush; wrap delicate specimens in paper or plaster jackets; large bones are excavated with the surrounding rock.
\item \textbf{Preparation:} in the laboratory, clean the fossil mechanically (needles, drills, air-abrasive tools) or chemically (dilute acid to free calcite fossils from limestone); harden fragile specimens with consolidant.
\item \textbf{Cataloguing:} give each specimen a reference number and record its name, exact locality, stratigraphic horizon, age, date of collection and collector --- without these data a fossil loses most of its scientific value.
\end{enumerate}
\end{methode}

\courssub{Fossil Occurrences in Cameroon}

Cameroon offers excellent illustrations of these principles:

\begin{itemize}
\item The \textbf{Douala Basin} and the \textbf{Rio del Rey Basin} (Atlantic coastal sedimentary basins) contain Cretaceous to Cenozoic marine sediments --- limestones, marls and sandstones --- rich in ammonites, bivalves, gastropods, echinoids and abundant foraminifera, the latter being essential for dating strata during petroleum exploration.
\item The \textbf{Figuil limestone} (North Region, Benue Trough domain) is a Cretaceous marine limestone exploited for cement; it yields bivalves (including rudists), gastropods and echinoids --- a typical fossiliferous horizon.
\item By contrast, the volcanic rocks of the \textbf{Cameroon Volcanic Line} (Mount Cameroon, Mount Manengouba, the Bambouto Mountains) contain no fossils, exactly as the theory predicts.
\end{itemize}

\begin{exoresolu}[Preservation potential and rock type]
\textbf{Statement.} During fieldwork, a student collects samples from four sites: (A) a granite inselberg near Ngaoundéré; (B) Cretaceous limestone at Figuil; (C) a recent basalt flow of Mount Cameroon; (D) marine marl of the Douala Basin. (1) At which sites can fossils reasonably be expected? Justify. (2) The student finds abundant bivalve shells but no worms at site D, although worms must have lived there. Explain.\\
\tcblower
\textbf{Solution.}\\
(1) Fossils can be expected at \textbf{B} (Figuil limestone) and \textbf{D} (Douala marl): both are \emph{sedimentary rocks} formed by gentle burial of sediment at low temperature, which allows organic remains to be preserved. No fossils can be found at \textbf{A} (granite) or \textbf{C} (basalt): these are \emph{igneous rocks} crystallized from hot magma, whose heat destroys all organic material.\\
(2) The bivalves possessed \emph{hard} calcium carbonate shells --- a high preservation potential: the shells resisted decay and were buried and fossilized. The worms were \emph{soft-bodied} --- a very low preservation potential: their tissues were destroyed by scavengers and bacterial decay before burial could protect them. The fossil assemblage is therefore \emph{biased} towards hard-part organisms and does not represent the complete original community.
\end{exoresolu}

\begin{aretenir}[Key points]
\begin{itemize}
\item \cle{Preservation potential} = probability that an organism becomes a fossil; it is controlled by hard parts, skeletal mineralogy, size, abundance and habitat.
\item High potential: molluscs, brachiopods, echinoderms, foraminifera, bones and teeth. Low potential: jellyfish, worms, insects, birds.
\item Fossils occur almost exclusively in \cle{sedimentary rocks}; igneous heat destroys remains, and metamorphism deforms or obliterates them (rare survival in low-grade rocks).
\item Fossils concentrate in fossiliferous horizons, shell beds and exceptional \cle{Lagerstätten} (Burgess Shale, Solnhofen).
\item Fieldwork follows four stages: prospecting, extraction, preparation, cataloguing.
\item Cameroonian examples: Cretaceous marine fossils of the Douala and Rio del Rey basins and of the Figuil limestone; none in the volcanic rocks of the Cameroon Volcanic Line.
\end{itemize}
\end{aretenir}

\begin{exercice}[Check your understanding]
\begin{enumerate}
\item Define \emph{preservation potential} and state four factors that control it.
\item Explain why foraminifera, although microscopic, have an excellent fossil record.
\item A candidate writes: ``Fossils can be found in any rock, including basalt, provided the organism was buried quickly.'' Correct this statement.
\item What is a Lagerstätte? Give two world-famous examples and state what makes each exceptional.
\item Why are the limestones of Figuil fossiliferous while the lavas of Mount Cameroon are not?
\end{enumerate}
\end{exercice}

\begin{corrige}[Exercise 7]
\begin{enumerate}
\item The probability that an organism or structure survives decay and destruction to become a fossil. Factors: presence and mineralogy of hard parts, size/robustness, abundance of the species, habitat (marine benthic versus terrestrial).
\item Each individual has a resistant calcite test, and each species produces enormous numbers of individuals living in marine settings where sediment accumulates; the combination of hard parts, abundance and habitat gives a high preservation potential.
\item False. Fossils are found essentially in sedimentary rocks. Basalt is igneous: magma temperatures destroy organic remains, so no fossil can survive in it, whatever the speed of burial. (Only remains in volcanic ash or amber, which are sedimentary/resinous materials, are exceptions.)
\item A deposit of exceptional fossil preservation, in quality (soft tissues, complete skeletons) or quantity. Burgess Shale (Cambrian, Canada): soft-bodied animals preserved as films. Solnhofen Limestone (Jurassic, Germany): fine lagoonal limestone preserving \emph{Archaeopteryx} with feathers, jellyfish and insects.
\item Figuil limestone is a marine sedimentary rock: shells were buried gently in lime mud at low temperature and fossilized. Mount Cameroon lavas crystallized from magma at very high temperature, which destroys any organic material; hence they are barren of fossils.
\end{enumerate}
\end{corrige}

\coursec{Uses of Fossils, Gaps in the Fossil Record and Extinction}

In the previous section we saw that only a small fraction of dead organisms escapes destruction and enters the rock record, and that fossils are found almost exclusively in sedimentary rocks, concentrated in particular beds. A natural question now arises: once we have found these fossils, \emph{what are they good for?} The answer is remarkable. Fossils are at the same time \cle{clocks} (they tell the age of rocks), \cle{thermometers and salinity meters} (they record ancient environments), \cle{maps} (they help reconstruct vanished continents), \cle{archives of evolution}, and \cle{tools of economic exploration}. But the archive is incomplete: it contains gaps, and it records the disappearance of species through \cle{extinction}. This section explores these three great themes: the uses of fossils, the gaps in the fossil record, and extinction.

\courssub{Stratigraphic Uses of Fossils: Index Fossils and Relative Dating}

\textbf{The problem.} Imagine two road cuttings, one near Bafoussam and one near Bertoua. Both expose layered sedimentary rocks. Were the layers deposited at the same time? The rocks look different, so lithology alone cannot answer. But if both cuttings contain the \emph{same} fossil species, and that species lived only during a short, well-defined interval of geological time, then the two layers must be of the same age. Fossils thus allow us to \emph{date rocks relatively} and to \emph{correlate} strata across distances.

This reasoning rests on a fundamental principle established by the English engineer William Smith around 1800.

\begin{propriete}[Principle of Faunal Succession]
Fossil organisms succeed one another through geological time in a \textbf{definite, irreversible and recognisable order}. Each interval of Earth history is characterised by a distinctive assemblage of fossil species, so that sedimentary strata can be identified, dated relatively and correlated by their fossil content. (The principle itself is examinable; no formal proof is required, but you must be able to apply it to a correlation exercise.)
\end{propriete}

Not every fossil is useful for dating. A species that lived for hundreds of millions of years, or that existed only in one small lake, tells us little. The ideal time-marker is the \cle{index fossil} (or \cle{zone fossil}).

\begin{definition}[Index Fossil]
An \textbf{index fossil} (zone fossil) is a fossil of an organism that satisfies four criteria:
\begin{enumerate}
\item \textbf{Wide geographic distribution} -- the species lived over a large area (ideally worldwide), so it can be found in many basins.
\item \textbf{Short vertical (stratigraphic) range} -- the species evolved and became extinct within a brief interval of geological time, so its presence pins the rock to a narrow age band.
\item \textbf{Abundance} -- the species is common, so it is likely to be found even in small samples.
\item \textbf{Ease of identification} -- the species has distinctive features, so it can be recognised quickly and unambiguously, even from fragmentary material.
\end{enumerate}
Classic examples: \emph{ammonites} (Mesozoic), \emph{graptolites} (Lower Palaeozoic), \emph{trilobites} (Cambrian--Ordovician), and planktonic \emph{foraminifera} (Cretaceous--Recent).
\end{definition}

\begin{methode}[Correlating Strata Between Two Outcrops Using Fossils]
\begin{enumerate}
\item Collect and identify the fossils in each bed of outcrop A and outcrop B.
\item Separate \textbf{index fossils} (short-ranging, widespread) from long-ranging or facies-controlled species.
\item Match the beds that contain the \emph{same} index fossil species or the same fossil assemblage (biozone).
\item Draw tie-lines between the matched beds on the two columns; these lines represent surfaces of equal age.
\item Interpret any missing biozone in one column as a gap (non-deposition or erosion) and any repeated assemblage as a possible fault or fold.
\end{enumerate}
\end{methode}

\begin{popfigure}
\begin{tikzpicture}[scale=0.9, every node/.style={font=\small}]
% Column A
\fill[popblueL] (0,0) rectangle (2,1.5);
\fill[popgreenL] (0,1.5) rectangle (2,3);
\fill[popgold!40] (0,3) rectangle (2,4.5);
\draw[thick] (0,0) rectangle (2,4.5);
\node[align=center] at (1,0.75) {shale\\ \emph{Grapt.~A}};
\node[align=center] at (1,2.25) {limestone\\ \emph{Tril.~B}};
\node[align=center] at (1,3.75) {sandstone\\ \emph{Amm.~C}};
\node[above] at (1,4.5) {\textbf{Outcrop A}};
% Column B
\fill[popblueL] (4,0.5) rectangle (6,2);
\fill[popgold!40] (4,2) rectangle (6,3.5);
\draw[thick] (4,0.5) rectangle (6,3.5);
\node[align=center] at (5,1.25) {shale\\ \emph{Grapt.~A}};
\node[align=center] at (5,2.75) {sandstone\\ \emph{Amm.~C}};
\node[above] at (5,3.5) {\textbf{Outcrop B}};
\node[below, align=center] at (5,0.4) {gap: \emph{Tril.~B}\\ zone missing};
% Column C
\fill[popgreenL] (8,0) rectangle (10,1.2);
\fill[popgold!40] (8,1.2) rectangle (10,2.7);
\fill[gray!30] (8,2.7) rectangle (10,3.9);
\draw[thick] (8,0) rectangle (10,3.9);
\node[align=center] at (9,0.6) {limestone\\ \emph{Tril.~B}};
\node[align=center] at (9,1.95) {sandstone\\ \emph{Amm.~C}};
\node[align=center] at (9,3.3) {younger\\ beds};
\node[above] at (9,3.9) {\textbf{Outcrop C}};
% Tie lines
\draw[thick, poporange, dashed] (2,1.5) -- (4,2);
\draw[thick, poporange, dashed] (2,3) -- (4,2);
\draw[thick, poporange, dashed] (2,4.5) -- (4,3.5);
\draw[thick, poppurple, dashed] (6,2) -- (8,1.2);
\draw[thick, poppurple, dashed] (6,3.5) -- (8,2.7);
\end{tikzpicture}
\legende{Correlation of three stratigraphic columns using index fossils. Dashed tie-lines join beds containing the same index fossil (\emph{Graptolite A}, \emph{Trilobite B}, \emph{Ammonite C}). In outcrop B the \emph{Trilobite B} zone is missing: a stratigraphic gap.}
\end{popfigure}

\begin{attention}[Facies Fossils Are Not Clocks]
A very common GCE mistake is to use \emph{any} common fossil for dating. A \cle{facies fossil} is tied to a particular \emph{environment} (e.g. a coral that lives only in warm shallow reefs) and may have existed for a very long time. Finding it in two places only proves the two places had similar environments -- \textbf{not} that they are the same age. Rule to remember: \textbf{index fossils give age; facies fossils give environment.}
\end{attention}

\courssub{Palaeoenvironmental and Palaeogeographic Reconstruction}

Because most organisms are adapted to specific conditions, fossils are sensitive recorders of ancient environments. This is the second great use of palaeontology.

\begin{definition}[Facies Fossil]
A \textbf{facies fossil} is a fossil (or fossil assemblage) restricted to a particular sedimentary environment. Because its occurrence is controlled by environment rather than by time, it is used to reconstruct \textbf{palaeoenvironments}: water depth, salinity, temperature, oxygenation and climate.
\end{definition}

Typical indicators used in palaeoenvironmental analysis:

\begin{center}
\rowcolors{2}{popblueL}{white}
\begin{tabularx}{\linewidth}{|l|X|}
\hline
\rowcolor{popblueL}\textbf{Fossil evidence} & \textbf{Environmental interpretation} \\
\hline
Reef corals, calcareous algae & Warm ($>20\ ^\circ$C), shallow, clear, well-lit marine water of normal salinity \\
\hline
Planktonic foraminifera only & Deep, open-marine conditions, offshore \\
\hline
Oysters, some bivalves, few species but many individuals & Brackish water (estuary, lagoon): reduced or fluctuating salinity \\
\hline
Thick-shelled, ribbed bivalves and gastropods & High-energy shallow water (surf zone) \\
\hline
Thin-shelled, delicate bivalves, graptolites & Quiet, deeper, possibly poorly oxygenated water \\
\hline
Coal seams with plant roots and tree trunks & Warm, humid continental climate; swamp environment \\
\hline
Evaporite minerals with salt-tolerant microfossils & Hot, arid climate; restricted basin \\
\hline
\end{tabularx}
\end{center}

\textbf{Palaeogeography and palaeoclimatology.} By plotting the distribution of climatic indicators (reefs, coals, evaporites, glacial deposits with their fossils) on maps for successive geological periods, geologists reconstruct the positions of ancient continents and climatic belts. For example, identical fossil land reptiles and plants (such as the fern \emph{Glossopteris}) on continents now separated by oceans were key evidence that these continents were once joined -- a palaeontological argument for continental drift. Similarly, finding warm-water coral fossils in rocks of regions that are cold today shows that continents have moved across climatic zones.

\begin{savaistu}[Fossils and the Cameroon Volcanic Line Region]
Sedimentary basins flanking the Cameroon Volcanic Line, such as the Douala and Rio del Rey basins, preserve rich Cretaceous and Tertiary marine microfossils. Planktonic foraminifera recovered from boreholes allow geologists to date the strata and to reconstruct the progressive opening of the South Atlantic and the marine incursions onto the West African margin.
\end{savaistu}

\courssub{Fossils as Evidence for Evolution}

The ordered succession of fossils (faunal succession) is itself the strongest historical evidence that life has \emph{changed} through time. Two kinds of evidence are examinable:

\begin{itemize}
\item \textbf{Transitional forms}: fossils displaying a mixture of ancestral and derived characters, linking major groups. The classic example is \emph{Archaeopteryx} (Late Jurassic), which combines reptilian teeth and a long bony tail with bird feathers and wings -- a bridge between reptiles and birds.
\item \textbf{Evolutionary lineages}: continuous fossil series showing gradual morphological change. The \textbf{horse lineage} (from the small, multi-toed \emph{Eohippus} of the Eocene to the large, single-toed modern \emph{Equus}) documents reduction of toes and enlargement of grinding teeth as forests gave way to grasslands. \textbf{Ammonites} show rapid, stepwise evolution of shell coiling and suture patterns, which is precisely why they are such excellent index fossils: fast evolution means short ranges.
\end{itemize}

\courssub{Economic Uses of Fossils}

Fossils are not only academic objects; they are working tools of the petroleum and mining industries.

\begin{itemize}
\item \textbf{Biostratigraphy in drilling.} When an oil well is drilled, cuttings and cores are examined for \textbf{microfossils} (foraminifera, ostracods, pollen, spores, nannofossils). Because microfossils are tiny, abundant and rapidly evolving, a few millimetres of rock suffice to date a layer and to correlate it between wells. This guides the driller toward reservoir horizons of known age.
\item \textbf{Palaeoenvironmental prediction.} Fossil assemblages indicate the depositional environment, helping to locate porous reservoir sands (shoreface environments) and organic-rich source rocks (quiet, oxygen-poor basins).
\item \textbf{Coal and limestone resources.} Coal is itself a fossil fuel made of accumulated plant remains; its quality is assessed through fossil spores and plant fragments. Many limestones are largely built of fossil shells, corals and calcareous algae, and their fossil content controls their industrial quality.
\end{itemize}

\begin{savaistu}[Microfossils and Oil in the Douala Basin]
In the Douala (Kribi--Campo) basin of Cameroon, petroleum exploration relies heavily on foraminifera and palynomorphs (pollen and spores) extracted from borehole samples. These microfossils date the Cretaceous--Tertiary strata, identify marine source-rock intervals, and allow correlation between wells drilled tens of kilometres apart -- a direct industrial application of the principle of faunal succession.
\end{savaistu}

\courssub{Gaps in the Fossil Record}

If fossils are archives of life, the archive is frustratingly incomplete. Charles Darwin already worried that the "missing links" of his time might simply reflect a poor record.

\begin{propriete}[Incompleteness of the Fossil Record]
The fossil record samples only a \textbf{tiny fraction of past life}. The overwhelming majority of organisms that ever lived left no fossil trace, and the majority of species that ever existed are unknown to science. The record is strongly biased toward organisms with hard parts, living in depositional (especially shallow-marine) environments. (Examinable as a discussion question: state and explain the causes of this incompleteness.)
\end{propriete}

\textbf{Causes of the gaps.}

\begin{enumerate}
\item \textbf{Non-deposition}: in uplifted or terrestrial areas, sediments are not laid down, so no record forms at all.
\item \textbf{Erosion}: previously deposited, fossil-bearing strata are worn away, deleting pages of the archive (unconformities).
\item \textbf{Dissolution}: shells of calcium carbonate or silica dissolve in acidic waters or during deep burial; aragonitic shells are especially vulnerable.
\item \textbf{Metamorphism and deformation}: heat and pressure recrystallise or obliterate fossils; tectonic deformation distorts them beyond recognition.
\item \textbf{Soft-bodied organisms}: organisms lacking hard parts (worms, jellyfish, most insects, many plants) fossilise only under exceptional conditions, so entire branches of the tree of life are nearly invisible.
\item \textbf{Incomplete sampling}: many fossiliferous rocks are buried deep, hidden under vegetation or water, or simply never examined by a palaeontologist; many collected fossils remain unstudied in museum drawers.
\end{enumerate}

\textbf{Consequences.} Evolutionary lineages appear discontinuous; transitional forms may be missing not because they never existed but because they were never preserved or never found; the apparent sudden appearance of groups (e.g. the Cambrian "explosion") may partly reflect preservational windows; and estimates of extinction rates and biodiversity through time are minimum values.

\courssub{Extinction}

Every fossil species we describe is, by definition, no longer alive. Extinction is the normal fate of species: more than 99\% of all species that ever existed are extinct.

\begin{definition}[Extinction and Mass Extinction]
\textbf{Extinction} is the complete and irreversible disappearance of a species (or higher taxon), when its last individual dies without descendants. \textbf{Background extinction} is the normal, low, continuous rate of species loss operating at all times. A \textbf{mass extinction} is a geologically brief event during which a large proportion of Earth's species (commonly more than half) disappears across many groups and environments, far above the background rate.
\end{definition}

\textbf{Causes of extinction.}

\begin{itemize}
\item \textbf{Climate change}: global cooling or warming, or drying, shifts habitats faster than species can adapt or migrate.
\item \textbf{Sea-level change}: marine regression destroys shallow-shelf habitats (the richest marine ecosystems); transgression can flood lowlands and reduce oxygen in deep waters.
\item \textbf{Volcanism}: massive, prolonged eruptions (flood basalts) inject ash and gases into the atmosphere, causing acid rain, darkness and climatic disturbance.
\item \textbf{Meteorite (bolide) impact}: a large impact causes instantaneous devastation, global wildfires, dust-induced darkness and collapse of food chains -- the accepted trigger of the \textbf{K--Pg event} (Cretaceous--Palaeogene boundary, about 66 million years ago) that ended the non-avian dinosaurs and the ammonites.
\item \textbf{Competition and biological factors}: newly evolved competitors, predators or diseases can drive less adapted species to extinction.
\item \textbf{Human activity}: habitat destruction, overhunting, pollution and introduced species are causing a present-day extinction crisis, sometimes called the "sixth extinction".
\end{itemize}

\begin{popfigure}
\begin{tikzpicture}[scale=1.0, every node/.style={font=\small}]
% timeline axis (Ma, from 550 to 0)
\draw[thick, ->] (0,0) -- (11.4,0) node[right]{time (Ma)};
\foreach \x/\age in {0/550, 2/440, 4/330, 6/220, 8/110, 10/0}{
  \draw (\x,0.08) -- (\x,-0.08);
  \node[below] at (\x,-0.1) {\age};
}
% extinction events: age Ma -> x = (550-age)/55
% Ordovician 444 -> x=1.93 ; Devonian 372 -> x=3.24 ; Permian 252 -> x=5.42 ; Triassic 201 -> x=6.35 ; KPg 66 -> x=8.8
\draw[very thick, poporange] (1.93,0) -- (1.93,1.0) node[above, align=center]{Ordovician\\ $\sim$444 Ma};
\draw[very thick, poporange] (3.24,0) -- (3.24,1.0) node[above, align=center]{Devonian\\ $\sim$372 Ma};
\draw[very thick, poporange] (5.42,0) -- (5.42,1.6) node[above, align=center]{Permian\\ $\sim$252 Ma\\ (greatest)};
\draw[very thick, poporange] (6.35,0) -- (6.35,1.0) node[above, align=center]{Triassic\\ $\sim$201 Ma};
\draw[very thick, poppurple] (8.8,0) -- (8.8,1.6) node[above, align=center, text=poppurple]{\textbf{K--Pg}\\ $\sim$66 Ma\\ impact};
\fill[poppurple] (8.8,0) circle (0.09);
\end{tikzpicture}
\legende{Timeline of the five major mass extinctions of the Phanerozoic. The Cretaceous--Palaeogene (K--Pg) event, linked to a meteorite impact, is highlighted; the end-Permian event was the most severe of all.}
\end{popfigure}

\begin{exoresolu}[Correlating Two Outcrops and Interpreting a Gap]
\textbf{Statement.} Outcrop X shows, from base to top: bed 1 with the trilobite \emph{Calymene}; bed 2 with the graptolite \emph{Monograptus}; bed 3 with the ammonite \emph{Perisphinctes}. Outcrop Y, 80 km away, shows: bed a with \emph{Calymene}; bed b with \emph{Perisphinctes}; bed c with oyster shells. (i) Correlate the beds. (ii) What happened to the \emph{Monograptus} zone in outcrop Y? (iii) Can the oysters of bed c be used to date it? Justify.
\tcblower
\textbf{Solution.}
(i) Applying the principle of faunal succession, beds containing the same index fossil are of equal age: bed 1 correlates with bed a (\emph{Calymene}), and bed 3 correlates with bed b (\emph{Perisphinctes}).

(ii) In outcrop Y the \emph{Monograptus} zone is absent between beds a and b. Either sediments of that age were never deposited there (non-deposition), or they were deposited and later eroded: in both cases there is a \textbf{stratigraphic gap} (unconformity) between bed a and bed b.

(iii) No. Oysters are \textbf{facies fossils}: they indicate a brackish, shallow coastal environment but range over a very long time. They tell us the \emph{environment} of bed c, not its \emph{age}. To date bed c we would need a true index fossil (short range, wide distribution, abundant, easily identified).
\end{exoresolu}

\begin{aretenir}[Key Points of This Section]
\begin{itemize}
\item Index fossils (wide distribution, short range, abundant, easily identified) date rocks; facies fossils indicate environment. Never confuse the two.
\item The principle of faunal succession allows relative dating and correlation of strata between distant outcrops.
\item Fossils reconstruct palaeoenvironments, palaeogeography and past climates, document evolution (transitional forms, lineages), and serve the oil, coal and limestone industries (e.g. microfossils in the Douala basin).
\item The fossil record is profoundly incomplete: non-deposition, erosion, dissolution, metamorphism, soft bodies and incomplete sampling all create gaps.
\item Extinction is the permanent disappearance of a species; background extinction is continuous, while mass extinctions (five major ones, including the K--Pg impact event) abruptly reset the history of life.
\end{itemize}
\end{aretenir}

\begin{exercice}[Fossils, Gaps and Extinction]
\begin{enumerate}
\item State the four criteria of a good index fossil and explain why ammonites satisfy them.
\item A borehole in the Douala basin yields only planktonic foraminifera at one level and oyster shells at another. Interpret the depositional environment of each level.
\item Give four causes of gaps in the fossil record and explain one consequence of these gaps for the study of evolution.
\item Distinguish background extinction from mass extinction, and name the event and probable cause of the extinction at the K--Pg boundary.
\end{enumerate}
\end{exercice}

\coursec{Classification of Fossils}

In the previous section we saw that fossils are powerful tools for dating rocks and reconstructing past environments, but also that the fossil record is full of gaps and marked by great extinctions. Before we can use fossils to tell the story of life, however, we must be able to \emph{talk about them} in a precise, universal way. A palaeontologist in Yaound\'e, another in London and a third in Tokyo must all understand exactly the same organism when one of them writes a name on a specimen label. This is only possible because fossils, like living organisms, are \cle{classified}. This section explains why classification is needed, how it is done, what difficulties fossils pose, and how the major fossil groups are organised.

\courssub{The Need for Classification}

Imagine walking into a market in Bafoussam where the traders had no names for their goods: ``give me some of that round thing, no, the other round thing, the greenish one\ldots'' Trade would be impossible. Palaeontology faces the same problem on a far larger scale: millions of fossil species have been described, from microscopic foraminifera to giant dinosaurs. Without an agreed system we could not:

\begin{itemize}
\item \textbf{Name} organisms unambiguously, so that scientists everywhere refer to the same fossil;
\item \textbf{Compare} fossils, recognising which specimens are identical, similar or different;
\item \textbf{Organise} the diversity of past life, grouping organisms that share features and are presumed to be related;
\item \textbf{Predict} properties: if a new fossil is placed in the class Cephalopoda, we immediately infer it was a marine, mobile, predatory mollusc.
\end{itemize}

\begin{definition}[Species, Genus and Taxon]
A \cle{species} is the basic unit of classification: a group of organisms that resemble one another closely and, in living forms, can interbreed to produce fertile offspring. A \cle{genus} (plural \emph{genera}) is a group of closely related species sharing a common basic body plan. A \cle{taxon} (plural \emph{taxa}) is any named group at any rank of the classification: a species, a genus, a family, a phylum, etc.
\end{definition}

\courssub{Binomial Nomenclature}

Before the 18th century, organisms were given long descriptive Latin phrases that varied from author to author, so the same animal could carry several different ``names''. The Swedish naturalist \textbf{Carl Linnaeus} (1707--1778) introduced the simple, universal system we still use: \cle{binomial nomenclature}, in which every species receives a name of \emph{two words}.

\begin{methode}[Writing a Scientific Name Correctly]
\begin{enumerate}
\item The name has two parts: the \textbf{genus name} first, then the \textbf{specific epithet}. Example: \emph{Homo sapiens}.
\item The genus name begins with a \textbf{capital letter}; the specific epithet is written entirely in \textbf{small letters}, even if it honours a person or place.
\item Both words are written in \emph{italics} (or underlined when hand-written): \emph{Homo sapiens}, \emph{Acanthoceras rhotomagense}.
\item After first use in a text, the genus may be abbreviated to its initial: \emph{H. sapiens}.
\item Names are Latin or Latinised, so they are identical in every language and every country.
\item When the species cannot be determined, write the genus followed by \emph{sp.} (not italicised): \emph{Ammonites} sp. means ``an undetermined species of the genus \emph{Ammonites}''.
\end{enumerate}
\end{methode}

\begin{attention}[Common Mistakes with Names]
Do not write ``Homo Sapiens'' (the epithet must be lower case), nor ``homo sapiens'' (the genus must be capitalised), and never write a species name as a single word: ``sapiens'' alone is meaningless, because the same epithet can occur in many different genera.
\end{attention}

\courssub{The Taxonomic Hierarchy}

Species and genera are only the two lowest rungs of a ladder of increasingly inclusive groups. Related species are grouped into a \textbf{genus}; related genera into a \textbf{family}; families into an \textbf{order}; orders into a \textbf{class}; classes into a \textbf{phylum}; and phyla into a \textbf{kingdom}. Each level is a taxon, and membership is nested: every member of a family is automatically a member of its order, class, phylum and kingdom. A useful mnemonic, from the largest group downwards, is: \emph{\textbf{K}ing \textbf{P}hilip \textbf{C}ame \textbf{O}ver \textbf{F}or \textbf{G}ood \textbf{S}oup} (Kingdom, Phylum, Class, Order, Family, Genus, Species).

\begin{popfigure}
\begin{center}
\begin{tikzpicture}[
box/.style={draw=popblue, thick, fill=popblueL, rounded corners=2pt, minimum width=4.6cm, minimum height=0.72cm, align=center, font=\small},
ex/.style={font=\small\itshape, text=popdark},
arr/.style={->, thick, poporange}]
\node[box] (k) at (0,6.3) {\textbf{Kingdom}};
\node[box] (p) at (0,5.4) {\textbf{Phylum}};
\node[box] (c) at (0,4.5) {\textbf{Class}};
\node[box] (o) at (0,3.6) {\textbf{Order}};
\node[box] (f) at (0,2.7) {\textbf{Family}};
\node[box] (g) at (0,1.8) {\textbf{Genus}};
\node[box] (s) at (0,0.9) {\textbf{Species}};
\node[ex, anchor=west] at (2.7,6.3) {Animalia};
\node[ex, anchor=west] at (2.7,5.4) {Mollusca};
\node[ex, anchor=west] at (2.7,4.5) {Cephalopoda};
\node[ex, anchor=west] at (2.7,3.6) {Ammonitida};
\node[ex, anchor=west] at (2.7,2.7) {Acanthoceratidae};
\node[ex, anchor=west] at (2.7,1.8) {Acanthoceras};
\node[ex, anchor=west] at (2.7,0.9) {Acanthoceras rhotomagense};
\draw[arr] (s) -- (g);
\draw[arr] (g) -- (f);
\draw[arr] (f) -- (o);
\draw[arr] (o) -- (c);
\draw[arr] (c) -- (p);
\draw[arr] (p) -- (k);
\node[font=\footnotesize, text=popmuted, rotate=90] at (-2.6,3.6) {increasing inclusiveness};
\end{tikzpicture}
\end{center}
\legende{The taxonomic hierarchy, with the classification of a Cretaceous ammonite as a worked example. Each rank contains all the ranks below it.}
\end{popfigure}

\begin{exemplebox}[Worked Example: Classifying an Ammonite]
A coiled, chambered shell with ribbed walls and sutured septa is collected from Cretaceous limestones. Its coiled shell places it among the cephalopod molluscs; its complex sutures place it in the order Ammonitida; its strong ribs and tubercles identify the family Acanthoceratidae and the genus \emph{Acanthoceras}. The full classification is: Kingdom Animalia; Phylum Mollusca; Class Cephalopoda; Order Ammonitida; Family Acanthoceratidae; Genus \emph{Acanthoceras}; species \emph{Acanthoceras rhotomagense}.
\end{exemplebox}

\courssub{Difficulties Specific to Fossil Classification}

Classifying a living snail is easy compared with classifying its fossil shell. Palaeontologists face four special problems:

\begin{itemize}
\item \textbf{Incomplete material.} Fossils are often broken, crushed or represented by a single fragment. A classification may rest on one jawbone or half a shell.
\item \textbf{No soft anatomy or DNA.} Living animals are classified largely on soft parts (gills, nervous system, reproductive organs) and increasingly on genetic sequences. Fossils almost never preserve these, so relationships must be inferred from hard parts alone.
\item \textbf{No breeding test.} The biological definition of a species (interbreeding) cannot be applied to extinct organisms, so fossil species are defined on morphology only (see below).
\item \textbf{Isolated parts of one organism get different names.} A fossil plant may be known from leaves, seeds, pollen and wood found separately; each part is given its own \cle{form genus} (one form genus for the leaf, another for the seed) until a rare complete specimen shows they belonged to a single plant.
\end{itemize}

\begin{propriete}[Morphological Basis of Fossil Classification]
Fossils are classified mainly on the \cle{morphology of their hard parts}: shell shape, ornament, suture patterns, number of rays, plates or segments, tooth structure, and so on. A fossil species is therefore a \emph{morphospecies}: a set of specimens whose hard-part morphology varies no more than is accepted within one living species. (This principle is examinable as a statement, not a proof.)
\end{propriete}

\courssub{The Major Fossil Kingdoms}

All fossil organisms are distributed among the great kingdoms of life:

\begin{center}
\begin{tabularx}{\linewidth}{|l|X|X|}
\hline
\rowcolor{popblueL} \textbf{Kingdom} & \textbf{Characteristics} & \textbf{Important fossils} \\
\hline
Monera & Prokaryotic cells, no nucleus & Bacteria; layered \cle{stromatolites} built by cyanobacteria, among the oldest evidence of life \\
\hline
Protista & Unicellular eukaryotes & \cle{Foraminifera} (calcareous tests) and \cle{radiolaria} (siliceous tests); key microfossils \\
\hline
Plantae & Multicellular, photosynthetic & Algae, ferns, gymnosperms (conifers, cycads), angiosperms (flowering plants) \\
\hline
Animalia & Multicellular, heterotrophic & All invertebrate phyla (Porifera to Echinodermata) and the vertebrates \\
\hline
\end{tabularx}
\end{center}

\courssub{Overview of the Main Fossil Animal Phyla}

Within the Animalia, the phyla that dominate the fossil record are those possessing mineralised hard parts:

\begin{itemize}
\item \textbf{Porifera} (sponges): porous bodies supported by spicules or a skeleton of silica, calcite or spongin; sessile marine filter feeders.
\item \textbf{Cnidaria}: includes the \textbf{corals}, with calcareous cups (solitary or colonial), radial symmetry and septa; major reef builders.
\item \textbf{Mollusca}: unsegmented soft body with a muscular foot and usually a shell; classes include \textbf{Gastropoda} (snails), \textbf{Bivalvia} (clams, oysters), \textbf{Cephalopoda} (ammonites, nautiloids, belemnites).
\item \textbf{Brachiopoda} (lamp shells): two-valved shell, but the valves are \emph{dorsal and ventral} (unequal), each valve being bilaterally symmetrical; fixed to the sea floor by a fleshy stalk (pedicle).
\item \textbf{Arthropoda}: jointed limbs and a segmented exoskeleton, moulted during growth; \textbf{trilobites} are the emblematic fossil group.
\item \textbf{Echinodermata}: five-fold (pentaradial) symmetry and a skeleton of calcite plates; crinoids (sea lilies), echinoids (sea urchins), starfish.
\item \textbf{Protochordates and Vertebrates}: chordates possess a notochord; vertebrates add a mineralised skeleton --- fish, amphibians, reptiles, birds and mammals, known from bones, teeth and footprints.
\end{itemize}

\begin{attention}[Brachiopods are NOT Bivalves]
The classic GCE trap: a brachiopod looks like a bivalve mollusc because both have two shells. But in a \textbf{bivalve} the two valves are \emph{left and right}, usually mirror images of each other, while the two sides of one valve are asymmetrical. In a \textbf{brachiopod} the two valves are \emph{dorsal and ventral} and differ from each other, but \emph{each single valve is symmetrical} about its own mid-line. They belong to \textbf{different phyla} (Brachiopoda vs Mollusca). Always check the plane of symmetry before naming the phylum.
\end{attention}

\begin{popfigure}
\begin{center}
\begin{tikzpicture}
\begin{axis}[
ybar, bar width=16pt,
width=0.95\linewidth, height=6.2cm,
ylabel={Relative abundance (qualitative scale)},
symbolic x coords={Porifera, Cnidaria, Brachiopoda, Mollusca, Arthropoda, Echinodermata, Vertebrata},
xtick=data, x tick label style={rotate=30, anchor=east, font=\footnotesize},
ymin=0, ymax=10, ytick={0,2,4,6,8,10},
nodes near coords, nodes near coords style={font=\footnotesize},
axis lines=left, ymajorgrids, grid style={popline},
every axis plot/.append style={fill=popblue, draw=popdark}]
\addplot coordinates {(Porifera,3) (Cnidaria,5) (Brachiopoda,6) (Mollusca,9) (Arthropoda,7) (Echinodermata,6) (Vertebrata,2)};
\end{axis}
\end{tikzpicture}
\end{center}
\legende{Qualitative comparison of the abundance of the major fossil phyla in the record. Molluscs dominate because their shells fossilise easily; vertebrates are comparatively rare.}
\end{popfigure}

\courssub{The Morphological Species Concept Applied to Fossils}

Because we cannot observe breeding in extinct organisms, palaeontologists apply the \cle{morphological species concept}: a fossil species is a population of specimens sharing a distinctive combination of morphological characters, whose variation is comparable to that seen within a single living species. In practice this means measuring and comparing many specimens: if the spread of shapes in a fossil assemblage is no greater than the spread within a known living species, the assemblage is treated as one species; if it splits into clearly separated morphological groups, several species are recognised.

\begin{exoresolu}[Defining Fossil Species from a Collection]
\textbf{Statement.} A bed of Cenomanian limestone yields 60 coiled cephalopod shells. They fall into two sets: set A (45 shells) has strong regular ribs with ventral tubercles; set B (15 shells) is smooth with only faint growth lines. A student claims all 60 belong to one species because they come from one bed. Discuss and conclude.
\tcblower
\textbf{Solution.}
\begin{enumerate}
\item Occurrence in one bed shows only that the animals lived at the same time and place; it says nothing about species identity.
\item Apply the morphological species concept: compare the hard-part characters. Ribbing and tuberculation are exactly the kind of stable, genetically controlled features used to separate ammonite species.
\item The difference between set A (strong ribs, tubercles) and set B (smooth) is far greater than the variation observed \emph{within} each set, and greater than normal variation within living cephalopod species.
\item Conclusion: the collection contains \textbf{two distinct morphospecies} (and possibly two genera), not one. The student's mistake was to use stratigraphic occurrence instead of morphology as the criterion.
\end{enumerate}
\end{exoresolu}

\begin{exercice}[Classification Practice]
\begin{enumerate}
\item Write out the full classification (kingdom to class) of a fossil coral, a trilobite and a crinoid.
\item Correct these names and justify each correction: ``homonculus Erectus'', ``Ammonites'', used alone as a species name.
\item A two-valved shell has a plane of symmetry passing \emph{between} the two valves. To which phylum does it belong? What if the plane of symmetry passes \emph{through} each valve?
\end{enumerate}
\end{exercice}

\begin{corrige}[Exercise 1]
\begin{enumerate}
\item Coral: Animalia, Cnidaria, Anthozoa. Trilobite: Animalia, Arthropoda, Trilobita. Crinoid: Animalia, Echinodermata, Crinoidea.
\item \emph{Homo erectus}: genus capitalised, epithet in lower case, both italicised. \emph{Ammonites} alone is only a genus name; a species requires a binomial, e.g. \emph{Ammonites} sp. if the species is undetermined.
\item Symmetry between the valves (left and right valves equal): phylum Mollusca, class Bivalvia. Symmetry through each valve (dorsal and ventral valves unequal): phylum Brachiopoda.
\end{enumerate}
\end{corrige}

\begin{aretenir}[Key Points]
\begin{itemize}
\item Classification allows us to \textbf{name, compare and organise} fossil diversity using the Linnaean binomial system: \emph{Genus species}, genus capitalised, both italicised.
\item The hierarchy runs: species $\rightarrow$ genus $\rightarrow$ family $\rightarrow$ order $\rightarrow$ class $\rightarrow$ phylum $\rightarrow$ kingdom.
\item Fossils are classified on the \textbf{morphology of hard parts} (morphospecies), because soft anatomy, DNA and breeding behaviour are lost.
\item Special difficulties: incomplete material, and \textbf{form genera} for isolated plant parts.
\item The main fossil kingdoms are Monera, Protista, Plantae and Animalia; the dominant fossil phyla are Mollusca, Arthropoda, Brachiopoda, Echinodermata, Cnidaria and Porifera.
\item Never confuse brachiopods with bivalves: check the plane of symmetry of the shell.
\end{itemize}
\end{aretenir}

\begin{savaistu}[Fossil Classification in Cameroon]
Marine Cretaceous sediments of the Douala and Rio del Rey basins have yielded ammonites, bivalves, gastropods and foraminifera that palaeontologists classify with the same Linnaean system used worldwide. Identifying these fossils to genus and species level is what allows Cameroonian strata to be correlated precisely with rocks of the same age in Nigeria, Brazil and Europe --- a direct application of everything learned in this section.
\end{savaistu}

\coursec{Description of Commonly Preserved Fossils I: Microfossils, Sponges, Corals and Graptolites}

In the previous section we learned how fossils are \emph{classified} --- placed within the great tree of life according to their morphology and biological affinities. Classification, however, is only a filing system. At the GCE Advanced Level, the examiner expects you to go further: to pick up a fossil, \emph{describe} it in a disciplined way, \emph{identify} its group, and \emph{state} its geological importance. In this section and the next, we therefore tour the groups of organisms most commonly preserved as fossils, applying one single method of description. We begin with the smallest and the simplest: the microfossils (foraminifera, radiolaria, diatoms), the sponges, the corals and the graptolites.

\courssub{A Standard Method for Describing a Fossil}

Before describing any group, we need a repeatable grid. A good palaeontological description is like a good police report: it answers the same questions in the same order, so that nothing is forgotten.

\begin{methode}[The standard fossil description grid]
For every fossil or fossil group, describe the following points in order:
\begin{enumerate}
\item \textbf{Name} of the fossil or group (common name and, where possible, a typical genus).
\item \textbf{Phylum} (and class or order where relevant).
\item \textbf{Morphology}: shape, symmetry, main structures (use a labelled sketch).
\item \textbf{Composition} of the hard parts (calcite, aragonite, silica, chitin, calcium phosphate).
\item \textbf{Size} (typical order of magnitude).
\item \textbf{Habitat and mode of life} (marine/freshwater/terrestrial; benthic, planktonic, sessile, mobile).
\item \textbf{Geological range} (first to last appearance).
\item \textbf{Use} in geology (index fossil, palaeoenvironmental indicator, economic value).
\end{enumerate}
\end{methode}

\begin{attention}[Describing is not naming]
A frequent examination weakness is to write only ``this is a coral''. A description must give \emph{observable} characters (shape, symmetry, structures, composition) \emph{before} the name. Train yourself to describe first, identify second.
\end{attention}

\courssub{Microfossils I: The Foraminifera}

Imagine drilling an oil well offshore near Kribi or Limbe. The rock cuttings brought up from thousands of metres down are washed and examined under a microscope: they are full of tiny shells, often smaller than a grain of sand. These are the tests of \cle{foraminifera}, and they tell the petroleum geologist the age and environment of each layer drilled.

\begin{definition}[Test]
The \cle{test} is the shell secreted by a foraminiferan. It is made of one or several \textbf{chambers} arranged in a line, a spiral or a coil, usually of \textbf{calcium carbonate} (low-magnesium calcite), sometimes of agglutinated sand grains or organic material. The chambers communicate through openings called \textbf{apertures} and \textbf{foramina} --- hence the name \emph{foraminifera}, ``bearers of little holes''.
\end{definition}

Foraminifera are \textbf{unicellular protists} (single-celled eukaryotes), yet they build remarkably elaborate shells. Two classic examples must be known:

\begin{itemize}
\item \emph{Nummulites}: a large foraminifer (up to several centimetres across), shaped like a lens or a coin (Latin \emph{nummus}, coin), with chambers coiled in a flat \textbf{planispiral}. Very abundant in \textbf{Eocene} limestones; the ``nummulitic limestones'' were used to build the pyramids of Egypt.
\item \emph{Globigerina}: a small \textbf{planktonic} (floating) foraminifer with a few globular chambers arranged in a low \textbf{trochospiral}. Its tests rain down onto the deep ocean floor today, forming \textbf{globigerina ooze}, and its fossil relatives are used to date Cretaceous--Cenozoic marine sediments.
\end{itemize}

\begin{popfigure}
\begin{center}
\begin{tikzpicture}[scale=1.05]
% Nummulites test in section: lens shape with spiral chambers
\draw[thick, popblue, fill=popblueL] (0,0) ellipse (3.2 and 1.15);
\draw[popdark, thick] (0,0) ellipse (2.6 and 0.85);
\draw[popdark, thick] (0,0) ellipse (1.9 and 0.6);
\draw[popdark, thick] (0,0) ellipse (1.25 and 0.38);
\draw[popdark, thick] (0,0) ellipse (0.65 and 0.2);
\fill[poporange] (0,0) circle (0.06);
% chamber partitions (septa)
\foreach \a in {30,60,120,150,210,240,300,330}{
  \draw[popmuted] (\a:0.65) -- (\a:1.25);
  \draw[popmuted] (\a:1.9) -- (\a:2.6);
}
% labels
\draw[->, popink] (3.2,0.9) -- (4.1,1.5) node[right, align=left]{\footnotesize outer wall (calcite)};
\draw[->, popink] (1.9,0.55) -- (2.6,1.7) node[right, align=left]{\footnotesize coiled chambers};
\draw[->, popink] (0.06,-0.06) -- (-1.2,-1.4) node[below, align=center]{\footnotesize initial chamber\\ \footnotesize (proloculus)};
\draw[->, popink] (-2.6,0.6) -- (-3.6,1.4) node[left, align=right]{\footnotesize septa between\\ \footnotesize chambers};
\end{tikzpicture}
\end{center}
\legende{Median section through the coin-shaped test of \emph{Nummulites} (Eocene), showing the planispiral arrangement of chambers around the initial chamber (proloculus).}
\end{popfigure}

Applying our grid to the foraminifera:

\begin{tabularx}{\linewidth}{|l|X|}
\hline
\rowcolor{popblueL} \textbf{Criterion} & \textbf{Foraminifera} \\
\hline
Phylum & Foraminifera (unicellular eukaryotes; often treated as class within Retaria) \\
\hline
Morphology & Chambered test; chambers in rows, coils or spirals (planispiral, trochospiral, milioline); apertures and foramina \\
\hline
Composition & Calcite (most), agglutinated particles (some), rarely aragonite or organic (tectin) \\
\hline
Size & Mostly \SI{0.1}{mm}--\SI{1}{mm}; \emph{Nummulites} up to several cm \\
\hline
Habitat & Exclusively or almost exclusively \textbf{marine}; benthic (bottom-dwelling) or planktonic (floating) \\
\hline
Range & \textbf{Cambrian to Recent} (still living today) \\
\hline
Use & Index fossils; palaeoenvironmental and palaeodepth indicators; dating of sediments in \textbf{oil exploration} \\
\hline
\end{tabularx}

\begin{propriete}[Foraminifera in oil exploration]
Because foraminifera are (i) extremely abundant, (ii) small enough to be recovered intact from drill cuttings, (iii) rapidly evolving (so each species has a short geological range) and (iv) sensitive to water depth, temperature and salinity, they are \textbf{the most useful microfossils for dating marine sediments and reconstructing depositional environments in oil wells}. (Proof not examinable; the justification above is.)
\end{propriete}

\begin{savaistu}[Foraminifera and the Cameroonian petroleum industry]
The Cretaceous--Cenozoic sedimentary basins of the Rio del Rey and Douala basins, which host Cameroon's oil fields, are dated and correlated largely by means of foraminifera and other microfossils recovered from boreholes. Micropalaeontology is thus a direct tool of the national petroleum industry.
\end{savaistu}

\courssub{Microfossils II: Radiolaria and Diatoms}

Two other groups of microfossils must be recognised, more briefly:

\begin{itemize}
\item \textbf{Radiolaria}: unicellular marine protists (phylum Retaria, class Radiolaria), planktonic, with an intricate skeleton of \textbf{silica} ($\mathrm{SiO_2}$) made of spines and lattice-like spheres. Range: Cambrian--Recent. Their skeletons accumulate as \emph{radiolarian ooze} in deep oceans and form the rock \textbf{radiolarite} (a siliceous sedimentary rock). They indicate deep marine environments.
\item \textbf{Diatoms}: unicellular \textbf{algae} (phylum Bacillariophyta) with a two-valved siliceous shell (frustule), often beautifully ornamented. They live in marine \emph{and} fresh waters; range essentially Cretaceous--Recent. Accumulations of frustules form \textbf{diatomite}, used industrially as a filter and abrasive.
\end{itemize}

\begin{attention}[Calcareous versus siliceous]
Do not confuse the composition: foraminifera have \textbf{calcareous} tests; radiolaria and diatoms have \textbf{siliceous} skeletons. Siliceous skeletons resist dissolution in deep, cold ocean water where calcite dissolves --- which is why radiolarian oozes occur at greater depths than globigerina oozes (below the \emph{calcite compensation depth}).
\end{attention}

\courssub{Phylum Porifera: The Sponges}

Sponges represent the \textbf{simplest grade of organisation} among multicellular animals: no true tissues, no organs, no nervous system. The body is a \textbf{porous sac or vase} pierced by innumerable pores (\emph{ostia}) through which water enters, and one or more larger openings (\emph{oscula}) through which it leaves. The sponge feeds by filtering microscopic particles from this water current.

The soft body is supported by a skeleton of microscopic needles called \cle{spicules}, made of \textbf{silica} (class Demospongiae, Hexactinellida) or \textbf{calcite} (class Calcarea). In some groups the skeleton is of the organic protein \textbf{spongin} (rarely fossilised). It is mainly the spicules --- or the fused spicule network --- that fossilise.

\begin{tabularx}{\linewidth}{|l|X|}
\hline
\rowcolor{popblueL} \textbf{Criterion} & \textbf{Porifera (sponges)} \\
\hline
Morphology & Porous vase-, cup- or ball-shaped body; ostia and osculum; skeleton of spicules (monaxons, triaxons, tetraxons, polyaxons) \\
\hline
Composition & Spicules of silica or calcite; sometimes spongin (rarely fossilised) \\
\hline
Size & A few mm to more than \SI{1}{m} \\
\hline
Habitat & Mostly marine, sessile (fixed to the sea floor), filter feeders; a few freshwater forms (e.g. \emph{Spongilla}) \\
\hline
Range & \textbf{Cambrian (even Precambrian) to Recent} \\
\hline
Use & Locally rock-builders (reef sponges, e.g. stromatoporoids); indicators of marine benthic environments \\
\hline
\end{tabularx}

Because the spicules are often scattered after death, complete fossil sponges are less common than their abundance in life would suggest --- a good illustration of \emph{preservation potential} studied earlier.

\courssub{Phylum Cnidaria: The Corals}

Corals are \textbf{cnidarian animals} --- relatives of jellyfish and sea anemones --- that secrete an external \textbf{calcareous skeleton} (a corallum). The living animal is a soft polyp, like a miniature sea anemone, sitting in a cup (\emph{calyx} or \emph{calice}) at the top of the skeleton. Inside the cup, vertical radial plates called \cle{septa} support the polyp; horizontal plates (\emph{tabulae}) and curved plates (\emph{dissepiments}) may also be present.

\begin{attention}[Corals are animals, not plants!]
The most frequent examination mistake is to describe corals as plants (or as ``rocks''). Corals are \textbf{colonial cnidarian animals}: each branch of a coral colony is covered with tiny polyps, each polyp an individual animal. The ``stony'' part is only the skeleton they secrete.
\end{attention}

Two great architectural types exist:

\begin{itemize}
\item \textbf{Solitary corals}: a single polyp secretes a horn-shaped (conical) skeleton --- the \emph{horn corals}. Example: the rugose genus \emph{Zaphrentis}.
\item \textbf{Colonial corals}: many polyps bud from one another and build a shared skeleton, forming massive, branching or chain-like colonies. These are the great \textbf{reef builders}.
\end{itemize}

Three orders of corals dominate the fossil record:

\begin{tabularx}{\linewidth}{|l|X|X|}
\hline
\rowcolor{popblueL} \textbf{Order} & \textbf{Characters} & \textbf{Range} \\
\hline
Tabulate corals & Exclusively colonial; corallites with strong horizontal tabulae; septa absent or very reduced & Ordovician--Permian (extinct) \\
\hline
Rugose (tetracorals) & Solitary (horn corals) or colonial; well-developed septa in multiples of four (tetramerous symmetry) & Ordovician--Permian (extinct) \\
\hline
Scleractinian (hexacorals) & Solitary or colonial; septa in multiples of six (hexamerous symmetry); the modern reef builders & Triassic--Recent \\
\hline
\end{tabularx}

\begin{popfigure}
\begin{center}
\begin{tikzpicture}[scale=0.95]
% Horn coral
\begin{scope}[xshift=-3.4cm]
\draw[thick, poppurple, fill=poppurple!12] (0,0) .. controls (-0.9,-0.5) and (-1.1,-2.2) .. (-0.55,-3.4)
   -- (0.35,-3.4) .. controls (0.9,-2.2) and (0.9,-0.5) .. (0,0);
\draw[thick, poppurple] (0,0) ellipse (0.95 and 0.28);
% septa inside calyx
\foreach \a in {20,50,80,110,140,160}{
  \draw[popdark] (\a:0.95 and 0.28) -- (0,0);
}
\draw[popdark] (200:0.95 and 0.28) -- (0,0);
\draw[popdark] (340:0.95 and 0.28) -- (0,0);
% growth lines
\draw[popmuted] (-0.72,-1.2) .. controls (-0.2,-1.05) and (0.3,-1.05) .. (0.62,-1.25);
\draw[popmuted] (-0.62,-2.2) .. controls (-0.15,-2.05) and (0.25,-2.05) .. (0.5,-2.25);
\draw[->, popink] (0.95,0.14) -- (2.0,0.9) node[right, align=left]{\footnotesize calyx with septa};
\draw[->, popink] (0.45,-2.6) -- (2.0,-2.2) node[right, align=left]{\footnotesize conical theca\\ \footnotesize (horn shape)};
\node[align=center] at (-0.1,-4.1) {\footnotesize \textbf{Solitary rugose coral}\\ \footnotesize (horn coral)};
\end{scope}
% Tabulate coral colony
\begin{scope}[xshift=3.6cm]
\draw[thick, popgreen, fill=popgreenL] (-1.9,-3.2) -- (-1.9,-0.4)
  .. controls (-1.9,0.5) and (-1.2,0.9) .. (0,0.9)
  .. controls (1.2,0.9) and (1.9,0.5) .. (1.9,-0.4) -- (1.9,-3.2) -- cycle;
% corallite tubes
\foreach \x in {-1.5,-1.0,-0.5,0,0.5,1.0,1.5}{
  \draw[popdark] (\x,-3.2) -- (\x,0.35);
}
% tabulae
\foreach \y in {-2.7,-2.1,-1.5,-0.9,-0.3}{
  \draw[popmuted] (-1.85,\y) -- (1.85,\y);
}
\draw[->, popink] (1.0,-0.3) -- (2.5,-0.1) node[right, align=left]{\footnotesize tabulae};
\draw[->, popink] (-1.0,-1.6) -- (-2.6,-1.3) node[left, align=right]{\footnotesize corallite tubes};
\node[align=center] at (0,-4.1) {\footnotesize \textbf{Colonial tabulate coral}\\ \footnotesize (massive colony)};
\end{scope}
\end{tikzpicture}
\end{center}
\legende{Left: a solitary rugose (horn) coral with radial septa in the calyx (tetramerous arrangement). Right: a massive colony of a tabulate coral, showing parallel corallite tubes crossed by horizontal tabulae.}
\end{popfigure}

\begin{propriete}[Rugose and tabulate corals as Palaeozoic index fossils]
Rugose and tabulate corals evolved rapidly, are abundant and widespread in marine limestones, and \textbf{both orders died out at the end of the Permian}. They are therefore excellent \textbf{index fossils for Palaeozoic (especially Ordovician--Permian) marine rocks}. Finding a horn coral immediately tells the field geologist that the rock is Palaeozoic, not Mesozoic or Cenozoic.
\end{propriete}

Corals are also \textbf{palaeoenvironmental indicators}: reef-building scleractinians (like their modern descendants in warm seas) require warm, clear, shallow, well-lit marine water, so a fossil reef signals an ancient tropical shallow sea.

\courssub{Graptolites: The Saw-Blade Fossils of Marine Shales}

Suppose you split a slab of dark Ordovician shale and find a delicate marking that looks like a tiny saw blade or a pencil scribble. You have found a \cle{graptolite} --- one of the most important index fossils of the early Palaeozoic.

Graptolites are \textbf{extinct colonial animals} belonging to the \textbf{hemichordates} (relatives of the ancestors of chordates). Each colony (\emph{rhabdosome}) consisted of one or several branches called \cle{stipes}. Along each stipe, tiny cup-shaped tubes, the \cle{thecae}, were arranged in one or two rows; each theca housed an individual zooid. The skeleton was made of a resistant organic (chitin-like) material called \emph{periderm}, which is why graptolites are commonly preserved as \textbf{carbonised films} in fine-grained marine shales.

\begin{popfigure}
\begin{center}
\begin{tikzpicture}[scale=1.0]
% stipe
\draw[very thick, popink] (0,0) -- (0,4.6);
% thecae: saw teeth on both sides (biserial)
\foreach \y in {0.25,0.65,1.05,1.45,1.85,2.25,2.65,3.05,3.45,3.85,4.25}{
  \draw[thick, poporange] (0,\y) -- (0.55,\y+0.28) -- (0.55,\y+0.02);
  \draw[thick, poporange] (0,\y) -- (-0.55,\y+0.28) -- (-0.55,\y+0.02);
}
% sicula
\draw[thick, popdark] (0,0) -- (0.18,-0.5) -- (-0.05,-0.45) -- cycle;
\draw[->, popink] (0.55,3.6) -- (1.7,4.2) node[right, align=left]{\footnotesize thecae (cups of zooids)};
\draw[->, popink] (0,2.3) -- (1.7,2.3) node[right, align=left]{\footnotesize stipe (branch)};
\draw[->, popink] (0.1,-0.45) -- (1.7,-0.7) node[right, align=left]{\footnotesize sicula (initial cone)};
\node[align=center] at (0,-1.3) {\footnotesize \textbf{A biserial graptolite stipe}\\ \footnotesize (saw-blade outline)};
\end{tikzpicture}
\end{center}
\legende{Sketch of a graptolite stipe bearing two rows of thecae (biserial), growing upwards from the initial sicula. Colonies could bear one (uniserial), two (biserial), four or more stipes.}
\end{popfigure}

\begin{tabularx}{\linewidth}{|l|X|}
\hline
\rowcolor{popblueL} \textbf{Criterion} & \textbf{Graptolites} \\
\hline
Phylum & Hemichordata (class Graptolithina) \\
\hline
Morphology & Colonial; one or more stipes bearing rows of thecae; saw-blade or tuning-fork outlines; uniserial, biserial, or multiserial arrangements \\
\hline
Composition & Organic, chitin-like skeleton (periderm) preserved as carbon films \\
\hline
Size & Colonies from a few mm to about \SI{1}{m} (mostly a few cm) \\
\hline
Habitat & Marine; many planktonic (floating), hence worldwide distribution; some benthic (dendroids) \\
\hline
Range & Cambrian--Carboniferous; \textbf{most useful in the Ordovician and Silurian} \\
\hline
Use & Outstanding \textbf{index fossils} for dating and correlating Ordovician--Silurian marine shales \\
\hline
\end{tabularx}

\begin{aretenir}[Why graptolites are superb index fossils]
Graptolites combine the four qualities of an ideal index fossil: (i) \textbf{rapid evolution} (short-lived species), (ii) \textbf{wide geographic distribution} (planktonic, spread by ocean currents), (iii) \textbf{abundance} in marine shales, and (iv) \textbf{easy identification}. They allow correlation of Ordovician--Silurian strata between continents.
\end{aretenir}

\begin{exoresolu}[Describing an unknown fossil]
\textbf{Statement.} A student finds, in a dark marine shale, a fossil \SI{3}{cm} long resembling a tiny double-edged saw, preserved as a black film. Describe it fully using the standard grid and identify it.
\tcblower
\textbf{Solution.}
\begin{itemize}
\item \textbf{Morphology}: an elongated colony (rhabdosome) made of a stipe bearing two rows of small cup-like thecae, giving the saw-blade outline; preserved as a carbonised film.
\item \textbf{Composition}: originally organic (chitin-like periderm), now a carbon film --- a carbonization mode of preservation.
\item \textbf{Size}: \SI{3}{cm}, typical of graptolite colonies.
\item \textbf{Habitat}: marine; the planktonic habit explains its occurrence in fine offshore shales.
\item \textbf{Identification}: a \textbf{graptolite} (phylum Hemichordata).
\item \textbf{Range and use}: if the shale is Ordovician--Silurian, the graptolite species can date it precisely, because graptolites are index fossils of early Palaeozoic marine shales.
\end{itemize}
\end{exoresolu}

\begin{exercice}[Applying the description grid]
For each of the following, give the phylum, composition of hard parts, habitat, geological range and main geological use: (a) \emph{Globigerina}; (b) a horn coral; (c) a tabulate coral; (d) a sponge.
\end{exercice}

\begin{corrige}[Exercise 1]
(a) \emph{Globigerina}: protist (foraminifer); calcareous chambered test; marine planktonic; Cretaceous--Recent; dating of marine sediments, especially in oil wells. (b) Horn coral: cnidarian, order Rugosa; calcite skeleton with septa in multiples of four; solitary benthic marine; Ordovician--Permian; Palaeozoic index fossil. (c) Tabulate coral: cnidarian, exclusively colonial; calcite corallites with tabulae; shallow marine, reef-building; Ordovician--Permian; Palaeozoic index fossil and reef indicator. (d) Sponge: poriferan; spicules of silica or calcite; sessile marine filter feeder; Cambrian--Recent; indicator of marine benthic environments, locally reef-builder.
\end{corrige}

\begin{aretenir}[Key points of the section]
\begin{itemize}
\item Describe every fossil with the grid: name, phylum, morphology, composition, size, habitat, range, use.
\item Foraminifera: unicellular protists with chambered \textbf{calcareous tests} (\emph{Nummulites}, \emph{Globigerina}); Cambrian--Recent; marine; the key microfossils of oil-well dating.
\item Radiolaria and diatoms: \textbf{siliceous} microfossils; radiolaria indicate deep marine settings; diatoms form diatomite.
\item Sponges: simplest multicellular animals; porous bodies with \textbf{spicules} of silica or calcite; Cambrian--Recent.
\item Corals: colonial or solitary \textbf{cnidarian animals} (never plants!) with calcareous skeletons and septa; rugose (tetramerous) and tabulate corals are Palaeozoic index fossils; scleractinians (hexamerous) build modern reefs.
\item Graptolites: extinct colonial hemichordates with saw-blade stipes bearing thecae; the supreme index fossils of Ordovician--Silurian marine shales.
\end{itemize}
\end{aretenir}

Having mastered the microfossils and the simplest invertebrate groups, we are ready, in the next section, to describe the larger and more complex invertebrates --- molluscs, brachiopods, arthropods and echinoderms --- as well as vertebrates and plants, which complete our tour of the commonly preserved fossils.

\coursec{Description of Commonly Preserved Fossils II: Molluscs, Brachiopods, Arthropods, Echinoderms, Vertebrates and Plants (Integration)}

In the previous section we described the smaller or more simply built fossil groups: microfossils, sponges, corals and graptolites. We now turn to the larger and more complex invertebrates --- the \cle{molluscs}, \cle{brachiopods}, \cle{arthropods} and \cle{echinoderms} --- and then to \cle{vertebrates} and \cle{fossil plants}. These groups supply many of the classic GCE index fossils, especially the ammonites and trilobites, so their morphology must be mastered in detail. For each group we shall follow the same plan: the \textbf{hard parts} that can be preserved, the \textbf{diagnostic morphology}, the \textbf{mode of life}, and the \textbf{stratigraphic range and use}.

\courssub{Phylum Mollusca}

Molluscs are soft-bodied, unsegmented animals usually protected by a calcareous shell secreted by a fold of the body wall called the \cle{mantle}. Because the shell is made of calcium carbonate, molluscs have a high preservation potential and are among the most abundant macrofossils. Three classes are of great palaeontological importance: \cle{Bivalvia}, \cle{Gastropoda} and \cle{Cephalopoda}.

\textbf{Class Bivalvia (lamellibranchs or pelecypods).} Bivalves such as oysters, clams, cockles and mussels possess \textbf{two hinged valves} joined along a dorsal hinge line furnished with interlocking \textbf{teeth and sockets}, and closed by one or two \textbf{adductor muscles} whose attachment scars may be preserved on the inside of the shell. The two valves are usually \textbf{equal in size and mirror images of each other}: the plane of symmetry of the shell lies \emph{between} the valves. The oldest part of each valve is the beak-like \cle{umbo}, situated near the hinge; growth is recorded by \textbf{concentric growth lines} around the umbo. The shell is made of calcium carbonate (calcite or aragonite). Bivalves are mostly marine and benthic (burrowers, attached forms or free livers), though some live in fresh water; they range from the \textbf{Ordovician to the Recent} and are excellent \textbf{environmental indicators}.

\begin{popfigure}
\begin{tikzpicture}[scale=1.0]
% Bivalve shell (single valve, exterior)
\draw[thick, popblue] (0,0) .. controls (-2.2,0.4) and (-2.4,2.2) .. (-0.6,2.6)
 .. controls (0.8,2.9) and (2.2,2.0) .. (2.0,0.6) .. controls (1.9,0.1) and (1.0,-0.1) .. (0,0);
\draw[popblue] (0,0) .. controls (-1.7,0.5) and (-1.9,2.0) .. (-0.5,2.3);
\draw[popblue] (0,0) .. controls (-1.1,0.6) and (-1.2,1.7) .. (-0.3,1.9);
\draw[popblue] (0,0) .. controls (-0.5,0.7) and (-0.5,1.3) .. (-0.1,1.5);
\fill[popblue] (0,0) circle (0.06);
\node[popink, align=center] at (0.6,1.4) {\small growth\\ \small lines};
\node[popink] at (0,-0.5) {\small umbo (hinge)};
\draw[->, popmuted] (0,-0.35) -- (0,-0.05);
\legende{Exterior of a bivalve valve: note the umbo near the hinge and the concentric growth lines.}
\end{tikzpicture}
\end{popfigure}

\textbf{Class Gastropoda.} Gastropods (snails and slugs) possess a \textbf{single, usually spirally coiled shell} (univalve), coiled in one plane or, more commonly, in a helical spiral around a central axis (the \cle{columella}). The successive turns of the shell are called \textbf{whorls}; the animal's body occupies the last and largest whorl (the \cle{body whorl}), which opens at the \cle{aperture}; the stacked earlier whorls form the \cle{spire}. Gastropods live in marine, freshwater and terrestrial environments and range from the \textbf{Cambrian to the Recent}.

\begin{popfigure}
\begin{tikzpicture}[scale=1.0]
% Gastropod shell
\draw[thick, poppurple] (0,0) .. controls (1.6,0.2) and (1.8,1.2) .. (1.0,1.8)
 .. controls (0.4,2.2) and (-0.4,2.0) .. (-0.6,1.4)
 .. controls (-0.8,0.9) and (-0.3,0.5) .. (0.1,0.7)
 .. controls (0.4,0.85) and (0.35,1.2) .. (0.1,1.25);
\draw[thick, poppurple] (0,0) .. controls (0.4,-0.5) and (0.2,-1.1) .. (-0.3,-1.3);
\draw[thick, poppurple] (-0.3,-1.3) .. controls (-0.9,-1.1) and (-1.0,-0.4) .. (-0.6,0.1);
\node[popink] at (1.3,0.4) {\small body whorl};
\node[popink] at (0.1,1.6) {\small spire};
\node[popink] at (-0.6,-1.6) {\small aperture};
\draw[->, popmuted] (-0.6,-1.45) -- (-0.45,-1.05);
\legende{A gastropod shell: single coiled valve with spire, body whorl and aperture.}
\end{tikzpicture}
\end{popfigure}

\textbf{Class Cephalopoda.} Cephalopods are exclusively marine, swimming or floating predators. Their shell (when present) is divided into \textbf{chambers} separated by partitions called \cle{septa}; the animal lives in the last (body) chamber, and a fleshy tube, the \cle{siphuncle}, passes through all the chambers and regulates buoyancy. Three groups are essential for the GCE.

\begin{itemize}
\item \textbf{Ammonites}: coiled (planispiral) chambered shells in which the line where each septum meets the shell wall forms a \cle{suture line}. They range from the \textbf{Devonian to the end of the Cretaceous} and are the outstanding \textbf{Mesozoic index fossils}.
\item \textbf{Belemnites}: squid-like cephalopods whose commonest fossil part is the solid, bullet-shaped internal \cle{guard} made of calcite; Mesozoic (especially Jurassic--Cretaceous).
\item \textbf{Nautiloids}: chambered shells, straight (orthoconic) or coiled, with \emph{simple} sutures; Cambrian--Recent (the modern \emph{Nautilus} is a living representative).
\end{itemize}

\begin{definition}[Suture line]
The \cle{suture line} of an ammonite is the line of junction between a septum and the inner wall of the shell, visible on the surface of the fossil when the outer shell wall is removed. Its pattern is of great \textbf{taxonomic value}: nautiloids have simple, smooth sutures; goniatites have simple angular sutures; ceratites have smooth saddles and frilled lobes; true ammonites have highly complex, finely frilled (ammonitic) sutures. Suture complexity increases through time and helps identify species and date rocks.
\end{definition}

\begin{popfigure}
\begin{tikzpicture}[scale=1.1]
% Ammonite spiral
\draw[thick, poporange] (0,0) circle (0.18);
\draw[thick, poporange] (0.05,0) circle (0.38);
\draw[thick, poporange] (0.12,0) circle (0.68);
\draw[thick, poporange] (0.22,0) circle (1.05);
\draw[thick, poporange] (0.35,0) circle (1.5);
% suture lines on outer whorl
\foreach \a in {20,60,100,140,200,250,300,340}{
 \draw[popdark, thick] (\a:1.42) .. controls (\a+6:1.55) and (\a-6:1.65) .. (\a:1.78);
}
\node[popink] at (2.4,1.3) {\small suture lines};
\draw[->, popmuted] (2.1,1.25) -- (1.35,1.15);
\node[popink] at (-1.9,-1.3) {\small coiled chambered shell};
\legende{Ammonite: planispiral chambered shell with frilled suture lines on the whorls.}
\end{tikzpicture}
\end{popfigure}

\begin{popfigure}
\begin{tikzpicture}[scale=1.0]
% Belemnite guard
\draw[thick, popgreen] (0,0) .. controls (0.35,0.4) and (0.4,2.4) .. (0.18,3.4)
 -- (-0.18,3.4) .. controls (-0.4,2.4) and (-0.35,0.4) .. (0,0);
\draw[popgreen] (0,3.4) ellipse (0.18 and 0.05);
\node[popink] at (1.7,1.6) {\small solid calcitic guard};
\draw[->, popmuted] (1.45,1.55) -- (0.3,1.55);
\node[popink] at (0,-0.4) {\small pointed apex};
\legende{Belemnite guard: bullet-shaped, solid, pointed at one end --- an \emph{internal} shell.}
\end{tikzpicture}
\end{popfigure}

\begin{propriete}[Ammonites and trilobites as index fossils]
Ammonites \textbf{evolved rapidly} (short-lived species), \textbf{spread widely} across the oceans, are abundant and easily recognised: they are therefore \textbf{ideal index fossils for the Mesozoic}. In the same way, trilobites, which evolved quickly and diversified enormously during the early Palaeozoic, are the \textbf{key index fossils of the Palaeozoic}, especially the Cambrian and Ordovician.
\end{propriete}

\courssub{Phylum Brachiopoda}

Brachiopods (lamp shells) are solitary, exclusively marine animals with \textbf{two unequal valves}: a larger \emph{pedicle (ventral) valve} and a smaller \emph{brachial (dorsal) valve}. A fleshy stalk, the \cle{pedicle}, used for attachment to the sea floor, emerges through an opening (the \cle{pedicle opening}) near the beak of the ventral valve. The crucial point is the symmetry: \textbf{each single valve is bilaterally symmetrical about a plane passing through both valves} (perpendicular to the hinge line), whereas in bivalves the symmetry plane lies between the two valves. The shell is calcareous (or, in a few groups, phosphatic). Brachiopods range from the \textbf{Cambrian to the Recent} and were \textbf{dominant in Palaeozoic} seas, where they are excellent index and environmental fossils.

\begin{popfigure}
\begin{tikzpicture}[scale=1.1]
% Brachiopod
\draw[thick, poppink] (0,0) .. controls (-1.8,0.3) and (-1.9,1.8) .. (0,2.2)
 .. controls (1.9,1.8) and (1.8,0.3) .. (0,0);
\draw[poppink] (0,0.25) .. controls (-1.3,0.5) and (-1.4,1.5) .. (0,1.9);
\draw[poppink] (0,0.25) .. controls (1.3,0.5) and (1.4,1.5) .. (0,1.9);
\draw[dashed, popdark, thick] (0,-0.3) -- (0,2.6);
\node[popink] at (1.9,2.5) {\small plane of symmetry};
\node[popink] at (0,-0.6) {\small pedicle opening};
\draw[->, popmuted] (0,-0.45) -- (0,-0.02);
\legende{Brachiopod: the plane of symmetry cuts \emph{through} each valve; note the pedicle opening at the beak.}
\end{tikzpicture}
\end{popfigure}

\begin{methode}[Distinguishing a brachiopod from a bivalve]
\begin{enumerate}
\item Place the shell with the hinge line away from you and identify the two valves.
\item In a \textbf{bivalve}, the two valves are equal mirror images: the symmetry plane passes \emph{between} the valves, but each valve alone is asymmetrical.
\item In a \textbf{brachiopod}, the two valves are unequal (large pedicle valve, small brachial valve), but each valve alone is symmetrical: the symmetry plane passes \emph{through} the valves, perpendicular to the hinge.
\item Confirm with the \textbf{pedicle opening} near the beak: present in brachiopods, absent in bivalves.
\end{enumerate}
\end{methode}

\courssub{Phylum Arthropoda: Class Trilobita}

Trilobites are extinct marine arthropods with a segmented, oval \cle{exoskeleton} divided \emph{longitudinally} into three lobes (a central \textbf{axial lobe} and two \textbf{pleural lobes} --- hence ``tri-lobite'') and \emph{transversely} into three parts: the \cle{cephalon} (head shield, bearing the eyes and the swollen central \cle{glabella}), the \cle{thorax} (several articulated segments) and the \cle{pygidium} (tail plate). The exoskeleton of calcite and chitin was \textbf{moulted} as the animal grew, so a single individual could leave several fossils. Trilobites were benthic crawlers or swimmers in Palaeozoic seas; they range from the \textbf{Cambrian to the Permian}, where they became extinct.

\begin{popfigure}
\begin{tikzpicture}[scale=1.1]
% Trilobite
\draw[thick, popgold] (0,2.6) arc (90:270:0.9 and 0.55);
\draw[thick, popgold] (0,2.6) arc (90:-90:0.9 and 0.55);
% cephalon
\draw[thick, popgold] (0,2.05) .. controls (-0.9,1.95) and (-0.9,1.55) .. (0,1.5);
\draw[thick, popgold] (0,2.05) .. controls (0.9,1.95) and (0.9,1.55) .. (0,1.5);
% thorax segments
\foreach \y in {1.5,1.2,0.9,0.6,0.3}{
 \draw[popgold] (-0.85,\y) .. controls (0,\y-0.12) .. (0.85,\y);
}
\draw[thick, popgold] (-0.85,1.5) -- (-0.85,0.3);
\draw[thick, popgold] (0.85,1.5) -- (0.85,0.3);
% pygidium
\draw[thick, popgold] (-0.85,0.3) .. controls (-0.6,-0.5) and (0.6,-0.5) .. (0.85,0.3);
% axial lobe
\draw[popdark] (-0.22,2.5) -- (-0.22,-0.2);
\draw[popdark] (0.22,2.5) -- (0.22,-0.2);
\node[popink] at (2.0,2.3) {\small cephalon};
\draw[->, popmuted] (1.7,2.25) -- (0.9,2.0);
\node[popink] at (2.0,0.9) {\small thorax};
\draw[->, popmuted] (1.7,0.9) -- (0.9,0.9);
\node[popink] at (2.0,-0.2) {\small pygidium};
\draw[->, popmuted] (1.7,-0.2) -- (0.7,-0.15);
\node[popink] at (-2.2,1.0) {\small axial lobe};
\draw[->, popmuted] (-1.6,1.0) -- (-0.25,1.0);
\legende{Trilobite: cephalon, segmented thorax and pygidium, with a central axial lobe.}
\end{tikzpicture}
\end{popfigure}

\begin{attention}[Trilobites and the Mesozoic]
A very common examination error is to state that trilobites occur in Mesozoic rocks. \textbf{Trilobites became extinct at the end of the Permian} (about 252 million years ago). Finding a trilobite in a bed immediately tells you the rock is \textbf{Palaeozoic}, never Mesozoic or Cenozoic.
\end{attention}

\courssub{Phylum Echinodermata}

Echinoderms are exclusively marine animals with \textbf{pentaradial (five-fold) symmetry} and a skeleton of calcareous plates (\cle{ossicles}). Two groups are common fossils.

\begin{itemize}
\item \textbf{Crinoids} (sea lilies): a cup-shaped body, the \cle{calyx}, bearing feathery arms used to filter food, held above the sea floor by a flexible \cle{stem} made of stacked disc-shaped \textbf{columnal ossicles}. Isolated stem ossicles (``sea-lily beads'') are among the most abundant fossils in Palaeozoic limestones.
\item \textbf{Echinoids} (sea urchins): a globular or flattened rigid \cle{test} of interlocking plates, showing clear five-fold symmetry, often bearing movable spines. They become especially common from the Mesozoic onwards and are useful index fossils in Cretaceous and Cenozoic rocks.
\end{itemize}

\begin{popfigure}
\begin{tikzpicture}[scale=1.0]
% Crinoid
\draw[thick, popblue] (0,1.6) ellipse (0.45 and 0.35);
\foreach \a in {30,90,150}{
 \draw[thick, popblue] ([shift={(0,1.6)}]\a:0.45) .. controls ([shift={(0,2.0)}]\a:0.9) .. ([shift={(0,2.4)}]\a:1.3);
}
\foreach \y in {1.3,1.0,0.7,0.4,0.1}{
 \draw[popblue] (-0.18,\y) -- (0.18,\y);
}
\draw[thick, popblue] (-0.18,1.45) -- (-0.18,-0.05);
\draw[thick, popblue] (0.18,1.45) -- (0.18,-0.05);
\node[popink] at (1.9,1.7) {\small calyx with arms};
\draw[->, popmuted] (1.6,1.65) -- (0.5,1.65);
\node[popink] at (1.7,0.6) {\small stem ossicles};
\draw[->, popmuted] (1.4,0.6) -- (0.22,0.6);
\legende{Crinoid (sea lily): calyx bearing arms, supported by a stem of columnal ossicles.}
\end{tikzpicture}
\end{popfigure}

\courssub{Vertebrates and Fossil Plants}

\textbf{Vertebrates.} Vertebrate hard parts are bones, teeth and scales made of calcium phosphate (apatite), which resists decay better than most tissues.
\begin{itemize}
\item \textbf{Fish}: commonly preserved as \textbf{scales, teeth and bones}; the earliest fish appear in the Palaeozoic.
\item \textbf{Dinosaurs}: known from \textbf{bones, footprints (trace fossils) and eggs}; they dominated land life in the Mesozoic and (except for their descendants, the birds) died out at the end of the Cretaceous.
\item \textbf{Mammals}: \textbf{teeth are the most durable remains} because enamel resists decay; mammal teeth are the commonest mammal fossils and are diagnostic of diet and species.
\end{itemize}

\textbf{Fossil plants.} Plants are preserved as imprints, carbonized films, petrified wood or microfossils (spores and pollen).
\begin{itemize}
\item \textbf{Ferns}: among the earliest land plants, common in Palaeozoic coal measures.
\item \textbf{Gymnosperms}: conifers and relatives; trunks may be preserved as \cle{petrified wood}, in which silica has replaced the woody tissue cell by cell.
\item \textbf{Angiosperms}: flowering plants, preserved as \textbf{leaves and pollen}; pollen grains have extremely resistant walls and are valuable in dating and palaeoenvironmental studies.
\item \textbf{Coal formation}: plant debris accumulating in oxygen-poor swamps is compacted and carbonized through peat and lignite to \textbf{coal} --- an economically vital fossil product.
\end{itemize}

\begin{savaistu}[Fossils in Cameroon]
Sedimentary basins in Cameroon, such as the Douala and Rio del Rey basins of the coastal region, yield marine molluscs, echinoids and abundant microfossils (foraminifera, pollen and spores) that are used to date the Cretaceous--Cenozoic strata and to guide petroleum exploration. Fossil pollen is also used to reconstruct past vegetation and climate in the region.
\end{savaistu}

\courssub{Integration: Comparative Table of the Fourteen Fossil Groups}

\begin{center}
\small
\begin{tabularx}{\linewidth}{|l|l|X|l|l|}
\hline
\rowcolor{popblueL} \textbf{Group} & \textbf{Phylum} & \textbf{Key morphology} & \textbf{Range} & \textbf{Main use} \\
\hline
Foraminifera & Protozoa & Chambered calcareous test & Cambrian--Recent & Index, oil exploration \\
\hline
Sponges & Porifera & Porous body, spicules & Precambrian--Recent & Reef indicators \\
\hline
Corals & Cnidaria & Calcareous cup, septa & Ordovician--Recent & Reefs, warm seas \\
\hline
Graptolites & Hemichordata & Saw-tooth stipes & Cambrian--Carboniferous & Lower Palaeozoic index \\
\hline
Bivalves & Mollusca & Two equal hinged valves & Ordovician--Recent & Environmental indicators \\
\hline
Gastropods & Mollusca & Single coiled shell & Cambrian--Recent & Environmental indicators \\
\hline
Ammonites & Mollusca & Coiled shell, sutures & Devonian--end Cretaceous & Mesozoic index \\
\hline
Belemnites & Mollusca & Bullet-shaped guard & Mesozoic & Jurassic--Cretaceous index \\
\hline
Brachiopods & Brachiopoda & Two unequal valves, pedicle & Cambrian--Recent & Palaeozoic index \\
\hline
Trilobites & Arthropoda & Cephalon, thorax, pygidium & Cambrian--Permian & Palaeozoic index \\
\hline
Crinoids & Echinodermata & Stem ossicles, calyx & Ordovician--Recent & Limestone builders \\
\hline
Echinoids & Echinodermata & Pentaradial test, spines & Ordovician--Recent & Cretaceous--Cenozoic index \\
\hline
Vertebrates & Chordata & Bones, teeth, scales & Ordovician--Recent & Evolution, correlation \\
\hline
Plants & Plantae & Leaves, wood, pollen & Silurian--Recent & Climate, coal formation \\
\hline
\end{tabularx}
\end{center}

\begin{attention}[Belemnites versus orthoconic nautiloids]
Do not confuse a \textbf{belemnite guard} with a \textbf{straight (orthoconic) nautiloid}. The belemnite guard is a \emph{solid, internal} bullet of calcite with \emph{no chambers}. An orthoconic nautiloid is an \emph{external} shell that is clearly \emph{chambered}, with septa and a siphuncle. If you see chambers, it is not a belemnite guard.
\end{attention}

\begin{exoresolu}[Identifying and dating a fossil assemblage]
A limestone bed in a field exercise yields abundant coiled shells with complex frilled suture lines, together with bullet-shaped solid calcite objects and a few sea-urchin tests. (a) Identify the three fossil types. (b) State the most probable age of the bed, justifying your answer.
\tcblower
\textbf{(a)} The coiled shells with complex frilled sutures are \textbf{ammonites}; the bullet-shaped solid objects are \textbf{belemnite guards}; the sea-urchin tests are \textbf{echinoids}. \par
\textbf{(b)} The bed is most probably \textbf{Mesozoic}, and more precisely Jurassic or Cretaceous. Justification: ammonites are excellent Mesozoic index fossils and became extinct at the end of the Cretaceous; belemnites are also characteristic of the Mesozoic (especially Jurassic--Cretaceous); echinoids became abundant from the Mesozoic onwards. The association of ammonites with belemnites is a classic Mesozoic marine assemblage.
\end{exoresolu}

\begin{exercice}[Symmetry test]
A student finds a shell with two valves of clearly different sizes, the larger one bearing a small hole near its pointed beak. (a) Is it a bivalve or a brachiopod? (b) Draw the plane of symmetry and justify your answer. (c) What was the function of the hole?
\end{exercice}

\begin{corrige}[Exercise 1]
\textbf{(a)} It is a \textbf{brachiopod}. \textbf{(b)} The plane of symmetry passes \emph{through} each valve, perpendicular to the hinge line, so that each valve alone is symmetrical; in a bivalve the two valves would be equal mirror images with the symmetry plane between them. \textbf{(c)} The hole is the \textbf{pedicle opening}, through which the fleshy pedicle passed to attach the animal to the sea floor.
\end{corrige}

\begin{exercice}[Dating with index fossils]
A shale contains trilobite fragments and graptolite stipes; an overlying limestone contains ammonites with complex sutures. State the era of each bed and name the principle that allows you to correlate the shale with rocks of the same age elsewhere.
\end{exercice}

\begin{corrige}[Exercise 2]
The shale is \textbf{Palaeozoic}: trilobites are restricted to the Palaeozoic (Cambrian--Permian) and graptolites to the Cambrian--Carboniferous. The limestone is \textbf{Mesozoic}, since ammonites with complex (ammonitic) sutures are Mesozoic index fossils. The principle used is the \textbf{principle of faunal succession} (correlation by index fossils): beds containing the same index fossil assemblage are of the same age, even in different regions.
\end{corrige}

\begin{aretenir}[Key points]
\begin{itemize}
\item Bivalves: two \emph{equal} valves, symmetry plane \emph{between} the valves; Ordovician--Recent.
\item Gastropods: single coiled shell (spire, body whorl, aperture); Cambrian--Recent.
\item Ammonites: coiled chambered shells with suture lines; fast evolution and wide spread make them ideal \textbf{Mesozoic index fossils}; extinct at end-Cretaceous.
\item Belemnites: solid internal guards, never chambered.
\item Brachiopods: two \emph{unequal} valves, symmetry plane \emph{through} the valves, pedicle opening; dominant in the Palaeozoic.
\item Trilobites: cephalon--thorax--pygidium; Cambrian--Permian only; key Palaeozoic index fossils.
\item Crinoids (stem ossicles, calyx) and echinoids (pentaradial test) are echinoderms.
\item Mammal teeth, fish scales and petrified wood are common vertebrate and plant fossils; plant debris in swamps forms coal.
\end{itemize}
\end{aretenir}

\newpage

\coursec{Integration activity}

\courssub{Aims of the activity}

This integration activity closes the chapter on palaeontology by asking you to work exactly as a professional palaeontologist does. Starting from a simulated collection of fossils said to come from the \cle{Figuil} limestone quarries (North Region of Cameroon), you will mobilise, in one single investigation, all the notions studied in this chapter:

\begin{itemize}
\item the \cle{types of fossils} (body fossils, trace fossils, chemical fossils);
\item the \cle{conditions} and \cle{modes of fossilization} (unaltered remains, permineralization, replacement, moulds and casts, carbonization);
\item the \cle{preservation potential} of different organisms;
\item the \cle{classification} and \cle{description} of commonly preserved fossils (ammonites, bivalves, gastropods, corals, sharks, plants);
\item the \cle{uses of fossils} (relative dating with index fossils, reconstruction of palaeoenvironments, economic applications);
\item the \cle{gaps in the fossil record} and the reasons for the absence of certain groups.
\end{itemize}

By the end of the activity you should be able to identify and classify an unknown fossil, infer its mode of fossilization, deduce the depositional environment and relative age of a rock from its fossil content, and defend your conclusions in a short scientific report --- precisely the skills examined in the GCE Advanced Level practical and essay papers.

\begin{aretenir}[What this activity tests]
This is an \emph{integration} exercise: no new notion is introduced. Every answer must combine several ideas from the chapter, and every conclusion must be \textbf{justified by a named specimen}. In the GCE examination, an unjustified statement earns little or no credit.
\end{aretenir}

\courssub{Situation}

\textbf{The Figuil Fossil Inquiry: Reconstructing an Ancient Cameroonian Sea}\par

The town of Figuil, in the North Region of Cameroon, is well known for its limestone deposits, exploited for the manufacture of cement. These limestones belong to the sedimentary formations associated with the \cle{Benue Trough}, a great rift structure that crosses Nigeria and extends into northern Cameroon. During the Cretaceous period, this trough was invaded several times by the sea (marine transgressions), and thick beds of limestone, marl and shale accumulated on its floor.

The Ministry of Secondary Education (MINESEC), in collaboration with a cement company, has received a box of specimens collected by quarry workers in a freshly exposed limestone bed at Figuil. Your class has been commissioned as a team of \emph{palaeontology consultants}. You receive the ``field box'' containing drawings and photographs of seven specimens, labelled A to G, together with the field notes of the collector.

\textbf{Contents of the field box}

\begin{center}
\begin{tabularx}{\linewidth}{|c|X|}
\hline
\rowcolor{popblueL} \textbf{Specimen} & \textbf{Description from the field notebook} \\
\hline
A & Coiled shell, \SI{12}{cm} across, divided internally into chambers separated by wavy partitions (sutures); a pearly layer is visible where the surface is broken. \\
\hline
B & Two symmetrical hinged valves, \SI{5}{cm} long, with concentric growth lines; the shell material is still pearly white. \\
\hline
C & Spirally coiled unchambered shell, \SI{4}{cm} high, shaped like a cone wound around an axis. \\
\hline
D & Massive colony of small cylindrical tubes, each about \SI{3}{mm} in diameter, grouped like a honeycomb; the tubes are made of calcium carbonate. \\
\hline
E & Three pointed, triangular teeth, \SI{3}{cm} high, with a broad root; the enamel is shiny and dark. \\
\hline
F & Several cylindrical tubes, \SI{1}{cm} in diameter, running horizontally inside the rock; no shell or skeleton is present, only the tubes filled with sediment of a different colour. \\
\hline
G & A thin black film on a bedding plane of shale, showing the detailed outline and veins of a leaf. \\
\hline
\end{tabularx}
\end{center}

\textbf{Field notes of the collector}
\begin{itemize}
\item The specimens come from a single grey limestone bed, about \SI{2}{m} thick, intercalated between marls and shales.
\item Specimens A, B, C and D are very abundant; E is rare; F occurs throughout the bed; G was found in a thin shale layer just above the limestone.
\item No bones of land animals and no flowers or fruits were observed anywhere in the quarry.
\end{itemize}

\begin{popfigure}
\begin{center}
\begin{tikzpicture}[scale=0.9]
% limestone bed
\fill[popblueL] (0,0) rectangle (10,2);
\fill[popgreenL] (0,2) rectangle (10,2.8);
\fill[popblueL!60] (0,-1) rectangle (10,0);
\draw (0,0) rectangle (10,2);
\draw (0,2) rectangle (10,2.8);
\draw (0,-1) rectangle (10,0);
\node[align=center] at (5,1) {Limestone bed (specimens A--F)};
\node[align=center] at (5,2.4) {Shale with leaf film (G)};
\node[align=center] at (5,-0.5) {Marl};
% ammonite sketch
\draw[popdark,thick] (1.5,1) circle (0.35);
\draw[popdark,thick] (1.5,1) circle (0.18);
\node[popdark] at (1.5,0.45) {\scriptsize A};
% bivalve
\draw[popdark,thick] (3.5,1) ellipse (0.3 and 0.2);
\node[popdark] at (3.5,0.5) {\scriptsize B};
% coral
\foreach \x in {5.3,5.6,5.9} \draw[popdark,thick] (\x,0.8) rectangle (\x+0.15,1.3);
\node[popdark] at (5.7,0.5) {\scriptsize D};
% burrows
\draw[poporange,thick] (7,1.2) -- (8.4,1.2);
\draw[poporange,thick] (7.2,0.8) -- (8.2,0.8);
\node[poporange] at (7.7,0.45) {\scriptsize F};
% tooth
\draw[popdark,thick] (9.2,0.7) -- (9.45,1.4) -- (9.7,0.7) -- cycle;
\node[popdark] at (9.45,0.45) {\scriptsize E};
% young direction arrow
\draw[->,thick,popdark] (10.4,-0.8) -- (10.4,2.6);
\node[popdark,rotate=90] at (10.7,0.9) {\scriptsize Younging direction};
\end{tikzpicture}
\end{center}
\legende{Simplified log of the fossiliferous bed at Figuil showing the position of the specimens.}
\end{popfigure}

Your team must analyse the assemblage and present its conclusions in a five-minute briefing to the ``MINESEC scientific committee'' (the rest of the class), supported by a one-page written report.

\courssub{Tasks}

Work in teams and answer the following questions in order. Each answer must be justified.

\begin{enumerate}
\item \textbf{Identification and classification.} Using the description grid studied in this chapter, identify each specimen A--G: give its common name, its \cle{phylum} (and class where possible), its main morphological features and the original composition of its hard parts. Present your answer as a table.
\item \textbf{Mode of fossilization.} For each specimen, state the most likely \cle{mode of fossilization} (unaltered remains, permineralization, replacement, mould/cast, carbonization, trace preservation) and justify your choice from the observations in the field notebook.
\item \textbf{Body fossils versus trace fossils.} Separate the specimens into \cle{body fossils} and \cle{trace fossils}. Which specimen is neither a body fossil nor a trace fossil in the strict sense, and why?
\item \textbf{Depositional environment.} Using the assemblage as a whole, deduce the \cle{palaeoenvironment} of deposition of the limestone bed (nature of the water, depth, temperature, energy). Justify each deduction with at least one specimen.
\item \textbf{Relative age.} Explain why specimen A is a good \cle{index fossil}. What type of geological information can it give about the age of the Figuil limestone, and how could this be used to correlate the Figuil beds with the sedimentary formations of the Douala basin?
\item \textbf{Gaps in the record.} Explain why no mammal bones and no flowers are found in the quarry, even though mammals and flowering plants existed later in Earth history. Mobilise at least three distinct reasons (evolutionary, environmental and preservational).
\item \textbf{Preservation potential.} Rank the organisms represented (ammonite, bivalve, gastropod, coral, shark, burrowing worm, leaf) from the highest to the lowest \cle{preservation potential}, and justify the ranking in terms of hard parts, composition and habitat.
\item \textbf{Economic interest.} Write the conclusion of your one-page report: why is this fossiliferous limestone of economic interest to Cameroon (cement industry, stratigraphic correlation, possible associated resources)?
\end{enumerate}

\begin{methode}[How to proceed]
\begin{enumerate}
\item Read the field notebook carefully and underline the diagnostic features of each specimen (chambers, sutures, valves, coiling, tubes, teeth, film).
\item Fill in the identification table first; the rest of the work depends on it.
\item For each deduction (environment, age, preservation), always name the specimen(s) that provide the evidence.
\item Share the tasks within the team, then cross-check: a good consultant never signs a report without verification.
\item Prepare the five-minute briefing: one speaker, one timekeeper, one person answering questions.
\end{enumerate}
\end{methode}

\courssub{Model answer}

\textbf{1. Identification and classification}

\begin{center}
\begin{tabularx}{\linewidth}{|c|l|l|X|}
\hline
\rowcolor{popblueL} \textbf{Sp.} & \textbf{Name} & \textbf{Phylum / Class} & \textbf{Diagnostic features and composition} \\
\hline
A & Ammonite & Mollusca / Cephalopoda & Coiled chambered shell with wavy sutures; shell originally of aragonite (a form of calcium carbonate). \\
\hline
B & Bivalve & Mollusca / Bivalvia & Two hinged symmetrical valves with concentric growth lines; shell of calcium carbonate. \\
\hline
C & Gastropod & Mollusca / Gastropoda & Single spirally coiled unchambered shell; calcium carbonate. \\
\hline
D & Colonial coral & Cnidaria / Anthozoa & Honeycomb colony of calcareous tubes (corallites); calcium carbonate. \\
\hline
E & Shark teeth & Chordata / Chondrichthyes & Pointed triangular teeth with enamel; calcium phosphate (apatite). \\
\hline
F & Burrows & Trace of a worm-like animal & Sediment-filled tubes; no preserved body. \\
\hline
G & Fossil leaf & Plantae (a land plant) & Thin carbon film showing outline and veins. \\
\hline
\end{tabularx}
\end{center}

\textbf{2. Modes of fossilization.} Specimen B retains its pearly shell, so it is preserved as \cle{unaltered hard parts} (or only slightly recrystallized). Specimens A, C and D, embedded in limestone, are most commonly preserved by \cle{permineralization} (mineral matter filling the pores of the shell) and \cle{replacement} (recrystallization of the original aragonite into the more stable calcite, or replacement by silica or pyrite); where the shell has dissolved away, \cle{moulds and casts} result. Specimen E (shark teeth) is preserved as unaltered hard parts, since apatite resists dissolution. Specimen F is a \cle{trace fossil}: preservation of an activity, not of a body. Specimen G is preserved by \cle{carbonization}: the volatile elements (hydrogen, oxygen, nitrogen) escaped during burial, leaving a thin film of carbon --- a carbonized imprint of the leaf.

\textbf{3. Body versus trace fossils.} A, B, C, D, E and G are \cle{body fossils} (remains of the organisms themselves). F is the only \cle{trace fossil} (burrows = evidence of the activity of an animal). Strictly, G is a body fossil preserved as a carbonized imprint; it is not a trace fossil because it represents the organism itself, not its activity.

\textbf{4. Depositional environment.} The assemblage indicates a \cle{warm, shallow, normal-salinity marine sea}:
\begin{itemize}
\item colonial corals (D) need warm, clear, shallow, well-lit water of normal salinity;
\item ammonites (A) and sharks (E) were free-swimming marine animals, showing open marine conditions;
\item bivalves (B) and gastropods (C) lived on or in the bottom sediment;
\item abundant burrows (F) show a soft, oxygenated sea floor;
\item the limestone itself, made of calcium carbonate, accumulates in warm shallow seas.
\end{itemize}
The shale with the leaf (G) above the limestone records slightly quieter water and an input of plant debris from a nearby landmass.

\textbf{5. Relative age.} The ammonite (A) is an excellent \cle{index fossil} because it satisfies the four criteria: it evolved rapidly (many short-lived species, each restricted to a narrow time interval), it is abundant, it is geographically widespread, and it is easily recognised. The identified ammonite species therefore characterises a precise interval of the Mesozoic Era (Cretaceous for the Benue Trough limestones). If the same species were found in the Douala basin, the beds there would be shown to be of the \emph{same relative age}: this is \cle{stratigraphic correlation} by fossils.

\textbf{6. Why no mammals or flowers?} Three complementary reasons:
\begin{itemize}
\item \emph{Evolutionary:} in the Cretaceous seas, large mammals did not exist in marine environments, and flowering plants were only beginning to spread; the common land plants belonged to other groups.
\item \emph{Environmental:} the bed is marine; land mammals and flowers lived far from the site of deposition, so their remains were rarely transported there.
\item \emph{Preservational:} bones and flowers have low preservation potential --- flowers are soft and decay quickly, and bones are destroyed unless buried rapidly; the fossil record is therefore incomplete (\cle{gaps in the fossil record}).
\end{itemize}

\textbf{7. Preservation potential (highest to lowest).} Coral, bivalve and ammonite (thick calcareous skeletons, marine habitat, rapid burial in sediment) $>$ gastropod (thinner shell) $>$ shark (only the teeth are preserved; the cartilaginous skeleton decays) $>$ leaf (soft tissue, preserved only in exceptional quiet, oxygen-poor conditions) $>$ burrowing worm (no hard parts at all: known only from its traces).

\textbf{8. Economic interest.} The fossiliferous limestone of Figuil is a high-purity carbonate rock exploited for \cle{cement manufacture}, a major industrial resource for northern Cameroon. Its index fossils allow \cle{correlation} with other Cretaceous basins (Benue Trough, Douala basin), which guides the search for associated resources (marl, clay, and hydrocarbon source rocks in analogous basins). The assemblage therefore has both scientific and economic value.

\begin{attention}[Common mistakes to avoid in the briefing]
\begin{itemize}
\item Do not call the burrows (F) ``worm fossils'': the worm itself is not preserved --- it is a trace fossil.
\item Do not say the leaf is ``petrified'': a carbon film results from carbonization, not from mineral replacement.
\item An index fossil must be abundant, widespread, short-lived (rapidly evolving) and easily identified --- quote all four criteria.
\item Absence of mammals is not proof that mammals did not exist anywhere: the fossil record is incomplete.
\end{itemize}
\end{attention}

\begin{exercice}[Going further]
A new bed, \SI{50}{cm} thick, is discovered \SI{3}{m} above the studied limestone. It contains the same bivalve as specimen B but a \emph{different} ammonite species, and no coral. Suggest two possible interpretations, and state what additional observation would help you choose between them.
\end{exercice}

\begin{corrige}[Exercise 1]
Two interpretations are acceptable. (1) \emph{Evolutionary change:} the new ammonite species indicates a younger age (faunal succession), while the long-lived bivalve persists; the absence of coral may reflect a local ecological change. (2) \emph{Environmental change:} the bed may record deeper, muddier or less clear water unsuitable for corals, with the ammonite difference being only a sampling effect. To decide, one should examine the rock type (limestone versus marl or shale), look for other depth-sensitive fossils, and check whether the new ammonite is a known index species of a younger stage.
\end{corrige}

\begin{savaistu}[Did you know?]
The limestone quarried at Figuil supplies one of Cameroon's cement plants, showing that palaeontology is not only an academic science: identifying fossiliferous limestone beds helps locate raw materials for industry, and index fossils from basins like the Benue Trough and the Douala basin guide the geological exploration of the whole country.
\end{savaistu}

\newpage

\coursec{Methods and worked examples}

\courssub{Method 1 — Identifying the Type of Fossil}

\begin{methode}[Recognising body, trace and chemical fossils]
When an examiner shows you a specimen, a photograph or a description and asks ``What type of fossil is this?'', proceed as follows:
\begin{enumerate}
\item \textbf{Ask: is it a part of the organism itself?} If yes (bone, shell, leaf, tooth), it is a \cle{body fossil}. Go to step 2. If no, go to step 3.
\item \textbf{Classify the body fossil}: is it \emph{hard} (shell, bone, exoskeleton) or \emph{soft} (skin, muscle, jellyfish — only preserved exceptionally)? State the material preserved.
\item \textbf{Ask: is it evidence of activity?} Footprints, burrows, borings, coprolites (fossil dung), trackways and feeding trails are \cle{trace fossils} (ichnofossils). The organism itself is absent.
\item \textbf{Ask: is it only a chemical signature?} Organic molecules (e.g. hydrocarbons, kerogen, biomarkers) detected in rock are \cle{chemical fossils} (molecular fossils).
\item \textbf{Justify your answer} with one observable feature (e.g. ``the repeated bilobed groove shows locomotion, not anatomy'').
\end{enumerate}
\textbf{Reflex:} trace fossils record \emph{behaviour}; body fossils record \emph{anatomy}; chemical fossils record \emph{biochemistry}.
\end{methode}

\begin{exoresolu}[Classifying fossil specimens]
A field party in the Benue Trough collects four specimens: (A) a coiled shell of an ammonite; (B) a sinuous groove on a sandstone bedding plane; (C) a rounded mass of phosphatic material containing fish scales; (D) black shale yielding hydrocarbons of biological origin. Classify each specimen as a body, trace or chemical fossil, justifying each answer.
\tcblower
\textbf{Specimen A — the ammonite shell.} This is the actual mineralised hard part (shell) of the organism. It is a \textbf{body fossil} (hard-part body fossil). Justification: the coiled, chambered shell is part of the animal's anatomy.

\textbf{Specimen B — the sinuous groove.} No part of an organism is present; the groove records the movement of an animal (probably a worm or mollusc) across the sediment surface. It is a \textbf{trace fossil} (a crawling trail). Justification: it records behaviour (locomotion), not anatomy.

\textbf{Specimen C — the phosphatic mass with fish scales.} This is a \textbf{coprolite} (fossilised faeces), hence a \textbf{trace fossil}. Justification: although it is a solid object, it is a product of the animal's activity (feeding and excretion), not a body part. The enclosed fish scales reveal the predator's diet.

\textbf{Specimen D — hydrocarbons in shale.} These are \textbf{chemical fossils} (biomarkers). Justification: only organic molecules derived from once-living organisms are preserved; no structure or activity is visible.

\textbf{Examiner's tip:} always give the justification — in the GCE, an unreasoned classification earns only partial credit.
\end{exoresolu}

\courssub{Method 2 — Analysing Conditions and Modes of Fossilization}

\begin{methode}[Explaining how a fossil formed]
To answer ``Describe the mode of fossilization of specimen X'':
\begin{enumerate}
\item \textbf{Identify what is preserved}: original hard parts, original soft tissue, replaced material, a cavity, a carbon film, or an impression.
\item \textbf{Match to the mode}:
\begin{itemize}
\item \emph{Unaltered remains}: original material survives (amber, freezing, mummification, unaltered shells).
\item \emph{Permineralization}: minerals (silica, calcite, pyrite) fill pores of bone or wood.
\item \emph{Replacement / petrification}: original material dissolved and substituted molecule by molecule (silicification, pyritization).
\item \emph{Carbonization}: volatile elements escape, leaving a carbon film (leaves, soft-bodied animals).
\item \emph{Mould and cast}: external/internal mould = the cavity or imprint; cast = sediment or mineral filling of the cavity.
\item \emph{Imprint/impression}: shallow mark of a thin organism on fine sediment.
\end{itemize}
\item \textbf{State the conditions} that made it possible: possession of hard parts, rapid burial, fine-grained sediment, quiet (low-energy) water, absence of oxygen (limiting decay and scavenging).
\item \textbf{Link mode to conditions}: e.g. carbonization requires fine sediment and anoxia; permineralization requires porous tissue and mineral-rich groundwater.
\end{enumerate}
\textbf{Reflex:} \emph{conditions} answer ``why was preservation possible?''; \emph{mode} answers ``by what process was it preserved?''. Do not confuse the two in an essay.
\end{methode}

\begin{exoresolu}[Mode of fossilization of a fish in shale]
A complete fish skeleton is found in finely laminated black shale. The bones are flattened, and a thin dark film outlines the body, including the fins. (a) State the mode of fossilization. (b) Explain the conditions that allowed this exceptional preservation. (c) Why are such fossils rare?
\tcblower
\textbf{(a) Mode of fossilization.} The dominant mode is \textbf{carbonization} (distillation): the soft tissues lost their volatile components (hydrogen, oxygen, nitrogen) under pressure during burial, leaving a thin film of \textbf{carbon} that outlines the body and fins. The bones themselves survive as \textbf{unaltered hard parts} (or may be permineralized). So the specimen combines carbonization of soft parts with preservation of the skeleton.

\textbf{(b) Conditions of preservation.}
\begin{enumerate}
\item \textbf{Rapid burial} by fine mud protected the carcass from scavengers and physical destruction.
\item \textbf{Fine-grained sediment} (mud $\rightarrow$ shale) recorded delicate details such as fin rays.
\item \textbf{Anoxic (oxygen-poor) bottom water} inhibited bacterial decay and excluded burrowing scavengers; the black colour and fine lamination of the shale indicate stagnant, undisturbed deposition.
\item \textbf{Quiet, low-energy water} prevented disarticulation of the skeleton — the bones remain in anatomical position.
\end{enumerate}

\textbf{(c) Rarity.} Most carcasses decay, are eaten, or are scattered by currents before burial; anoxic bottom conditions are uncommon; and even buried soft tissues are usually destroyed during compaction. The coincidence of rapid burial, fine sediment and anoxia is rare, so soft-part preservation is exceptional.

\textbf{Key idea:} exceptional preservation sites of this kind are called \emph{Lagerstätten}.
\end{exoresolu}

\courssub{Method 3 — Distinguishing Moulds and Casts}

\begin{methode}[Mould versus cast — a classic GCE trap]
\begin{enumerate}
\item \textbf{Define the mould}: the \emph{negative} impression or cavity left by the organism in the sediment (external mould = outside surface; internal mould = infilling of the inside of a shell, e.g. a steinkern).
\item \textbf{Define the cast}: the \emph{positive} replica formed when the mould is later filled with sediment or mineral matter.
\item \textbf{Memory aid}: the mould is the \emph{hole}; the cast is the \emph{filling}. Mould = concave/negative; cast = convex/positive.
\item \textbf{In an exam diagram}: shading or hatching usually indicates the rock; the empty or differently shaded region is the fossil. Trace the sequence: burial $\rightarrow$ dissolution of shell $\rightarrow$ mould $\rightarrow$ infilling $\rightarrow$ cast.
\end{enumerate}
\end{methode}

\begin{exoresolu}[Interpreting a mould-and-cast sequence]
The diagram below represents stages in the preservation of a bivalve shell. Describe what happens at each stage and name the structures produced.
\tcblower
\begin{popfigure}
\begin{tikzpicture}[scale=0.9]
% Stage 1
\draw[fill=popblueL] (0,0) rectangle (2.6,1.6);
\draw[thick,fill=white] (1.3,0.8) ellipse (0.7 and 0.35);
\node at (1.3,-0.4) {\small Stage 1: Shell buried};
% Stage 2
\draw[fill=popblueL] (3.4,0) rectangle (6,1.6);
\draw[thick,dashed] (4.7,0.8) ellipse (0.7 and 0.35);
\node at (4.7,-0.4) {\small Stage 2: Shell dissolves};
% Stage 3
\draw[fill=popblueL] (6.8,0) rectangle (9.4,1.6);
\draw[thick,fill=popgreenL] (8.1,0.8) ellipse (0.7 and 0.35);
\node at (8.1,-0.4) {\small Stage 3: Cavity filled};
\end{tikzpicture}
\legende{Stages in the formation of a mould and cast of a bivalve shell.}
\end{popfigure}

\textbf{Stage 1 — Burial.} The bivalve dies and its shell is rapidly buried in sediment. The hard shell resists decay while sediment hardens around and inside it.

\textbf{Stage 2 — Dissolution.} Groundwater percolates through the porous sediment and dissolves the calcium carbonate shell. What remains is a hollow cavity whose walls reproduce the outer surface of the shell: this is the \textbf{external mould}. If sediment had filled the inside of the shell before dissolution, the hardened internal filling is the \textbf{internal mould} (steinkern).

\textbf{Stage 3 — Infilling.} Mineral matter (e.g. silica or calcite) or fine sediment enters the cavity and solidifies, producing a \textbf{cast} — a positive three-dimensional replica of the original shell.

\textbf{Conclusion:} the fossil collector may find the mould (the cavity showing external ornament in negative) or the cast (the solid replica). Neither contains any original shell material — an important point: moulds and casts are modes of preservation in which the \emph{form} survives but the \emph{substance} is lost.
\end{exoresolu}

\courssub{Method 4 — Assessing Preservation Potential}

\begin{methode}[Predicting which organisms fossilize]
\begin{enumerate}
\item \textbf{Hard parts?} Organisms with shells, bones, teeth, tests or woody tissue have \emph{high} preservation potential; soft-bodied organisms (jellyfish, worms) have \emph{low} potential.
\item \textbf{Composition of hard parts}: calcite and silica resist better than aragonite (which recrystallizes or dissolves easily); chitin and phosphate (bones) are intermediate.
\item \textbf{Habitat}: marine organisms living in or near sediment (benthos) are far more likely to be buried than land organisms or open-water swimmers.
\item \textbf{Population size and durability in time}: abundant, widespread, long-ranging species are common as fossils.
\item \textbf{Conclude}: preservation potential = f(hard parts, mineralogy, habitat, abundance).
\end{enumerate}
\end{methode}

\begin{exoresolu}[Comparing preservation potentials]
Arrange the following in order of decreasing preservation potential and justify: (a) a jellyfish; (b) an oyster with a thick calcite shell living on the sea floor; (c) an aragonite-shelled gastropod; (d) a forest bird.
\tcblower
\textbf{Order (highest to lowest): (b) oyster $>$ (c) gastropod $>$ (d) bird $>$ (a) jellyfish.}

\textbf{(b) Oyster — highest potential.} It possesses thick, durable \textbf{calcite} hard parts (calcite is chemically stable), lives \textbf{on the sea floor} where sediment accumulates, and is therefore buried rapidly after death — often in life position.

\textbf{(c) Gastropod.} It has a hard shell and a marine benthic habitat, but the shell is \textbf{aragonite}, which is metastable and commonly dissolved or recrystallized during diagenesis; such shells are often preserved only as moulds.

\textbf{(d) Forest bird.} Bones are phosphatic and reasonably durable, but the \textbf{terrestrial habitat} is dominated by erosion rather than deposition; carcasses are scavenged, weathered and rarely buried. Land vertebrates are correspondingly rare as fossils.

\textbf{(e) Jellyfish — lowest potential.} Entirely \textbf{soft-bodied}; it decays within days unless buried instantly in anoxic, fine sediment. Its fossil record is negligible.

\textbf{General conclusion:} the fossil record is strongly biased toward marine, hard-shelled, bottom-dwelling organisms — a bias that must be remembered when interpreting past life.
\end{exoresolu}

\courssub{Method 5 — Using Fossils: Index Fossils and Relative Dating}

\begin{methode}[Applying index fossils to date and correlate strata]
\begin{enumerate}
\item \textbf{Recall the criteria} of a good index (zone) fossil: wide geographic distribution, short vertical (time) range, abundance, easy identification, independence of facies.
\item \textbf{Read the stratigraphic column}: note which fossils occur in each bed.
\item \textbf{Apply the principle of faunal succession}: fossil assemblages succeed one another in a definite, irreversible order; identical assemblages = same age.
\item \textbf{Correlate}: match beds in different localities that share the same index fossil(s); they are contemporaneous even if the rock types differ.
\item \textbf{Date}: a bed is as old as the narrowest-ranging index fossil it contains; an assemblage gives a tighter age than a single fossil.
\end{enumerate}
\textbf{Reflex:} correlation matches \emph{time-equivalence}, not rock similarity.
\end{methode}

\begin{exoresolu}[Correlating two sections with index fossils]
Two measured sections, 80 km apart in the Benue Trough, give the following successions (oldest at base):

\begin{center}
\begin{tabular}{|c|l|l|}
\hline
\rowcolor{popblueL} Bed & Locality 1 & Locality 2 \\
\hline
4 (top) & Sandstone, no fossils & Shale with ammonite \emph{X} \\
3 & Shale with ammonite \emph{X} & Limestone with \emph{X} + \emph{Y} \\
2 & Limestone with \emph{X} + \emph{Y} & Sandstone, no fossils \\
1 (base) & Shale with trilobite \emph{Z} & Shale with trilobite \emph{Z} \\
\hline
\end{tabular}
\end{center}

(a) Correlate the beds between the two localities. (b) Which bed at Locality 2 is missing at Locality 1, and what does this imply? (c) Why is ammonite \emph{X} a better index fossil than trilobite \emph{Z} here?
\tcblower
\textbf{(a) Correlation.} Using the principle of faunal succession, beds containing the same fossil assemblage are time-equivalent:
\begin{itemize}
\item Bed 1 (Locality 1) = Bed 1 (Locality 2): both contain trilobite \emph{Z}.
\item Bed 2 (Locality 1) = Bed 3 (Locality 2): both contain the assemblage \emph{X} + \emph{Y}.
\item Bed 3 (Locality 1) = Bed 4 (Locality 2): both contain ammonite \emph{X} alone.
\end{itemize}
Note that lithologies differ (shale vs limestone) yet correlation is still valid — correlation is based on \emph{age}, not rock type.

\textbf{(b) Missing bed.} Bed 2 of Locality 2 (unfossiliferous sandstone) has no counterpart at Locality 1. This implies either that sandstone was deposited only at Locality 2 (a lateral facies change — the sandstone passes laterally into the fossiliferous beds), or that an equivalent bed at Locality 1 was eroded or not deposited (a hiatus). The fossil evidence shows deposition at Locality 1 was continuous through this interval, so a \textbf{facies change} is the better interpretation.

\textbf{(c) Quality as index fossil.} Ammonite \emph{X} has a \textbf{shorter time range} (it appears only in the upper beds) and, being a free-swimming marine animal, is \textbf{widely distributed and independent of facies} (found in both shale and limestone). Trilobite \emph{Z} ranges through the lower beds of both sections and gives only a broad age. The narrower the range and the wider the distribution, the better the index fossil.

\textbf{Examiner's tip:} always state the principle used (faunal succession) explicitly — it carries marks.
\end{exoresolu}

\courssub{Method 6 — Interpreting Gaps in the Fossil Record}

\begin{methode}[Explaining why the record is incomplete]
When asked to account for gaps or missing links:
\begin{enumerate}
\item \textbf{Biological causes}: soft-bodied organisms rarely fossilize; small or delicate skeletons are destroyed.
\item \textbf{Geological causes}: erosion removes strata (unconformities); metamorphism destroys fossils; deep burial makes rocks inaccessible; non-deposition creates hiatuses.
\item \textbf{Collection bias}: many rocks are unexposed, unstudied, or in inaccessible regions; rare species are seldom found.
\item \textbf{Consequence}: transitional forms are scarce, lineages appear to jump, and the record is biased toward shallow-marine hard-shelled faunas.
\item \textbf{Exam skill}: link each gap to a specific cause rather than writing a vague list.
\end{enumerate}
\end{methode}

\begin{exoresolu}[Explaining an apparent sudden appearance]
In a stratigraphic column, a diverse ammonite fauna appears abruptly above an unconformity, with no transitional forms below. A student concludes that the ammonites ``evolved suddenly''. Criticise this conclusion using your knowledge of gaps in the fossil record.
\tcblower
The student's conclusion confuses a \textbf{gap in the record} with a \textbf{gap in evolution}. Several explanations must be considered:

\textbf{1. The unconformity itself.} An unconformity represents a period of erosion or non-deposition. The strata that would have recorded the gradual evolution of the ammonites were either never deposited or were eroded away. The ``sudden appearance'' is an artefact of the missing time interval.

\textbf{2. Destruction of fossils.} Even where rocks exist, fossils may have been destroyed by metamorphism, recrystallization or dissolution during diagenesis.

\textbf{3. Preservation bias.} The earliest, transitional ammonites may have been rare, small, or inhabited environments where burial was unlikely, so they left few fossils.

\textbf{4. Collection failure.} Transitional forms may exist in unexposed or unsampled rocks elsewhere; absence of evidence is not evidence of absence.

\textbf{Conclusion.} The correct interpretation is that the fossil record here is \textbf{incomplete}: the unconformity and preservation biases removed the record of gradual change. Evolution should be inferred from the most complete sections available worldwide, not from a single gappy column.
\end{exoresolu}

\courssub{Method 7 — Describing and Classifying a Fossil Systematically}

\begin{methode}[Writing a full palaeontological description (Integration skill)]
For the extended ``describe the fossil'' questions (Lessons 36--49), use a fixed plan:
\begin{enumerate}
\item \textbf{Classification}: kingdom, phylum, class (e.g. Phylum Mollusca, Class Cephalopoda).
\item \textbf{Morphology}: shape, symmetry, size, main parts (for an ammonite: coiled chambered shell, sutures, siphuncle; for a trilobite: cephalon, thorax, pygidium, three longitudinal lobes).
\item \textbf{Composition and mode of preservation}: calcite/aragonite/chitin; mould, cast, replaced shell, etc.
\item \textbf{Mode of life and habitat}: benthic/nektonic/planktonic; marine/freshwater/terrestrial; feeding type.
\item \textbf{Geological range and use}: stratigraphic range, value as index fossil, palaeoenvironmental indicator.
\end{enumerate}
\textbf{Reflex:} a labelled diagram earns marks — always sketch and label.
\end{methode}

\begin{exoresolu}[Full description of a trilobite]
Give a systematic description of trilobites under the headings: classification, morphology, preservation, mode of life, and geological importance. Include a labelled diagram.
\tcblower
\textbf{1. Classification.} Kingdom Animalia; Phylum \textbf{Arthropoda}; Class \textbf{Trilobita} (extinct marine arthropods).

\textbf{2. Morphology.} The body is oval, dorsoventrally flattened, and divided \emph{longitudinally} into three lobes (a central axial lobe and two pleural lobes — hence ``trilobite'') and \emph{transversely} into three regions:
\begin{itemize}
\item \textbf{Cephalon} (head): bears the eyes, antennae and mouth; often with genal spines.
\item \textbf{Thorax}: numerous similar articulated segments allowing enrolment (rolling into a ball for protection).
\item \textbf{Pygidium} (tail): fused posterior segments.
\end{itemize}

\begin{popfigure}
\begin{tikzpicture}[scale=1.0]
\draw[thick,fill=popgreenL] (0,0) ellipse (1.6 and 2.2);
% axial lobe
\draw[thick,fill=popblueL] (0,0) ellipse (0.45 and 2.0);
% cephalon / thorax / pygidium divisions
\draw[thick] (-1.45,1.15) -- (1.45,1.15);
\draw[thick] (-1.35,-1.35) -- (1.35,-1.35);
% thoracic segments
\foreach \y in {0.85,0.5,0.15,-0.2,-0.55,-0.9} {\draw (-1.5,\y) -- (1.5,\y);}
\node at (3.1,1.5) {\small Cephalon}; \draw[->] (2.2,1.45) -- (0.9,1.5);
\node at (3.1,0.1) {\small Thorax}; \draw[->] (2.3,0.1) -- (1.1,0.1);
\node at (3.1,-1.7) {\small Pygidium}; \draw[->] (2.2,-1.65) -- (0.8,-1.7);
\node at (-3.3,0.6) {\small Axial lobe}; \draw[->] (-2.4,0.55) -- (-0.4,0.5);
\node at (-3.3,-0.5) {\small Pleural lobe}; \draw[->] (-2.4,-0.5) -- (-1.2,-0.5);
\end{tikzpicture}
\legende{Labelled sketch of a trilobite (dorsal view).}
\end{popfigure}

\textbf{3. Preservation.} The dorsal exoskeleton of \textbf{chitin} impregnated with calcium carbonate fossilizes well; commonly found as dorsal shields, moulted exoskeletons, or enrolled specimens; soft ventral parts (legs, gills) are preserved only exceptionally.

\textbf{4. Mode of life.} Exclusively \textbf{marine}; mostly \textbf{benthic}, crawling on or burrowing in the sea floor; some could swim or float; feeding on small particles and organic debris.

\textbf{5. Geological importance.} Trilobites range through the \textbf{Palaeozoic Era} (Cambrian to Permian) and evolved rapidly; individual species have \textbf{short ranges and wide distribution}, making them outstanding \textbf{index fossils}, especially for the Cambrian and Ordovician. They also indicate shallow-marine palaeoenvironments.

\textbf{Examiner's tip:} in integration questions, marks are distributed across all five headings plus the labelled diagram — never omit the diagram or the geological importance.
\end{exoresolu}

\courssub{Method 8 — Explaining Extinction}

\begin{methode}[Analysing causes of fossil extinction]
\begin{enumerate}
\item \textbf{Distinguish background extinction} (slow, continuous loss of species through competition and environmental change) from \textbf{mass extinction} (rapid, global loss of many groups, e.g. end-Permian, end-Cretaceous).
\item \textbf{List possible causes with mechanisms}: climatic change (loss of habitat), sea-level change (loss of shallow seas), volcanic activity (dust, gases, climate forcing), meteorite impact (dust blocking sunlight, collapse of food chains), competition and predation, disease.
\item \textbf{Use evidence}: sudden disappearance in the record, iridium anomalies, soot layers, synchronous global losses.
\item \textbf{Structure the answer}: cause $\rightarrow$ environmental effect $\rightarrow$ biological consequence.
\end{enumerate}
\end{methode}

\begin{exoresolu}[Extinction of the dinosaurs]
The non-avian dinosaurs disappear abruptly from the fossil record at the end of the Cretaceous Period, together with ammonites and many marine plankton. (a) Is this a background or a mass extinction? Justify. (b) Describe one widely accepted cause and its mechanism. (c) Why did some groups (e.g. small mammals) survive?
\tcblower
\textbf{(a) Type of extinction.} This is a \textbf{mass extinction}: many unrelated groups (dinosaurs, ammonites, marine reptiles, much plankton) disappeared \textbf{simultaneously and globally} within a geologically short time, far above the normal background rate.

\textbf{(b) Cause and mechanism.} A widely accepted cause is the \textbf{impact of a large meteorite} (supported by a global clay layer enriched in iridium, an element rare on Earth but common in meteorites, together with shocked quartz and soot). The mechanism:
\begin{enumerate}
\item The impact threw enormous quantities of dust and aerosols into the atmosphere.
\item Sunlight was blocked for months to years, causing global cooling and halting photosynthesis.
\item Food chains collapsed from the base: plankton and plants died, then herbivores, then carnivores — explaining the simultaneous loss of marine plankton, ammonites and dinosaurs.
\end{enumerate}
Massive volcanic activity (e.g. the Deccan Traps) may have contributed additional climatic stress.

\textbf{(c) Survival of some groups.} Small mammals survived because they were \textbf{small} (low food requirements), could \textbf{burrow} or shelter, were often \textbf{nocturnal}, and could feed on seeds, insects and detritus that persisted when plant production failed. Survival in mass extinctions favours small, generalist, sheltered organisms.

\textbf{Key idea:} extinction removes dominant groups and opens ecological space — the disappearance of dinosaurs allowed the later radiation of mammals, showing that extinction is a driving force in the history of life.
\end{exoresolu}

\newpage

\coursec{Exercises}

\courssub{I check my knowledge}

\begin{exercice}[Multiple choice questions]
For each question, choose the single correct answer and justify your choice in one sentence.
\begin{enumerate}
\item The branch of palaeontology that studies fossil micro-organisms such as foraminifera is:
\begin{enumerate}
\item palaeobotany \item micropalaeontology \item palynology \item ichnology
\end{enumerate}
\item A fossil formed when the original shell dissolves and the resulting cavity is later filled with sediment or mineral matter is called:
\begin{enumerate}
\item a carbonized film \item a cast \item a permineralized bone \item an imprint
\end{enumerate}
\item Which of the following is a \emph{trace fossil}?
\begin{enumerate}
\item an ammonite shell \item a crinoid stem \item a dinosaur footprint \item a silicified tree trunk
\end{enumerate}
\item The best index fossil is one that is:
\begin{enumerate}
\item abundant, widespread, easily identified, and of short geological range
\item rare, localized, and of long geological range
\item large, heavy, and resistant to weathering
\item found only in one type of rock
\end{enumerate}
\end{enumerate}
\end{exercice}

\begin{exercice}[True or false — justify]
State whether each statement is true or false, and justify your answer in one or two sentences.
\begin{enumerate}
\item Fossils are commonly found well preserved in igneous rocks.
\item Carbonization preserves the soft tissues of organisms as a thin film of carbon.
\item The fossil record is a complete and unbiased archive of all past life.
\item A mass extinction is a short geological event during which a large proportion of species disappears worldwide.
\end{enumerate}
\end{exercice}

\begin{exercice}[Definitions]
Define each of the following terms in two to three lines, giving one example where appropriate:
\begin{enumerate}
\item fossil;
\item permineralization;
\item index fossil;
\item background extinction;
\item mould and cast.
\end{enumerate}
\end{exercice}

\begin{exercice}[Quick classification]
Classify each of the following fossils by stating its \textbf{phylum} (and class where possible): foraminifera, rugose coral, crinoid, belemnite, trilobite, brachiopod, graptolite, ammonite. Then state which of them lived only during the Palaeozoic era.
\end{exercice}

\courssub{I practise}

\begin{exercice}[Comparing modes of fossilization]
Construct a table with four columns (mode of fossilization, process involved, type of organism/part commonly preserved, one example) to distinguish \cle{permineralization}, \cle{replacement}, \cle{carbonization} and \cle{mould-and-cast} preservation.
\end{exercice}

\begin{exercice}[The case of the empty limestone]
A limestone bed in a quarry contains no original shell material: only external moulds, internal moulds and casts of bivalve shells are found. Reconstruct, step by step, the diagenetic history of the original shells, and state what this tells you about the composition of the original shells and the chemistry of the pore waters.
\end{exercice}

\begin{exercice}[Footprints in sandstone]
Dinosaur footprints are discovered in a sandstone bed.
\begin{enumerate}
\item State the type of fossil represented and the category of fossils to which it belongs.
\item Explain how such footprints can be preserved.
\item Give two kinds of information that footprints provide which bones alone cannot give.
\end{enumerate}
\end{exercice}

\begin{exercice}[Correlating the Douala basin]
Three stratigraphic columns (A, B, C) were logged in wells of the Douala sedimentary basin. Column A contains, from base to top: shale with foraminifera \emph{F1}; sandstone; limestone with ammonite \emph{Am1}. Column B contains: limestone with \emph{Am1}; shale with foraminifera \emph{F2}. Column C contains: shale with \emph{F1}; sandstone; limestone with \emph{Am1}; shale with \emph{F2}.
\begin{enumerate}
\item Draw the three columns side by side and correlate the beds.
\item Justify each correlation line, explaining why \emph{F1}, \emph{F2} and \emph{Am1} are useful for this exercise.
\item State two properties that make a fossil suitable for such correlations.
\end{enumerate}
\end{exercice}

\courssub{I prepare for the GCE}

\begin{exercice}[Structured question — fossils and fossilization \hfill (20 marks)]
\begin{enumerate}
\item Define the term \emph{fossil}. \hfill (2 marks)
\item List and briefly explain \textbf{four} conditions that favour the fossilization of an organism. \hfill (8 marks)
\item Explain why fossils are almost never found in igneous rocks, and why they are only exceptionally preserved in metamorphic rocks. \hfill (6 marks)
\item Name two sedimentary environments in which fossilization potential is high, justifying each answer. \hfill (4 marks)
\end{enumerate}
\end{exercice}

\begin{exercice}[Specimen identification \hfill (20 marks)]
The figure below (to be drawn by the candidate from specimens studied in class) shows four fossils: an \textbf{ammonite}, a \textbf{trilobite}, a \textbf{brachiopod} and a \textbf{bivalve}. For \emph{each} specimen:
\begin{enumerate}
\item give the phylum (and class for the molluscs); \hfill (4 marks)
\item describe three key morphological features visible on the specimen (symmetry, shell structure, segmentation, ornamentation); \hfill (8 marks)
\item state its mode of life and habitat; \hfill (4 marks)
\item give its geological range and indicate whether it can serve as an index fossil, justifying your answer. \hfill (4 marks)
\end{enumerate}
\end{exercice}

\begin{exercice}[Essay — the fossil record \hfill (25 marks)]
\emph{``The fossil record is incomplete, yet it remains indispensable to the geologist.''}

Discuss this statement in a structured essay. Your answer should:
\begin{enumerate}
\item explain the causes of gaps in the fossil record (preservation potential, conditions of fossilization, destruction by erosion, metamorphism and non-deposition); \hfill (9 marks)
\item show how the different modes of fossilization determine what can and cannot be preserved; \hfill (6 marks)
\item demonstrate, with at least four developed examples, the uses of fossils (stratigraphic correlation and relative dating, palaeoenvironmental and palaeogeographic reconstruction, evidence for evolution, economic applications such as petroleum exploration in the Douala and Rio del Rey basins); \hfill (8 marks)
\item end with a balanced conclusion. \hfill (2 marks)
\end{enumerate}
\end{exercice}

\newpage

\coursec{Solutions to the exercises}

\courssub{Solutions to \emph{I check my knowledge}}

\begin{corrige}[Exercise 1 — Multiple choice questions]
\begin{enumerate}
\item \textbf{Answer: (b) micropalaeontology.} Micropalaeontology is the branch devoted to microscopic fossils (foraminifera, ostracods, conodonts), whereas palaeobotany studies fossil plants, palynology studies spores and pollen, and ichnology studies trace fossils.
\item \textbf{Answer: (b) a cast.} When the original shell dissolves, it leaves a \cle{mould}; when this cavity is later filled by sediment or mineral matter, the infilling reproduces the shell's form and is called a \cle{cast}.
\item \textbf{Answer: (c) a dinosaur footprint.} A footprint records the \emph{activity} of an organism, not its body, so it is a trace fossil; the ammonite shell, crinoid stem and silicified trunk are all body fossils.
\item \textbf{Answer: (a) abundant, widespread, easily identified, and of short geological range.} A good index fossil must combine wide geographic distribution and abundance (so it is found often) with a short vertical range (so it dates the rock precisely).
\end{enumerate}
\end{corrige}

\begin{corrige}[Exercise 2 — True or false]
\begin{enumerate}
\item \textbf{False.} Igneous rocks crystallize from magma or lava at very high temperatures (typically $>800\ ^\circ\mathrm{C}$), which destroy any organic remains; fossils are found essentially in sedimentary rocks.
\item \textbf{True.} During carbonization, volatile elements (H, O, N) are driven off by pressure and heat during burial, leaving a thin film of carbon that can preserve the outline and fine details of soft tissues (e.g. leaves, insects, graptolites).
\item \textbf{False.} The fossil record is highly incomplete and biased: only organisms with hard parts buried quickly in suitable environments fossilize, and many fossil-bearing rocks are later eroded or metamorphosed.
\item \textbf{True.} A mass extinction is a geologically brief event (e.g. the end-Cretaceous event) during which a large proportion of species disappears worldwide, well above the normal background extinction rate.
\end{enumerate}
\end{corrige}

\begin{corrige}[Exercise 3 — Definitions]
\begin{enumerate}
\item \textbf{Fossil:} any preserved remain, impression or trace of an organism that lived in the geological past, found in sedimentary rocks (or occasionally other materials such as amber or ice). Example: an ammonite shell in Jurassic limestone.
\item \textbf{Permineralization:} a mode of fossilization in which mineral-rich groundwater deposits minerals (silica, calcite, pyrite) inside the pores and cavities of a porous hard part (bone, wood) without destroying the original structure. Example: petrified wood.
\item \textbf{Index fossil:} a fossil species that is abundant, geographically widespread, easily identified and restricted to a short interval of geological time, making it ideal for dating and correlating rock layers. Example: the ammonite genus \emph{Acanthoceras} for part of the Cretaceous.
\item \textbf{Background extinction:} the normal, continuous, low-rate extinction of species caused by ordinary ecological pressures (competition, predation, local environmental change), as opposed to catastrophic mass extinctions.
\item \textbf{Mould and cast:} a mould is the negative impression (external or internal) left in sediment after dissolution of a shell; a cast is the positive replica formed when that mould is later filled with sediment or mineral matter.
\end{enumerate}
\end{corrige}

\begin{corrige}[Exercise 4 — Quick classification]
\begin{center}
\begin{tabularx}{\linewidth}{|l|X|X|}\hline
\rowcolor{popblueL} \textbf{Fossil} & \textbf{Phylum} & \textbf{Class} \\ \hline
Foraminifera & Foraminifera (traditionally Protozoa/Protista) & Globothalamea / Tubothalamea (modern) \\ \hline
Rugose coral & Cnidaria (Coelenterata) & Anthozoa \\ \hline
Crinoid & Echinodermata & Crinoidea \\ \hline
Belemnite & Mollusca & Cephalopoda \\ \hline
Trilobite & Arthropoda & Trilobita \\ \hline
Brachiopod & Brachiopoda & Rhynchonellata (Articulata) / Lingulata (Inarticulata) \\ \hline
Graptolite & Hemichordata & Graptolithina \\ \hline
Ammonite & Mollusca & Cephalopoda \\ \hline
\end{tabularx}
\end{center}
\textbf{Note:} Modern classification places Foraminifera in their own phylum (Foraminifera or Retaria); the traditional ``Protozoa'' is shown for GCE compatibility. Fossils that lived \textbf{only during the Palaeozoic era}: the \textbf{rugose corals} (Ordovician--Permian), the \textbf{trilobites} (Cambrian--Permian) and the \textbf{graptolites} (Cambrian--Carboniferous, essentially Ordovician--Silurian for the planktonic forms). The others survived into the Mesozoic (ammonites, belemnites) or are still living today (foraminifera, crinoids, brachiopods).
\end{corrige}

\courssub{Solutions to \emph{I practise}}

\begin{corrige}[Exercise 5 — Comparing modes of fossilization]
\begin{center}
\begin{tabularx}{\linewidth}{|l|X|X|X|}\hline
\rowcolor{popblueL} \textbf{Mode} & \textbf{Process} & \textbf{Organism/part preserved} & \textbf{Example} \\ \hline
Permineralization & Minerals (silica, calcite, pyrite) precipitate from groundwater into pores and cavities; original material largely retained & Porous hard parts: bones, wood, shells & Petrified (silicified) wood \\ \hline
Replacement & Original material (aragonite, calcite, apatite) dissolved and substituted molecule by molecule by another mineral (silica, pyrite, calcite) & Shells, corals, bones; fine structure may be retained & Pyritized or silicified ammonite shells \\ \hline
Carbonization & Volatile elements (H, O, N) expelled by burial pressure and heat, leaving a thin carbon film & Soft-bodied organisms, leaves, insects, graptolites & Carbon film of a fern leaf in shale \\ \hline
Mould and cast & Shell buried, dissolved away leaving a mould; the cavity is later filled by sediment or cement forming a cast & Shells of molluscs and brachiopods in limestone & Bivalve casts in a dissolved limestone bed \\ \hline
\end{tabularx}
\end{center}
\textbf{Key distinctions:} permineralization \emph{adds} mineral matter to existing pores; replacement \emph{substitutes} the original material; carbonization \emph{removes} everything except carbon; mould-and-cast \emph{destroys} the original shell but preserves its shape.
\end{corrige}

\begin{corrige}[Exercise 6 — The case of the empty limestone]
\textbf{Step-by-step diagenetic reconstruction:}
\begin{enumerate}
\item \textbf{Life and death:} bivalves lived on or in the carbonate mud; after death their shells were buried by lime mud before scavenging or transport could destroy them.
\item \textbf{Early lithification:} the surrounding sediment was compacted and cemented into limestone, enclosing the shells; external moulds of the outer shell surfaces formed against the sediment.
\item \textbf{Internal filling:} sediment or sparry calcite cement filled the interiors of some shells, forming internal moulds (steinkerns) or future casts.
\item \textbf{Dissolution:} percolating pore waters, which must have been \emph{acidic and undersaturated with respect to calcium carbonate}, dissolved the shell material completely, leaving empty cavities (external and internal moulds).
\item \textbf{Infilling:} later fluids deposited sediment or crystalline calcite in the cavities, producing casts that reproduce the shells' external form.
\end{enumerate}
\textbf{Conclusions:}
\begin{itemize}
\item The original shells were made of a \emph{soluble carbonate}, most probably \cle{aragonite} (the commonest bivalve shell mineralogy), which is less stable than calcite and dissolves readily during diagenesis.
\item The pore waters were \emph{acidic and undersaturated in \ce{CaCO3}} (aggressive meteoric or burial fluids), capable of dissolving the shells while the already-cemented calcitic matrix survived; the later infilling indicates a subsequent change to \emph{supersaturated, cement-precipitating} fluids.
\end{itemize}
\end{corrige}

\begin{corrige}[Exercise 7 — Footprints in sandstone]
\begin{enumerate}
\item \textbf{Type of fossil:} the footprints are \cle{trace fossils} (ichnofossils); they belong to the category of \emph{traces of activity} (tracks/trails), as opposed to body fossils.
\item \textbf{Preservation:} the dinosaur walked across soft, wet, fine-grained sediment (e.g. a floodplain or tidal-flat mud); its feet impressed hollows; the prints were then rapidly covered by a layer of sand of different grain size, which filled and protected them; burial and lithification turned the filling into sandstone and the underlying layer into the counterpart. Rapid burial before rain, wind or the next tide erased the prints was essential.
\item \textbf{Information that bones cannot give:}
\begin{itemize}
\item \emph{Behaviour and locomotion:} stride length and trackway spacing allow estimation of walking or running speed, and whether the animal moved alone or in herds.
\item \emph{Soft anatomy and activity:} the number and arrangement of toes, presence of claws or pads, and the fact that the animal was alive and moving at that precise place and time (bones are often transported far from the habitat).
\end{itemize}
\end{enumerate}
\end{corrige}

\begin{corrige}[Exercise 8 — Correlating the Douala basin]
\textbf{(a) Correlation of the columns:}

\begin{popfigure}
\begin{tikzpicture}[scale=0.95]
% Column A
\draw[fill=popblueL] (0,0) rectangle (1.6,1.6); \node[align=center] at (0.8,0.8) {\small shale F1};
\draw[fill=popgold!40] (0,1.6) rectangle (1.6,2.8); \node[align=center] at (0.8,2.2) {\small sandstone};
\draw[fill=popgreenL] (0,2.8) rectangle (1.6,4.0); \node[align=center] at (0.8,3.4) {\small lst. Am1};
\node[below] at (0.8,0) {\textbf{A}};
% Column B
\draw[fill=popgreenL] (3.4,1.2) rectangle (5.0,2.4); \node[align=center] at (4.2,1.8) {\small lst. Am1};
\draw[fill=popblueL] (3.4,2.4) rectangle (5.0,4.0); \node[align=center] at (4.2,3.2) {\small shale F2};
\node[below] at (4.2,1.2) {\textbf{B}};
% Column C
\draw[fill=popblueL] (6.8,0) rectangle (8.4,1.2); \node[align=center] at (7.6,0.6) {\small shale F1};
\draw[fill=popgold!40] (6.8,1.2) rectangle (8.4,2.0); \node[align=center] at (7.6,1.6) {\small sandstone};
\draw[fill=popgreenL] (6.8,2.0) rectangle (8.4,3.0); \node[align=center] at (7.6,2.5) {\small lst. Am1};
\draw[fill=popblueL] (6.8,3.0) rectangle (8.4,4.0); \node[align=center] at (7.6,3.5) {\small shale F2};
\node[below] at (7.6,0) {\textbf{C}};
% correlation lines
\draw[thick,poporange,dashed] (0,0) -- (6.8,0); % base F1 A-C
\draw[thick,poporange,dashed] (0,1.6) -- (6.8,1.2); % top F1
\draw[thick,poporange,dashed] (0,2.8) -- (6.8,2.0); % base lst
\draw[thick,poporange,dashed] (0,4.0) -- (3.4,1.2); % Am1 A-B base
\draw[thick,poporange,dashed] (3.4,2.4) -- (6.8,3.0); % top Am1 B-C
\draw[thick,poporange,dashed] (3.4,4.0) -- (6.8,4.0); % top F2
\end{tikzpicture}
\legende{Lithostratigraphic correlation of wells A, B and C using the marker fossils F1, Am1 and F2 (dashed lines show time-equivalent horizons).}
\end{popfigure}

\textbf{(b) Justification:}
\begin{itemize}
\item The \emph{F1 shale} of A is correlated with the basal shale of C because both contain the same foraminifer \emph{F1}, which characterizes the same time interval.
\item The \emph{Am1 limestone} is the key marker bed: it is recognized in all three wells by the ammonite \emph{Am1}, proving that the limestone of B is the same bed as the top limestone of A and the middle limestone of C.
\item The \emph{F2 shale} of B is correlated with the top shale of C by the foraminifer \emph{F2}; it is younger than the Am1 limestone, consistent with its position above it.
\item \emph{F1}, \emph{F2} and \emph{Am1} are useful because each is restricted to a short stratigraphic interval and is easily recognized, so each identifies a precise time horizon across the basin.
\end{itemize}
\textbf{(c) Two properties of a good correlation fossil:}
\begin{enumerate}
\item a \textbf{short geological (vertical) range}, so that its presence dates the bed precisely;
\item a \textbf{wide geographic distribution} (with abundance and easy identification), so that it occurs in many sections and can be found in every well.
\end{enumerate}
\end{corrige}

\courssub{Solutions to \emph{I prepare for the GCE}}

\begin{corrige}[Exercise 9 — Structured question: fossils and fossilization (20 marks)]
\textbf{(a) Definition (2 marks).} A fossil is any preserved remain, impression or trace of activity of an organism that lived in the geological past, preserved in the rocks of the Earth's crust (usually sedimentary rocks). \emph{[1 mark for ``remains or traces of past life''; 1 mark for ``preserved in rocks / geological past''.]}

\textbf{(b) Four conditions favouring fossilization (8 marks — 2 each: 1 for naming, 1 for explanation).}
\begin{enumerate}
\item \textbf{Possession of hard parts:} shells, bones, teeth and tests resist decay and mechanical destruction, so organisms with mineralized skeletons fossilize far more readily than soft-bodied ones.
\item \textbf{Rapid burial:} quick covering by sediment protects the remains from scavengers, weathering, oxidation and bacterial decay.
\item \textbf{Absence of oxygen (anoxic conditions):} oxygen-poor bottom waters or sediment inhibit decomposing bacteria and scavengers, allowing even soft tissues to be preserved.
\item \textbf{Quiet (low-energy) depositional environment:} in calm waters (lagoons, deep marine basins, lake bottoms) remains are not broken or dispersed by waves and currents.
\end{enumerate}
\emph{[Also acceptable: fine-grained sediment, early mineralization/cementation, absence of later metamorphism.]}

\textbf{(c) Fossils and crystalline rocks (6 marks).}
\begin{itemize}
\item \emph{Igneous rocks (3 marks):} they form by crystallization of magma or lava at temperatures of several hundred to over $1000\ ^\circ\mathrm{C}$; any organism caught in lava is burned and destroyed, and no organism can live within a magma, so fossils are essentially absent (rare exceptions: moulds of tree trunks at the cooled surface of lava flows, or ash beds burying remains).
\item \emph{Metamorphic rocks (3 marks):} fossils may originally exist in the sedimentary parent rock, but the heat and pressure of metamorphism recrystallize minerals and deform the rock, distorting or obliterating fossil structures; only under very weak (low-grade) metamorphism can deformed fossils exceptionally survive, e.g. stretched brachiopods in slate.
\end{itemize}

\textbf{(d) Two high-potential environments (4 marks — 2 each).}
\begin{enumerate}
\item \textbf{Deep or quiet marine basins / lagoons:} continuous, rapid deposition of fine mud in calm, often oxygen-poor water buries shells and carcasses quickly and intact (e.g. the Cretaceous shales and limestones of the Douala basin rich in foraminifera and ammonites).
\item \textbf{Lake bottoms and deltas (e.g. swampy delta plains):} rapid input of fine sediment and frequent anoxic bottom conditions preserve fishes, insects and plant remains; volcanic ash falls can also bury entire communities instantly.
\end{enumerate}
\end{corrige}

\begin{corrige}[Exercise 10 — Specimen identification (20 marks)]
\emph{Indicative mark scheme: 1 mark per correct phylum/class, about 1 mark per valid morphological feature (max.\ 3 per specimen), 1 mark per mode of life/habitat, 1 mark per range/index assessment.}

\textbf{1. Ammonite.}
\begin{itemize}
\item \emph{Phylum} Mollusca, \emph{class} Cephalopoda.
\item \emph{Morphology:} planispiral coiled shell in one plane (bilateral symmetry); shell divided into chambers separated by septa with complex (folded) suture lines; external ornament of ribs, sometimes tubercles; body chamber occupied the last whorl.
\item \emph{Mode of life:} free-swimming (nektonic) marine carnivore, moving by jet propulsion.
\item \emph{Range:} Devonian to end-Cretaceous; excellent index fossil because of rapid evolution (short ranges), wide distribution and abundance.
\end{itemize}

\textbf{2. Trilobite.}
\begin{itemize}
\item \emph{Phylum} Arthropoda, \emph{class} Trilobita.
\item \emph{Morphology:} dorsal exoskeleton of calcite/chitin divided longitudinally into three lobes (two pleural lobes and an axial lobe) and transversely into head (cephalon with eyes), segmented thorax and tail (pygidium); bilateral symmetry; jointed appendages beneath.
\item \emph{Mode of life:} benthic marine crawler, some burrowers or swimmers; shallow seas.
\item \emph{Range:} exclusively Palaeozoic (Cambrian--Permian); outstanding index fossil, especially for the Cambrian, owing to rapid evolution and worldwide distribution.
\end{itemize}

\textbf{3. Brachiopod.}
\begin{itemize}
\item \emph{Phylum} Brachiopoda (class Rhynchonellata for hinged forms).
\item \emph{Morphology:} bivalved shell with two \emph{unequal} valves (pedicle and brachial valve) but each valve bilaterally symmetrical about its own midline; pedicle opening (foramen) for attachment stalk; ornament of radial ribs or concentric growth lines.
\item \emph{Mode of life:} benthic marine, fixed to the sea floor by a fleshy pedicle; filter feeder.
\item \emph{Range:} Cambrian to Present, dominant in the Palaeozoic; many genera are good Palaeozoic index fossils.
\end{itemize}

\textbf{4. Bivalve.}
\begin{itemize}
\item \emph{Phylum} Mollusca, \emph{class} Bivalvia (Pelecypoda/Lamellibranchia).
\item \emph{Morphology:} two \emph{equal} (equivalve) valves, left and right, joined by a hinge with teeth and sockets; plane of symmetry between the valves; concentric growth lines; no head, no segmentation.
\item \emph{Mode of life:} benthic marine or freshwater; burrowers, attached forms or free livers; filter feeders.
\item \emph{Range:} Cambrian to Present; generally long-ranged, so most are poor index fossils, though some Mesozoic--Cenozoic genera are locally useful.
\end{itemize}
\emph{Examiner's expectations:} the candidate must clearly contrast the symmetry of brachiopod (symmetry through each valve) and bivalve (symmetry between the valves) — the classic GCE trap — and justify index-fossil status using range, distribution and abundance.
\end{corrige}

\begin{corrige}[Exercise 11 — Essay: the fossil record (25 marks)]
\emph{Indicative plan and mark scheme.}

\textbf{Introduction (within conclusion marks / general impression):} define the fossil record as the totality of fossils preserved in rocks; state the thesis — it is fragmentary, yet irreplaceable.

\textbf{(a) Causes of gaps (9 marks — about 2 marks per developed cause).}
\begin{itemize}
\item \emph{Preservation potential:} organisms with hard parts (shells, bones) fossilize easily; soft-bodied organisms (jellyfish, worms) almost never do, so most of past biodiversity is invisible.
\item \emph{Conditions of fossilization:} fossilization requires rapid burial, anoxia and quiet water; organisms dying on land, in oxygenated or high-energy settings are usually destroyed before burial.
\item \emph{Non-deposition and erosion:} sediments, and the fossils they contain, accumulate only in basins; uplifted areas record nothing, and previously deposited fossil-bearing rocks are eroded away.
\item \emph{Metamorphism and tectonic destruction:} burial metamorphism, deformation and subduction recrystallize or destroy fossils and their host rocks.
\item \emph{Collection bias:} many fossiliferous rocks lie at depth or in inaccessible regions and have never been sampled.
\end{itemize}

\textbf{(b) Modes of fossilization and what is preserved (6 marks).}
\begin{itemize}
\item Unaltered preservation (amber, freezing, desiccation) is exceptional and restricted to special settings and young ages.
\item Permineralization and replacement preserve hard parts, sometimes with microscopic detail, but require porous skeletons and suitable fluids.
\item Carbonization can record soft tissues, but only as flattened carbon films in fine, anoxic sediments.
\item Moulds and casts preserve only the \emph{shape} of shells whose original material has dissolved; trace fossils record \emph{behaviour} but never the animal itself. Hence each mode selects a different, partial fraction of the original biosphere.
\end{itemize}

\textbf{(c) Uses of fossils — at least four developed examples (8 marks — 2 each).}
\begin{enumerate}
\item \emph{Relative dating and correlation:} index fossils (ammonites, graptolites, foraminifera) define biozones and allow beds of the same age to be matched between distant sections, e.g. correlating Cretaceous strata across Cameroon using ammonite and foraminiferal assemblages.
\item \emph{Palaeoenvironmental reconstruction:} reef corals indicate warm, shallow, clear marine water; coal-forming plant fossils indicate swampy tropical lowlands; benthic versus planktonic foraminifera indicate water depth.
\item \emph{Evidence for evolution:} successive fossil assemblages document the appearance, diversification and extinction of groups (e.g. the evolution of ammonite sutures, the horse lineage), providing direct evidence of biological change through time.
\item \emph{Economic applications:} micropalaeontology (foraminifera, ostracods, palynomorphs) is routinely used to date and correlate subsurface strata during petroleum exploration in the Douala and Rio del Rey (Niger delta margin) basins, guiding the location of reservoir and source rocks.
\end{enumerate}
\emph{[Also creditable: palaeogeography and plate tectonics — identical fossil faunas on now-separated continents.]}

\textbf{(d) Balanced conclusion (2 marks).} The fossil record is undeniably biased towards hard-bodied, marine, rapidly buried organisms, and vast chapters of life's history are lost; nevertheless, because fossils alone give direct, datable evidence of past life, environments and climates, and because they have proven practical value in stratigraphy and resource exploration, they remain an indispensable tool of the geologist — provided their limitations are always kept in mind.
\end{corrige}

\newpage

\coursec{Summary sheet}

\begin{popfigure}
\begin{tikzpicture}[mindmap, grow cyclic,
every node/.style={concept, text=white, align=center},
root concept/.append style={concept color=popdark, font=\bfseries, minimum size=3.2cm},
level 1 concept/.append style={level distance=4.4cm, sibling angle=60, font=\small, minimum size=2.6cm},
level 2 concept/.append style={level distance=2.9cm, font=\scriptsize, minimum size=1.9cm}]
\node[root concept]{PALAEONTOLOGY}
child[concept color=popblue]{ node{Scope \& Branches}
  child[concept color=popblue!70]{ node{Palaeozoology} }
  child[concept color=popblue!70]{ node{Palaeobotany} }
  child[concept color=popblue!70]{ node{Micropalaeontology} }
}
child[concept color=popgreen]{ node{Types of Fossils}
  child[concept color=popgreen!70]{ node{Body fossils} }
  child[concept color=popgreen!70]{ node{Trace fossils} }
  child[concept color=popgreen!70]{ node{Chemical fossils} }
}
child[concept color=poporange]{ node{Fossilization}
  child[concept color=poporange!70]{ node{Rapid burial} }
  child[concept color=poporange!70]{ node{Hard parts} }
  child[concept color=poporange!70]{ node{No decay} }
}
child[concept color=popgold]{ node{Modes of Preservation}
  child[concept color=popgold!70]{ node{Unaltered} }
  child[concept color=popgold!70]{ node{Replacement} }
  child[concept color=popgold!70]{ node{Mould \& cast} }
  child[concept color=popgold!70]{ node{Carbonization} }
}
child[concept color=poppurple]{ node{Record \& Uses}
  child[concept color=poppurple!70]{ node{Index fossils} }
  child[concept color=poppurple!70]{ node{Palaeoenvironment} }
  child[concept color=poppurple!70]{ node{Gaps \& extinction} }
}
child[concept color=poppink]{ node{Classification \& Groups}
  child[concept color=poppink!70]{ node{Binomial name} }
  child[concept color=poppink!70]{ node{Corals, trilobites} }
  child[concept color=poppink!70]{ node{Ammonites, plants} }
};
\end{tikzpicture}
\legende{Mind map of the chapter ``Palaeontology'': scope, types of fossils, conditions and modes of fossilization, the fossil record and its uses, classification and the main fossil groups.}
\end{popfigure}

\begin{aretenir}[Palaeontology -- Summary Sheet]
\textbf{1. Key definitions}
\begin{itemize}
\item \cle{Palaeontology}: the scientific study of ancient life through fossils preserved in sedimentary rocks.
\item \cle{Fossil}: any preserved remain, impression or trace of an organism that lived in the geological past (conventionally older than about 10\,000 years, i.e.\ pre-Holocene).
\item \cle{Body fossil}: actual preserved part of an organism (bone, shell, tooth, leaf).
\item \cle{Trace fossil} (ichnofossil): evidence of the activity of an organism -- footprints, tracks, burrows, borings, coprolites (fossilized faeces).
\item \cle{Chemical fossil} (biomarker): organic molecule preserved in a rock that proves the former presence of life.
\item \cle{Taphonomy}: the study of everything that happens to an organism between its death and its discovery as a fossil (decay, burial, diagenesis).
\item \cle{Index fossil}: a fossil used to date and correlate rocks -- it must have a wide geographic range, a short stratigraphic range, be abundant, and be easily identified.
\item \cle{Extinction}: the permanent disappearance of a species from the whole Earth.
\end{itemize}

\textbf{2. Conditions for fossilization (to know by heart)}
\begin{itemize}
\item Possession of \cle{hard parts} (shell, bone, teeth, wood).
\item \cle{Rapid burial} by sediment, protecting the remains from scavengers, physical destruction and decay.
\item \cle{Anoxic (oxygen-poor) conditions}, which slow down bacterial decomposition.
\item A low-energy depositional environment (lagoon, deep sea floor, lake bottom, swamp).
\item Absence of later destructive processes (metamorphism, intense weathering, erosion, dissolution by ground-water).
\end{itemize}

\textbf{3. Modes of fossilization}
\begin{itemize}
\item \cle{Unaltered remains}: frozen mammoths in permafrost, insects trapped in amber, desiccated (mummified) carcasses, shells preserved with their original mineralogy.
\item \cle{Permineralization / petrification}: mineral matter (silica, calcite, iron compounds) deposited by ground-water fills the pores of bone or wood while the original structure is retained.
\item \cle{Replacement}: the original material is dissolved and replaced molecule by molecule by another mineral (silicification, pyritization).
\item \cle{Recrystallization}: an unstable mineral (e.g.\ aragonite of shells) changes into a stable polymorph (calcite) with loss of fine detail.
\item \cle{Carbonization}: volatile elements (H, O, N) are driven off, leaving a thin carbon film that preserves the outline (leaves, fish, graptolites).
\item \cle{Mould and cast}: the external or internal mould is the impression (negative) left in the sediment; the cast is the sediment or mineral fill that reproduces the original shape (positive).
\item \cle{Imprints and traces}: footprints, trails, burrows and coprolites preserved in sediment.
\end{itemize}

\textbf{4. Preservation potential and occurrence}
\begin{itemize}
\item High preservation potential: marine organisms with hard calcareous shells or skeletons (molluscs, brachiopods, echinoderms, corals).
\item Low preservation potential: soft-bodied organisms, terrestrial animals, and organisms of tropical forests where decay is rapid.
\item Fossils occur almost exclusively in \cle{sedimentary rocks} (limestone, shale, sandstone); igneous crystallization and metamorphic recrystallization destroy them.
\end{itemize}

\textbf{5. Uses of fossils}
\begin{itemize}
\item \cle{Stratigraphic correlation} and relative dating of rock strata (index fossils, zones).
\item Reconstruction of \cle{palaeoenvironments} and palaeoclimates (e.g.\ corals indicate warm, shallow, clear marine water).
\item Evidence for \cle{evolution} (transitional forms, succession of faunas) and for plate tectonics (identical fossil assemblages on continents now separated by oceans).
\item Economic geology: microfossils (foraminifera, ostracods, pollen and spores) guide \cle{petroleum} exploration, as in the Rio del Rey and Douala basins of Cameroon; fossiliferous limestones are raw material for cement.
\end{itemize}

\textbf{6. Gaps in the record and extinction}
\begin{itemize}
\item Gaps in the fossil record arise from non-deposition, erosion of strata, metamorphism, the non-fossilization of soft-bodied organisms, and fossil localities not yet discovered.
\item Causes of extinction: climate change, sea-level change, volcanic activity, meteorite impact, competition, predation and disease; mass extinctions include the end-Permian and end-Cretaceous events.
\end{itemize}

\textbf{7. Classification and main groups}
\begin{itemize}
\item Linnaean hierarchy: Kingdom, Phylum, Class, Order, Family, Genus, Species; \cle{binomial nomenclature} gives every species a two-part italicized name (\emph{Genus species}).
\item Know the morphology, mode of life and stratigraphic range of: foraminifera, sponges, corals, graptolites, bivalves, gastropods, ammonites, brachiopods, trilobites, crinoids and echinoids, vertebrates (fish, reptiles, mammals) and fossil plants.
\end{itemize}

\textbf{8. Common mistakes (examiner's warnings)}
\begin{itemize}
\item Confusing the \emph{mould} (impression, negative) with the \emph{cast} (fill, positive).
\item Saying that fossils occur in ``all rocks'' -- they are essentially restricted to sedimentary rocks.
\item Confusing trace fossils with body fossils; a coprolite is a trace fossil, not a body fossil.
\item Giving the index-fossil criteria incompletely: all four criteria (wide geographic range, short stratigraphic range, abundance, easy identification) are required.
\item Forgetting that extinction is permanent and global; do not confuse it with the local disappearance of a species.
\end{itemize}

\textbf{9. GCE tips}
\begin{itemize}
\item In any fossil description always give: classification, morphology (with a labelled sketch), composition of the hard parts, mode of life, geological range, and use.
\item Draw neat, labelled diagrams (ammonite sutures, the three longitudinal lobes of a trilobite, coral septa) -- marks are awarded for accurate labels.
\item Use Cameroonian examples where possible: Cretaceous and Tertiary fossils of the Douala and Rio del Rey sedimentary basins, and the limestones associated with the Benue Trough.
\end{itemize}
\end{aretenir}

\findecours

\end{document}
